\documentclass[a4paper,12pt]{article}
\usepackage{fullpage}
\usepackage{newtxtext} % Times-like text font
\usepackage{newtxmath} % Times-like math font
\usepackage[utf8]{inputenc}
\usepackage{setspace}
\usepackage[english]{babel}
\usepackage{amsmath}
\usepackage{pgfplots}
\usepackage{comment}
\usepackage{amsfonts}
\usepackage{amssymb}
\usepackage{caption}
\usepackage{pdflscape}
\usepackage{algorithm}
\usepackage{algpseudocode}
\usepackage{afterpage}
\PassOptionsToPackage{hyphens}{url}\usepackage{hyperref}
\usepackage{xurl}
\usepackage{natbib}
\usepackage{chngcntr}
\usepackage[toc,page,header]{appendix}
\usepackage{minitoc}
\usepackage{csquotes}
\usepackage[noabbrev,capitalize]{cleveref}
\setcitestyle{citesep={;}}
\usepackage{longtable}
\usepackage{array, makecell}
\usepackage{appendix}
\usepackage{amsmath,siunitx,booktabs}
\usepackage{booktabs}
\usepackage{threeparttable}
\usepackage{ltablex}
\usepackage{array, makecell}
\usepackage{etoc}
\usepackage{import}
\usepackage{xcolor}
\definecolor{darkblue}{rgb}{0.0, 0.0, 0.55}
\definecolor{darkred}{rgb}{0.55, 0.0, 0.0}
\usepackage{hyperref}
\hypersetup{
colorlinks=true, % false: boxed links; true: colored links
linkcolor=darkred, % color of internal links
citecolor=darkblue, % color of links to bibliography
filecolor=magenta, % color of file links
urlcolor=darkblue, % color of external links
anchorcolor=blue
}

\renewcommand{\algorithmicrequire}{\textbf{Input:}}
\renewcommand{\algorithmicensure}{\textbf{Output:}}

\usepackage{subcaption}
\usepackage{mathtools}
\usepackage{tablefootnote}
\usepackage{setspace}
\usepackage{graphicx}
\usepackage{bm}
\usepackage{multirow}
\usepackage{lscape}
%\usepackage{caption, floatrow}
% \captionsetup{labelfont={sc, small}}
%\DeclareFloatVCode{myrowsep}{\vskip 4ex}
\onehalfspacing
\newcommand\independent{\protect\mathpalette{\protect\independenT}{\perp}}
\def\independenT#1#2{\mathrel{\rlap{$#1#2$}\mkern2mu{#1#2}}}
\let\oldemptyset\emptyset
\let\emptyset\varnothing
\usepackage{xcolor}
\newtheorem{theorem}{Theorem}
\newtheorem{lemma}{Lemma}
\newtheorem{condition}{Condition}
\newtheorem{proposition}{Proposition}
\newtheorem{corollary}{Corollary}
\usepackage{natbib}

% Make the "Part I" text invisible
\renewcommand \thepart{}
\renewcommand \partname{}


\makeatletter
\@addtoreset{theorem}{section}
\makeatother

\usepackage{apptools}
\AtAppendix{\counterwithin{lemma}{section}}
\AtAppendix{\counterwithin{remark}{section}}

\newtheorem{assumption}{Assumption}
\newtheorem{remark}{Remark}

\DeclareMathOperator*{\argmax}{arg\,max}
\DeclareMathOperator*{\argmin}{arg\,min}


\usepackage[most]{tcolorbox}
\usepackage{enumitem}
\usepackage{listings}

\usepackage{ltablex}  % longtable + tabularx
\keepXColumns         % keep X-column widths consistent across pages
\usepackage{makecell} % for line breaks inside cells

\newtcolorbox{promptblock}[1][]{
  enhanced, breakable,
  colback=darkred!4!white,      % subtle red-tinted background
  colframe=darkred!80!black,    % dark-red border
  colbacktitle=darkred!90!black,% title bar background
  coltitle=white,               % title text color
  boxrule=0.5pt,
  arc=1.5mm,
  left=6pt, right=6pt, top=4pt, bottom=4pt,
  boxsep=2pt,
  fonttitle=\bfseries,
  fontupper=\small,
  title=#1
}
% shrink JSON listing text and vertical skips
\lstset{
  basicstyle=\ttfamily\footnotesize,
  columns=fullflexible, breaklines=true, showstringspaces=false,
  frame=single, aboveskip=4pt, belowskip=4pt, framexleftmargin=2pt, framexrightmargin=2pt
}

\newcommand{\redblock}[1]{%
  \begingroup\setlength{\fboxsep}{2pt}%
  \colorbox{red!20}{\textbf{\textcolor{red!80!black}{#1}}}%
  \endgroup
}

\title{The Impact of Banning Online Gambling Livestreams: \\Evidence from Twitch.tv\footnote{We thank Pinar Yildirim, Artem Timoshenko, Andreas Lanz, Tong Guo, Carl Mela, Oded Netzer, Hiroaki Kaido, Marc Rysman, Runhong Ma, Yuxiao Hu, and seminar participants at Marketing Science, New Data for Consumer Insights Conference at Chicago Booth, 2024 Conference on Artificial Intelligence, Machine Learning, and Business Analytics, Boston University, Columbia University, Chinese University of Hong Kong, Shenzhen, and Shanghai Jiao Tong University for helpful comments and suggestions.}}
% \title{The Impact of Banning Online Gambling Livestreams: \\Evidence from Twitch.tv}
\author{Qifan Han\thanks{Chinese University of Hong Kong, Shenzhen. Email: \href{mailto:hanqifan@cuhk.edu.cn}{hanqifan@cuhk.edu.cn}} \and Jasmine Yang\thanks{Chinese University of Hong Kong and Meta. Email: \href{mailto:yyang24@gsb.columbia.edu}{yyang24@gsb.columbia.edu}} \and Andrey Simonov\thanks{Columbia University and CEPR. Email: \href{mailto:as5443@gsb.columbia.edu}{as5443@gsb.columbia.edu}}}
% \author{}
\date{\today}
%\date{October 15, 2025}

\begin{document}
\doparttoc % Tell to minitoc to generate a toc for the parts
\faketableofcontents % Run a fake tableofcontents command for the partocs

\maketitle
\begin{spacing}{1.5}

\begin{abstract}
\setstretch{1}

\noindent How effective is platform self-regulation at eliminating harmful content? We examine Twitch’s ban on unlicensed online gambling livestreams implemented in October 2022. Using a novel panel dataset covering the top 6,000 Twitch streamers, we identify banned content and affected streamers by leveraging video analysis of historical clips, high-frequency stream titles, and in-stream chats. On the supply side, the policy led to a 63.2\% reduction in weekly gambling streams and a 44.3\% decrease in overall streams among streamers whose content was banned. Streamers whose gambling content was not banned also reduced their gambling and overall content, indicating spillover effects from the policy. This reaction was more pronounced among popular streamers and those with greater reputation concerns. On the demand side, while the policy reduced total viewership, low-tier subscriptions, and engagement, revenue from high-tier subscriptions (reflecting more loyal viewers) remained unaffected. The policy also reduced the gambling viewership for viewers classified as minors, although the magnitude was slightly smaller than for other viewers. Our findings highlight both the potential and limitations of platform self-regulation: although targeted bans can substantially reduce targeted content and its engagement, complementary actions may be necessary to mitigate unintended side effects and ensure effective protection for specific user groups. \\

\noindent \textbf{Keywords}: content regulation, platform self-regulation, online gambling, live streaming, minor protection, unstructured data
\end{abstract}

\thispagestyle{empty}

\newpage 
\setcounter{page}{1}
\pagenumbering{arabic}

\section{Introduction}\label{section-introduction}
The increasing prevalence of misinformation \citep[e.g.][]{beck2023guardians,gandhi2023misinformation,ananthakrishnan2020tangled} and harmful content \citep[e.g.][]{beknazar2022toxic,aridor2024economics} on digital platforms has sparked a significant debate over the necessity of content regulation. As technological advancements expand the range of regulatory challenges for governments, content regulation has increasingly shifted to platforms adopting self-regulation measures. For example, Facebook exploits automatic removal of comments that are classified as toxic content.\footnote{Source: \url{https://www.nytimes.com/2021/08/31/technology/facebook-accenture-content-moderation.html}.} However, the optimal form of self-regulation often depends on how the platform profits \citep{liu2022implications}, potentially limiting its success due to the misalignment between the platform's economic interests and regulation target \citep{cusumano2021can}. Furthermore, the success of platform self-regulation hinges on content producers’ compliance with the policy and their avoidance of circumventing it by altering their content production \citep{aridor2024economics}. %  While platforms may strive to balance between these outcomes, it remains unclear whether a balance-oriented policy can achieve regulatory success without causing unintended consequences.

In this paper, we investigate the regulatory effects and economic consequences of content self-regulation on Twitch, a major game streaming platform. The content of interest are streams of online gambling (e.g., slots, virtual casino) that have recently surged on Twitch, experiencing a 166\% increase in viewership between the first quarters of 2020 and 2022 \citep{Streamscheme,Streamhatchetb}. This increase in online gambling popularity has been accompanied by several high-profile scandals, with popular streamers borrowing large sums of money and defrauding followers to support gambling pursuits, along with instances of viewers falling victim to scam roulette games.\footnote{Source: \url{https://www.wired.com/story/twitch-streamers-crypto-gambling-boom/}.} The situation is further exacerbated when viewers are minors, who lack maturity to assess risks and are susceptible to developing addictive habits when streamers leverage their influence to advertise these services through streaming.

In response to the concerns, Twitch announced a policy on September 20, 2022 that banned streamers from broadcasting online gambling websites that include slots, roulette or dice games and lack official licenses in the U.S. or other jurisdictions with sufficient consumer protection.\footnote{Source: \url{https://x.com/Twitch/status/1572347129192132611?lang=en}.} Starting from October 18, 2022, four major websites -- Stake, Rollbit, Duelbits, and Roobet -- were banned by this policy, whereas streaming of gambling on other websites was not restricted.\footnote{The limited scope of this policy partially aligns with the findings in \citet{liu2022implications}, suggesting that platforms relying on advertising are less aggressive in content moderation to maintain their user base.} Interestingly, the four targeted websites did not differ significantly in their popularity on Twitch or in their website designs aimed at encouraging online gambling. As a result, the limited scope of the policy potentially allowed streamers to substitute to gambling content from unrestricted websites, leading to uncertainty of the policy's impact on regulating gambling content. 

%However, Twitch presents a unique case as it generates revenue from both advertising and subscriptions, while its content providers (streamers) mainly profit from subscription income. 
%Therefore, Twitch's banning policy provides a lens to understand the effect of content self-regulation in reducing the supply of regulated content, highlighting the challenges posed by misaligned incentives among platforms, streamers, and policymakers in balancing economic interests, user engagement, and consumer protection.

To unveil this uncertainty, we examine how content producers and consumers react to Twitch's ban on livestreams of gambling content, and whether the policy ban had any unintended impact on the platform. On the one hand, content producers might substitute their banned livestreams with other unbanned gambling or gambling-like content (e.g., gambling-like elements in video games) for their viewers. In this case, the ban incurs minimal costs to the platform but risks not being effective in reducing harmful effects of gambling, merely shifting an exposure of viewers to gambling to another form. On the other hand, content producers might significantly decrease their production or switch to other streaming services with more lenient gambling rules. Consumers may subsequently leave the platform, either following their favored streamers or due to unmet content needs. If this occurs, the misalignment between the policy's target and its actual economic consequences makes the self-regulation costly to the platform. Understanding the policy impacts and broader implications helps platforms and policymakers evaluate the impact and costs of such regulatory measures. 

We assemble a novel high-frequency streaming dataset covering the top 6,000 livestreamers from August 2022 to December 2022 to study the impact of Twitch's policy ban. Alongside stream activity data, we also collect extensive game-level data, including the presence of gambling-like features in over 5,000 video games streamed during our study period. To tackle the key challenges of detecting banned versus unbanned content within online gambling livestreams and identifying streamers who have streamed banned content, we leverage video analysis on historical video clips, as well as text analysis on stream titles and in-stream chat logs. We model the policy impact using a two-way fixed effects difference-in-differences approach, and show the robustness of our results by proposing a novel procedure that incorporates network analysis in our estimation to internalize potential interference of policy at the treatment group level or the streamer community level. This novel method helps us mitigate concerns over potential violations of the Stable Unit Treatment Values Assumption (SUTVA).

All our empirical approaches yield qualitatively similar results and insights. First, we find that the banning policy led to a significant decrease of approximately 63.2\% in weekly streams of gambling content among streamers who have streamed banned content (``banned streamers" hereafter), as well as a 12\% reduction in weekly gambling streams among streamers who have only streamed unbanned gambling content (``unbanned streamers" hereafter). These estimates show that the policy successfully reduced the overall supply of gambling content on Twitch by completely removing banned gambling content and decreasing unbanned gambling content. Through a heterogeneous treatment effects analysis, we find that only streamers who relied on gambling content the most increased their livestreams of video games with gambling-like content, suggesting a mild substitutablity between real gambling and those video games. More importantly as a potentially unexpected consequence of the policy, streamers of both banned and unbanned content produce less content overall, with a reduction of weekly streams of 44.3\% and 17.6\%, respectively. Finally, the policy had a more significant impact on streamers with higher popularity, which can be suggestively explained by the concerns from these streamers for their reputation.

On the demand side, we find that both banned and unbanned streamers experienced reductions in total hours watched by viewers, whereas the magnitude of reduction in content consumption was even higher compared to the reduction in content creation. Our analysis of subscription levels further reveal that only the lowest tier (cheapest) subscriptions decreased after the policy announcement for treated streamers. This finding indicates that while treated streamers saw a drop in revenue from casual viewers, they did not experience significant losses among their loyal viewers or in engagement from their core communities. The banning policy also affected both engagement volume and engagement quality from viewers of the treated streamers. These streamers had fewer chatters and chat messages, and each message contained fewer emotes, characters, and tokens.

We are particularly interested in the effect of banning gambling content on Twitch viewers who are minors (under 18 years old), since minors might be more susceptible to risks coming from gambling. While Twitch does not disclose the share of viewers who are minors, they are likely prominent on the platform due to its gaming focus; e.g. 17\% of teens say they use Twitch \citep{PewResearch2023}. To understand the effect of the policy on viewers who are likely to be minors, we classify viewers in our sample into ``likely minors'' or ``likely adults'' by checking whether they watch predominantly children-focused content on the platform. We find that viewers classified as minors were still exposed to gambling content in the pre-policy period, and that the policy was effective in reducing these viewers exposure to gambling content on Twitch, although the effect size was slightly smaller compare to non-minor viewers. 

% Finally, we also examine whether the policy has an effect on traffic to banned gambling websites, whose services may be promoted in streams by streamers or through advertisements on Twitch. We compiled a list of online gambling websites and checked the effect of the policy on the traffic of banned and non-banned websites. Our results suggest that there is no significant or economically meaningful decline across different traffic sources of banned gambling websites, suggesting that the policy had little effect on the gambling websites themselves.

To the best of our knowledge, our research is the first to quantify the causal effects of a banning policy targeting online gambling on a livestreaming platform and to explore the decision-making mechanisms of streamers, viewer demands for a variety of content, as well as its impact on minor viewers. While previous research has shown that gambling-like content in video games provides similar psychological experiences to real gambling \citep{drummond2018video,larche2021rare,amano2024makes}, and that consumption of this content is correlated with future gambling behavior \citep{zendle2019adolescents,kristiansen2020loot,close2021secondary}, there remains a gap in understanding whether content producers and consumers differentiate this content from genuine online gambling. Our findings provide at least threefold insights for content providers, livestreaming platforms and regulators. First, banning policies with a limited scope can generate spillover effects and reduce non-targeted content production, highlighting the economic costs of such policies for platforms and a need for more careful policy design. Second, we find no evidence of substitution from watching gambling to watching video games with gambling-like features (e.g. loot boxes). Third, we show that viewers who are more likely to be minors are exposed to gambling content on Twitch, and that banning polices also limit the exposure of these viewers to gambling content. Our findings indicate that platforms and policymakers can use data on both engagement volume and intensity to evaluate regulatory impacts and refine policies for more effective protection.

\section{Related Literature}\label{section-literature}

This paper contributes to several strands of existing literature. First, we contribute to the literature on platform self-regulation and content moderation in digital environments, which typically focuses on misinformation \citep{beck2023guardians, gandhi2023misinformation, ananthakrishnan2020tangled} and harmful content \citep{howard2019free,beknazar2022toxic,beknazar2024model}. Particularly, \citet{liu2022implications} explores the relationship between a platform's monetization channel and its optimal strategy for content moderation. They find that platforms under advertising behave less aggressively than those under subscriptions, because the former relies on a larger user base, whereas the latter focuses more on consumers' willingness to pay. Our context aligns with this result as Twitch mainly profits from advertising and adopted a less aggressive banning policy. We contribute to this literature in two additional ways: first, most existing studies analyze centralized moderation where platforms directly remove or downrank content, and users decide whether to stay or leave. In contrast, Twitch relies on individual content providers to re-optimize their own content decisions and on consumers to shift their viewing choices. Second, we provide empirical evidence for the impact of content moderation on both advertising\footnote{We cannot directly collect data on Twitch advertisements. However, since Twitch primarily uses pre-roll and mid-roll ads, its advertisement income is significantly tied to total hours watched by viewers, which we discuss in the demand-side analysis.} and subscription income as Twitch makes profits through both channels, extending the theoretical insights to a setting with a more complex and realistic revenue structure. Many theoretical works have also discussed the unintended consequences of self-regulation policies on digital platforms, including regulations on ad transparency \citep{wu2022regulating}, search self-preferencing \citep{zou2024self}, and product safety certification \citep{iyer2018voluntary}. We provide new evidence to this literature by showing that self-content regulation on Twitch led to an unintended reduction in the production of unregulated content, highlighting a potential misalignment between regulation policy targets and the platform's economic interests.

While our primary focus is on livestreams of online gambling through virtual casino games, viewers of these livestreams may also develop addictions to sports betting, which is also allowed to be streamed on Twitch. Therefore, our paper connects to the growing literature that examines how the legalization of sports gambling affects household financial health \citep{baker2024gambling,hollenbeck2024financial}, tax revenue, and irresponsible gambling behavior \citep{taylor2024online}. Our settings provides causal evidence of real-time consumer engagement, which is complementary to existing work on household responses. Our findings may also inform potential regulations on sports gambling livestreams to mitigate their negative financial and social impacts.

Second, we add to the literature of digital platform policies and consumer protection \citep{blaseg2020consumer,hunold2020rankings}, especially targeting minors and young adults \citep{braghieri2022social}. Unlike studies that focus on regulations targeting child-directed content alone \citep{johnson2024coppacalypse,cecere2025privacy}, we study a regulation that applied to all viewers but discuss heterogeneous effects across age groups. We find that the banning policy was less effective in reducing minors' gambling content consumption and engagement intensity, calling for more cautious consumer protection strategies or complementary actions to better safeguard minors.

Third, our paper contributes to the growing literature on content creation, consumption, and user engagement on livestreaming platforms \citep{zhao2023direct, qian2025express, huang2025promotional}. 
\cite{hu2017audiences} shows that both streamer and audience group identification are positively associated with viewers' continuous watching intention. \cite{lin2021happiness} finds that the happiness of streamers intensifies various viewer engagement, including chatting and tipping. We extend this literature by studying how heterogeneous viewer groups, such as loyal versus peripheral viewers, and minors versus non-minors, differed in their responses to exogenous regulatory shocks. We show substantial heterogeneity in both demand and engagement metrics, emphasizing the importance of understanding different motivations and behavioral patterns among viewers on livestreaming platforms. Our paper also connects to the broader literature on platform economies of social media platforms \citep{allcott2020welfare,aridor2024economics}. Particularly, we provide new empirical evidence on how policy intervention affect consumer engagement intensity and their monetization incentives. 

The remainder of the paper is structured as follows. Section \ref{section-empirical-setting} introduces the research context and describes our data sources. Section \ref{section-treatment-detection} describes the procedure of detecting banned content and streamers. Section \ref{section-empirical-strategy} discusses the main identification challenges and our empirical strategies. Section \ref{section-results} presents our main results. Section \ref{section-minor} discusses heterogeneous policy impact on minor versus non-minor viewers. Section \ref{section-conclusion} concludes.


\section{Empirical Context and Data}\label{section-empirical-setting}
\subsection{Gambling Streams on Twitch}\label{subsection-platform-overview}


Twitch is a live video platform where \textit{streamers} broadcast real-time content to \textit{viewers}. Streams are categorized by game or activity (e.g., ``Slots,'' ``Just Chatting,'' ``Fortnite''). Each stream has a real-time \textit{chat} window; participation requires a registered account, providing a proxy for engaged viewership.

Viewers can \textit{subscribe} to streamers at three price tiers: Tier~1 (\$4.99/month), Tier~2 (\$9.99/month), and Tier~3 (\$24.99/month). Streamers receive approximately 50\% of subscription revenue.

In a typical banned gambling stream, a streamer plays slots or roulette on websites such as Stake.com while commenting on results and encouraging viewers to sign up via affiliate links. The streamer's screen displays the gambling website interface, which our video classification pipeline detects (Section~\ref{section-treatment-detection}). Unbanned gambling streams look visually similar but feature branding from licensed platforms such as PokerStars.


\subsection{Data}\label{subsection-data}
We combine several novel data sources to study the effect of the banning policy and to identify banned gambling content and streamers. 

\noindent \textbf{Twitch streaming data.} We constructed a novel high-frequency dataset covering the top 6,000 Twitch streamers over five months (August 1--December 31, 2022) using two sources. First, we obtained the top 6,000 streamers and their stream IDs from \textit{Sullygnome}, a third-party website that reports stream statistics from Twitch. Based on the stream IDs, for August 1--October 26, we obtained the status of the streamer (online/offline), streamed content, viewer count and the title of the streams, recorded every 10 minutes, from \textit{Twitchtracker}, a third-party website that monitors streaming activity on Twitch. For October 26 onward, we used high-frequency data on the same streamers collected by \cite{yangwiptwitch} through sending requests to the Twitch API every 15 minutes (see Appendix~\ref{data-collection} for details).\footnote{We compile the first dataset from \textit{Sullygnome} and \textit{Twitchtracker}, since the second dataset does not cover data prior to October 26, 2022. To ensure the quality of the first data source, we checked a few days of overlapping data from both sources and found them consistent, except that the first dataset occasionally misses stream title information and merges streams that live for a short duration into one stream.} Together, the combined data provides a comprehensive view of the streaming activities before and after the policy. 

% , recording streaming metrics including total streaming hours, total hours watched by viewers, weekly average viewership, and total streaming hours of each content within a week for a streamer. 

For our main analyses, we aggregate these data to weekly observations for each streamer. To identify the causal effects of the policy, we exclude streamers active only before or only after the policy announcement.\footnote{Removing these streamers means that we do not consider the policy's impact on market entry or exit in this paper. This is because most of the removed accounts are official accounts of video game developers, or of the banned gambling websites Twitch permanently deactivated after the policy. Since these accounts did not provide individual user-generated content like unofficial streamers in our sample, removing them does not compromise the paper's main focus.} We classify streamers into three groups based on pre-announcement activity: banned, unbanned and untreated. We define \textit{banned streamers} as those who had streamed content from banned websites before the policy announcement, viewing them as directly impacted by the policy. \textit{Unbanned streamers} are those who streamed online gambling, but only from unbanned gambling websites. Although not directly targeted, these streamers may alter their content given changes in the competitive environment or perceived regulatory risk. Finally, \textit{untreated streamers} are those that never streamed gambling content. These people are less likely to be affected by the policy and may behave similarly to non-gambling streamers. Although high-frequency data allows us to recover streaming activities, our data do not label which gambling streams (and streamers) feature banned versus unbanned content. We discuss how we identify the treatment status of streamers in Section \ref{section-treatment-detection}. 

\noindent \textbf{Game attributes data}. We also collected data on video game attributes (e.g., developer, genre, platform, and age rating) from \textit{IGDB}, one of the most extensive online video game databases owned by Twitch. We use these data in Section \ref{section-minor} to detect content targeting minors. 

\noindent \textbf{Games with gambling-like features.} To study the substitution patterns between online gambling and games with gambling-like mechanics (e.g., loot boxes, gacha systems), we compiled a unique, comprehensive game list from three sources given there is no readily available data. First, 
we extracted games tagged with “Loot boxes” or “Gacha system” from \textit{GiantBomb}, an American video game website and wiki with a dedicated concepts taxonomy. These concepts capture the practice of randomly obtaining in-game items with predetermined odds through in-game currencies purchasable with real money, mimicking aspects of genuine gambling. Second, we 
conducted text analysis on stream titles in our high-frequency dataset, identifying games whose title contain keywords closely connected to gambling-like activities, such as “loot box opening”, or “insane pulls”. Third, we supplemented the list with games carrying in-game purchase labels that disclose the presence of loot boxes from \textit{Pan-European Game Information (PEGI)} and \textit{The Entertainment Software Rating Board (ESRB)}, two well-known age rating organizations.\footnote{We manually cleaned the additional games using both websites' public search tools with filters following \cite{xiao2023beneath}. We note that all games with gambling-like features carry in-game purchase labels, especially older titles; only 53\% of our detected games are also covered by the third source.} See Appendix Figure \ref{lootbox_ratings} for examples. Our final list covers 464 unique video games streamed on Twitch during our studied period. 

Eventually, we classified all streamed content into four categories: online gambling games, games with gambling-like features (i.e., ``loot box games''), games without         
gambling-like features (i.e., ``other games''), and non-gaming content (e.g., ``Just Chatting''). The online gambling category covers all streams labeled under `Slots'' (the largest gambling category on Twitch), ``slots!'', ``Blackjack'', ``Casino" and ``Virtual Casino''. As the policy ban is not on all online gambling content, we detail how we distinguish banned from unbanned content within this category in Section \ref{section-treatment-detection}.

\noindent \textbf{Video data.} We collected all video clips and available Video on Demand from online gambling streams of individual streamers before the policy announcement from Twitch API. We use these video clips to identify banned content and latent treatment status of streamers in Section \ref{section-treatment-detection}.

\noindent \textbf{Chat logs and emotes data.} We collected in-stream chat logs from \textit{Streams Charts}, whose coverage aligns with our high-frequency streaming data. For each stream, we observe the text messages, sender account names and IDs, and timestamps. We use chat data for three purposes: an additional source to identify the latent treatment status of streamers (Section~\ref{treatment-detection-chats}), constructing viewer-level treated and untreated groups (Section~\ref{section-demand-results}), and identifying minor viewers (Section~\ref{section-minor}).\footnote{The Twitch API reports aggregate viewer counts but not individual viewer lists; chat participation proxies for engaged viewership.}

From chat logs, we also constructed engagement measures such as chat counts, session counts, message length, and number of emotes (including emojis). For emotes data, in particular We
constructed the full list of 199,984 emotes on Twitch via its API, covering both global emotes and channel-specific emotes earned through following, subscribing, cheer bitting, or participating in channel events (see Appendix~\ref{appendix-chat-engagement-data} for details). We use these data to study the effect of the policy ban on viewer engagement (see Section \ref{subsection-total-chat-engagement} for details). 

\noindent \textbf{Subscriptions and revenue.} We collected per-stream subscribers gained by streamer from \textit{Streams Charts}, broken down by type (new, re-subscriptions, gifted subscriptions) and tier (Tier 1, Tier 2, Tier 3).\footnote{Only streamers who enabled channel tracking are covered. Streams Charts does not count offline or auto-renewed subscriptions unless the viewer explicitly notifies in chat. While we use this subscription data for revenue discussion, we note its limitation of possibly not being fully representative, as streamers self-select into disclosing this information.} We use the data to analyze streamer revenue changes due to the policy.

\noindent \textbf{Other streamer-level data.} For each streamer, we also collected their country and city of residence and channel prohibition histories (count, dates, duration of each prohibitions) from \textit{Streams Charts}. Residence data assist in matching video clips to streams (see Appendix~\ref{appendix-matching-clips-to-streams}); prohibition histories help explain the mechanisms behind heterogeneous treatment effects.

\noindent \textbf{Website traffic data.} We collected estimated visits and engagement for websites at the domain and subdomain level from \textit{Semrush}\footnote{\url{https://www.semrush.com/kb/997-semrush-data}}, with breakdowns by channel, device, and country. We use these data to examine whether banned gambling websites experienced traffic changes after the policy (see Section~\ref{section-website-traffic}). 

\section{Detection of Banned Content and Streamers} \label{section-treatment-detection}
Our streamer-level data does not specify whether a gambling stream was about banned or unbanned websites, nor does it indicate whether a streamer has streamed banned gambling content.
We combine several novel data sources and methods (video analysis, stream title analysis and chat log classification) to identify banned gambling content and streamers, which we overview in this section and provide full procedural details in Appendix~\ref{appendix-banned-content-detection}.

\noindent \textbf{Video analysis.}
We first construct a list of streamers and define them as ``treated'' if they streamed gambling content at least once before the policy is implemented. For these streamers, we analyze whether their pre-policy gambling video clips (over 4TB) and available Video on Demand (VOD) archives using optical character recognition (OCR)\footnote{\url{https://github.com/tesseract-ocr/tesseract.}} to extract on-screen text and match keywords from banned websites (Stake, Rollbit, Duelbits, Roobet). To ensure no misclassifications, we use a restrictive criterion by classifying a video as containing banned content only when such keywords appear in \textit{all} sampled frames of a video. See Appendix~\ref{appendix-treatment-detection-clips} for details on the video analysis and detection procedure. 

\noindent \textbf{Stream title analysis.} Next, we complement our detection of banned streams and streamers using real-time stream titles captured by our high-frequency data. This detection strategy is based on a recognized Twitch convention that gambling streamers often include referral-link commands (e.g., ``!Stake'', ``!Roobet'') in stream titles, directing viewers to gambling websites in exchange for commissions.\footnote{This is common practice within the community of gambling content streams. It became even more prevalent after Twitch banned the sharing of links to gambling sites in 2021, in response to concerns from the National Council on Problem Gambling regarding the risk of under-age gambling and gambling addiction. Source: \url{https://www.ncpgambling.org/news/ncpg-responds-to-twitch-banning-gambling-sites-links/}.} We define a gambling stream as containing banned content if both (1) a banned website's referral link is observed in the stream title and (2) the game being streamed at that time is also online \textit{gambling} content. We demonstrate the validity of this approach in Appendix~\ref{appendix-treatment-detection-title} by showing that non-gambling streamers almost never include banned websites' referral links in their stream titles, and that gambling streamers do not use referral links in their non-gambling streams. 

\noindent \textbf{Chat log analysis.} While some streamers may neither make video clips accessible nor include referral links in stream titles, we complement the identification of streamers' latent treatment status with a third approach using in-stream chat logs. We first construct a ground truth sample of 534 banned and 821 unbanned gambling streams based on the first two data sources using restrictive, multiple-source criteria. For the \textit{banned} ground truth, a gambling stream must satisfy both: (1) its matched video clips contain banned website names in all sampled frames via OCR, and (2) its stream title contains referral links to banned websites. For the \textit{unbanned} ground truth, a gambling stream must satisfy all three: (1) its matched video clips contain no banned website names detected in any sampled frame via OCR, (2) the streamer has no pre-policy gambling video clips with any banned content detected by OCR, and (3) the stream title contains neither referral links nor keywords of banned websites (see Appendix~\ref{appendix-treatment-detection-chats} for justifications on the sample validity). 

Next, we use the average number of referral links per hour as a test statistic and implement a threshold-based classifier. 
This detection strategy is based on the key industry insight that gambling streamers commonly use chat bot commands to automatically respond to viewers' referral link requests, producing a measurable difference in referral link proportions between banned and unbanned streams. We conduct several validation analyses in Appendix~\ref{appendix-treatment-detection-chats} and demonstrate that a threshold-
based classification approach based on referral links per hour is adequate. We use grid search with 10-fold cross validation to identify the optimal threshold that jointly minimizes Type~I and Type~II errors.\footnote{In our context, Type~I error is defined as incorrectly classifying an unbanned stream as banned, and Type~II error is defined as incorrectly classifying a banned stream as unbanned.} The out-of-sample classification accuracy is 97.5\% on average, demonstrating strong predictive performance (see Appendix~\ref{inference-step-chats} for full performance metrics, confusion matrices and ROC curves).

\noindent \textbf{Identification of treated groups using all sources.} We combine detection results from all three approaches to construct a list of streams containing banned content. Then, we use these identified banned streams to classify treated streamers into banned and unbanned groups.  Appendix Figure~\ref{figure:detection_banned_source} illustrates the number of streamers detected by each approach alone and in combination. 

\noindent \textbf{Measurement error.}  We note that we have used multiple data sources and a conservative approach with multiple
criteria when we approach the problem of detecting banned and unbanned gambling content.
While we have demonstrated the validity and extremely high classification accuracy of our
approach, one may still concern about potential misclassifications which affect causal estimates.
In Appendix \ref{appendix_misclassification}, we conduct a formal sensitivity analysis and prove that we have correctly
identified ATT for the banned streamers.




\section{Descriptive Evidence and Empirical Strategy}\label{section-empirical-strategy}

\subsection{Descriptive Evidence}

Our final dataset comprises 155 banned streamers, 320 unbanned streamers, and 4,626 untreated streamers. Table \ref{table-summary} reports summary statistics for these streamers before the policy announcement. Streamers across the three groups exhibit similar average hours of total content and games with loot boxes. However, banned streamers devote substantially more time to gambling content than unbanned streamers and are also more popular on average, as reflected in total hours watched and subscriptions. These patterns suggest that banned and unbanned streamers may differ systematically, which raises identification challenges. We discuss these challenges and our empirical responses below.

\begin{table}[htb!]
\centering
\renewcommand{\arraystretch}{1.2}
\begin{threeparttable}
\resizebox{\textwidth}{!}{%
\begin{tabular}{lcccccc}
\hline
\multirow{2}{*}{} & \multirow{2}{*}{Mean} & \multirow{2}{*}{Frac.\ Zero} & \multirow{2}{*}{Gini} & \multicolumn{3}{c}{Group Mean} \\ 
\cline{5-7}
 & & & & Banned & Unbanned & Untreated \\ 
\hline
\textbf{Log Supply-Side Variables} & & & & & & \\
Streaming Hours   & 2.577 & 12.5\% & 0.266 & 2.900 & 2.636 & 2.573 \\
Gambling Hours    & 0.090 & 95.6\% & 0.968 & 1.608 & 0.898 & 0.000 \\
Loot Box Hours    & 1.438 & 43.1\% & 0.556 & 1.299 & 1.361 & 1.502 \\
Other Games Hours & 1.135 & 51.3\% & 0.637 & 0.487 & 0.566 & 1.134 \\
Non-gaming Hours  & 0.739 & 57.4\% & 0.701 & 0.991 & 0.976 & 0.727 \\
\textbf{Log Demand-Side Variables} & & & & & & \\
Total Hours Watched & 8.431 & 12.5\% & 0.203 & 9.858 & 8.782 & 8.421 \\
Tier 1 Subscription & 3.272 & 23.6\% & 0.388 & 3.508 & 2.569 & 3.424 \\
Tier 2 Subscription & 0.376 & 68.9\% & 0.771 & 0.210 & 0.167 & 0.429 \\
Tier 3 Subscription & 0.387 & 68.3\% & 0.767 & 0.215 & 0.159 & 0.434 \\
\hline
\end{tabular}%
}
\caption{\small \textbf{Summary Statistics for Streamers.} Summary statistics for the main sample of 5,101 Twitch streamers between the start date of the sample and the policy announcement date (August 1--September~20, 2022). The unit of observation is the streamer-week level. The definition of livestream types of the supply-side variables is discussed in Section \ref{subsection-data}. The \textit{Gini} coefficient is computed on the log-transformed outcome $y = \log(1 + x)$ as $G = \frac{2}{n\bar{y}}\sum_{i=1}^{n} i\,y_{(i)} - \frac{n+1}{n}$, where $y_{(i)}$ are the ascending order statistics.}
\label{table-summary}
\end{threeparttable}
\end{table}   

\subsection{Main Specification}\label{subsection-model}

Our primary goal is to estimate how streamers whose content was banned adjusted their content-production decisions, and how those changes affected viewer demand and streamer revenue. We are also interested in whether these effects extended to streamers who broadcast unbanned online gambling content and therefore faced a heightened risk of future sanctions. Our baseline specification is a two-way fixed-effects difference-in-differences (TWFE-DiD) model with streamer and week fixed effects:
\begin{flalign}
\begin{split}
    y_{it} = \alpha_i + \gamma_t &+ \beta_1 \cdot \text{Banned}_i \cdot \text{Post}_t + \beta_2 \cdot \text{Banned}_i \cdot \text{Announcement}_t\\
    &+ \beta_3 \cdot \text{Unbanned}_i \cdot \text{Post}_t + \beta_4 \cdot \text{Unbanned}_i \cdot \text{Announcement}_t + \varepsilon_{it}
    \label{twfe}
\end{split}
\end{flalign}
where $\text{Banned}_i$ and $\text{Unbanned}_i$ are indicators for the two treated groups, and $\text{Post}_t$ and $\text{Announcement}t$ indicate the post-implementation period (from October 18, 2022 onward) and the interim period between the policy announcement and implementation (from September 20, 2022 to October 18, 2022), respectively. We separate these two periods because streamers may adjust their content in response to the announcement itself, which would otherwise bias our estimates. On the supply side, the dependent variable $y{it}$ includes the log of weekly streaming hours for each content category described in Section \ref{section-data}, as well as the log of total weekly streaming hours. On the demand side, we mainly use the log of total hours watched and the log of subscriptions in each of Twitch’s three tiers. Our main coefficients of interest are $\beta_1$ and $\beta_3$, which capture the post-implementation effects for banned and unbanned streamers, respectively. We also estimate an event-study specification to trace the dynamics and persistence of the policy’s effects.

Nonetheless, this design faces several identification challenges. First, the four banned websites were not randomly selected.\footnote{Twitch announced that the four websites were banned due to “the lack of U.S. license or sufficient consumer protections.” Source: \url{https://safety.twitch.tv/s/article/Prohibiting-Unsafe-Slots-Roulette-and-Dice-Gambling-Sites?language=en\_US}.} Consequently, the groups of treated and untreated streamers may also not have been randomly assigned, potentially introducing a selection problem that could bias causal estimates. Table \ref{table-summary} is consistent with this concern. At the same time, we do not view this issue as fatal for three reasons. First, most online gambling websites offered highly similar content and followed similar pricing rules, making them close substitutes. Second, we identified dozens of smaller gambling websites that also appeared in Twitch livestreams and were likely unlicensed and hosted outside the United States. Third, the banned websites were not the four largest sites in terms of streaming popularity or market share. Taken together, these facts suggest that Twitch’s criteria were not unique to the four banned sites and that streamers and viewers associated with banned versus unbanned sites were unlikely to differ in fundamental ways other than the average level of reliance on gambling livestreams. To further address concerns over non-random treatment assignment, we conduct two sets of robustness checks: first, we use multiple matching methods to reduce systematic differences across streamers in different groups. Second, we instrument for the treatment indicators of banned and unbanned streamers using the share of banned gambling content and the share of gambling content in all livestreams, and then estimate an IV-DiD specification. Both approaches yield results similar to those from our baseline specification (see Appendix \ref{appendix-robustness}).

% Secondly, the classic two-way fixed effects DiD model usually assumes for parallel trend assumption (PTA) between streamers with different treatment status. Although our descriptive evidence suggests that different groups of streamers tend to follow similar trends over time, the assumption still remains questionable because streamers might pursue strategic changes in their streaming plan in response to the policy. The strategic interaction may potentially affect both supply- and demand-side outcome variables in our identification. 

Second, the policy may induce streamers to strategically reallocate their content, which in turn may lead viewers to substitute across channels in search of similar content. As a result, the banning policy may generate equilibrium spillovers in observed demand, potentially biasing the DiD estimates of demand-side outcomes upward and possibly affecting supply-side estimates as well. In this setting, using untreated streamers who share viewers with treated streamers as the control group raises a potential violation of the Stable Unit Treatment Value Assumption (SUTVA). As discussed in Section \ref{section-introduction}, however, the complexity of the viewer network and its possible evolution after the policy prevent us from cleanly estimating spillover effects themselves. Instead, we use network analysis to construct more comparable treated and untreated groups, thereby internalizing spillovers and assessing the robustness of our main results. We discuss this approach below.

\subsection{Network Analysis and Attention Spillover}\label{subsection-network}

To ensure that TWFE-DiD identifies the policy's causal impact, a key assumption is that the post-policy outcome of each streamer only depends on her own treatment, commonly referred to as SUTVA \citep{rubin1980randomization, rubin1990comment}. In reality,  viewers can switch between channels, resulting in potential attention spillovers in the livestreaming market, both within and across treated and untreated groups. This threatens the validity of both our supply- and demand-side estimates. Our strategy of mitigating the concern for SUTVA exploits the pattern of attention shift by leveraging the network structure of streamers to internalize demand spillovers.\footnote{Our context presents unique operational challenges in using existing methods aiming at identifying causal effects under network interference \citep{AronowSamii2017,li2022random} due to: (i) the complexity of the network structure, characterized by significant heterogeneity in streamer connectivity, and (ii) our treatment assignment is not random, nor can interference be assumed to arise randomly from unobserved streamer characteristics \citep{li2022random}.} 

We first construct a network of shared registered viewers among streamers prior to policy announcement. Each node represents a streamer and each edge indicates shared registered viewers, with edge weights equal to the number of shared viewers.\footnote{Registered viewers are a subset of the total viewers of a stream. Since Twitch does not provide individual viewer lists, we use registered viewers in chats as a proxy. See Appendix~\ref{data-collection} for discussions on validity.} To reduce network complexity which scales up with the number of streamers, the network covers all streamers from the two treated groups and a random sample of 800 (out of 4,626) untreated streamers.\footnote{Using a different random sample or varying the sample size for untreated streamers does not qualitatively affect our results.} After excluding a few streamers with missing data, the final network comprises $n$=1,270 streamers (155 banned, 315 unbanned, 800 untreated) with $n(n-1)/2 \simeq$ 805,815 edges. Figure~\ref{streamer-network} visualizes the network and reveals cross-group viewer overlap, suggesting potential interference. An interactive version of the same network is available at \url{https://twitch-gambling.github.io/streamer-network/network/} to aid readers with understanding.  


\begin{figure}[ht!]
    \centering
    \includegraphics[scale=0.16]{Figures/network/network-w-examples.png}
    \caption{\small \textbf{Network of streamers based on common viewers}.
    %This figure shows the network of streamers based on common viewers they share, before the policy announcement. The central panel shows the overall streamer network.
    The network is plotted by force-directed placement \citep{fruchterman1991graph}, which tries to place connected nodes close to each other and in doing so, helps reveal the basic structure of the network. A node denotes a streamer, and an edge indicates the presence of common viewers between two streamers. The size of a node is the number of viewers of a streamer, and the weight of an edge corresponds to the number of common viewers shared two streamers. The displayed network filters out edge weight that are below 1000 for visibility and identification of clusters. The threshold is selected heuristically based on the observed 5th quantile of unique viewers within a group as shown in Appendix Table \ref{table:network-group-statistics}. The nodes are colored such that red denotes the untreated group, blue denotes the unbanned (treated) group, and green denotes the banned (treated) group. See Appendix \ref{local-network} for examples of the communities.}
    \label{streamer-network}
\end{figure}


To mitigate interference, we combine community detection with a breath-first search based filtering algorithm to select large subgroups of closely connected, same-treatment-status streamers with minimal to no cross-group viewer overlap. First, we apply the Louvain method \citep{blondel2008fast} to partition the network
into 55 non-overlapping communities that maximize within-community shared viewership. These communities range from homogeneous (all members share one
treatment status) to highly mixed (see Appendix Table~\ref{table:detected-communities} and Figure~\ref{local-network} for examples). Selecting untreated subgroups is straightforward, as many communities consist entirely of untreated streamers. For banned and unbanned groups---which typically reside in mixed communities---we apply a Breadth-First Search procedure \citep{moore1959shortest}: starting from the most-connected streamer of the target
treatment status, we iteratively traverse neighbors and exclude those of different status, yielding a maximal connected subgroup of homogeneous
treatment before encountering cross-group streamers (see Algorithm~\ref{algorithm-1} and Appendix~\ref{appendix_network} for pseudo-code and
implementation details).

However, this procedure reduces interference at a cost. Because the network filter disproportionately removes streamers who are strongly connected to streamers in other treatment groups, the resulting selection is non-random. In practice, this procedure is more likely to exclude popular streamers who provide more diverse content and may respond more strategically to the policy. As a result, treatment effects estimated in the network-based sample may mechanically be smaller in magnitude than those from the main specification. In Section \ref{section-supply-side-het-results}, we show that this conjecture is consistent with the heterogeneous-treatment-effect evidence. One may also think of using the estimated network to identify spillover estimands in a network-DiD specification. However, we choose not to do so for several reasons. First, the network structure is highly complicated given the number of streamers in our sample, and the connection status are highly unbalanced. Second, because the network is constructed by chatters instead of all viewers, it likely reflects the directions of attention shift between streamers but fails to capture the real connection intensity in the network. Both features hurt the identifiability of existing estimation methods. Therefore, we use the network analysis to help internalize spillover effects and to demonstrate the robustness of our main specification.

  
% To reduce interference among streamers, our strategy combines community detection and breath-first search based filtering algorithms to automate the selection of large enough subgroups of closely connected streamers who share the same treatment status but have minimal to no overlap in viewers for each treatment group. We achieve this in 2 steps. First, we perform community detection using the Louvain method \citep{blondel2008fast} to identify non-overlapping communities that maximizes shared viewership within a community.\footnote{The method works by iteratively optimizing modularity, which measures how much more connected the nodes within communities are with respect to those outside communities. We also experimented the Girvan-Newman algorithm for community detection, but found Louvain to be more effective in unfolding communities for large networks.} Using the method, we identified 55 community clusters from the network. Appendix Table \ref{table:detected-communities} shows the shares of streamers across the three treatment groups for each community. Figure \ref{local-network} illustrates examples of different community clusters. In general, communities can be classified into three types: (1) \textit{homogeneous communities}, where all streamers belong to the same treatment status (e.g., community 14), (2) \textit{dominantly homogeneous communities}, where one treatment status is dominant (e.g., communities 10, 13, 18), and (3) \textit{high mixed communities}, where streamers of different treatment statuses are interwoven and there is no dominant treatment status (e.g., communities 8, 16, 32, 51). 

% While it is straightforward to select streamers for the untreated group (i.e., there are many homogeneous communities that consist entirely of untreated streamers and they are well separated from other groups), it is operationally challenging to select streamers for the banned and unbanned groups since most belong to communities of mixed treatment status. To address this, we propose a procedure to automate the selection of a large enough subgroup of closely connected streamers within a community for each treatment status. Given a community, we operationlize this by first identifying the streamer of a target treatment status who has the most edges to other streamers (i.e., the start node). We then apply filtering algorithm based on Breadth-First Search \citep{moore1959shortest} to iteratively traverse the network from the start streamer, search direct neighbors of the streamer, and exclude neighboring streamers who have a different treatment status from the streamer. For each community that can be operationally disentangled (see Appendix \ref{appendix_network} for details), we repeat this process iteratively through the layers until all streamers are processed. This procedure helps render a large enough subset of closely connected streamers of homogeneous status for each treatment group before encountering streamers of other status, thus reducing interference. We outline the pseudo-code in Algorithm \ref{algorithm-1} and illustrate the process intuitively in Figure \ref{bfs-streamer-filter} in the Appendix.  

\subsection{Additional Robustness Checks}

It has been well documented that baseline DiD specifications with logged outcome variables can generate substantial identification bias when outcome distributions differ sharply across groups \citep{mcconnell2024can}, for example when one group is highly skewed or when the fraction of zeros is large \citep{chen2024logs}. Because some of our outcomes exhibit these features, we estimate two additional specifications: a multiplicative model using Poisson pseudo-maximum-likelihood (PPML) estimation and a log specification that uses average concurrent viewers (ACV) in the pre-announcement period as weights for streamer size. Both alternatives yield qualitatively similar results (see Appendix \ref{appendix-robustness}).

\section{Results}\label{section-results}

\subsection{Supply-side Results}\label{section-supply-side-results}

We begin with estimates of the policy's impact on streamers' production from the baseline model. Because identification relies on the parallel trend assumption, Figure \ref{fig:es_supply} presents the event-study estimates as visual evidence. We further evaluate the assumption with linear pre-trend tests in Appendix. Although some pre-treatment coefficients are not perfectly flat, these deviations would, if anything, make our main estimates conservative rather than rejecting our qualitative conclusions.

\begin{figure}[!htb]
    \centering
    % Row 1: Banned streamers
    \begin{subfigure}[b]{0.32\textwidth}
        \centering
        \includegraphics[width=\textwidth]{Figures/event/Event_ban_log_slots_hrs.png}
        \caption{Gambling, Banned}
        \label{fig:es_supply_ban_gambling}
    \end{subfigure}
    \hfill
    \begin{subfigure}[b]{0.32\textwidth}
        \centering
        \includegraphics[width=\textwidth]{Figures/event/Event_ban_log_lootboxes_hrs.png}
        \caption{Loot Box, Banned}
        \label{fig:es_supply_ban_lootbox}
    \end{subfigure}
    \hfill
    \begin{subfigure}[b]{0.32\textwidth}
        \centering
        \includegraphics[width=\textwidth]{Figures/event/Event_ban_log_broadcast_hrs.png}
        \caption{Total Streaming, Banned}
        \label{fig:es_supply_ban_total}
    \end{subfigure}

    \vspace{0.3cm}
    % Row 2: Unbanned streamers
    \begin{subfigure}[b]{0.32\textwidth}
        \centering
        \includegraphics[width=\textwidth]{Figures/event/Event_unban_log_slots_hrs.png}
        \caption{Gambling, Unbanned}
        \label{fig:es_supply_unban_gambling}
    \end{subfigure}
    \hfill
    \begin{subfigure}[b]{0.32\textwidth}
        \centering
        \includegraphics[width=\textwidth]{Figures/event/Event_unban_log_lootboxes_hrs.png}
        \caption{Loot Box, Unbanned}
        \label{fig:es_supply_unban_lootbox}
    \end{subfigure}
    \hfill
    \begin{subfigure}[b]{0.32\textwidth}
        \centering
        \includegraphics[width=\textwidth]{Figures/event/Event_unban_log_broadcast_hrs.png}
        \caption{Total Streaming, Unbanned}
        \label{fig:es_supply_unban_total}
    \end{subfigure}
    \caption{\small \textbf{Time-varying treatment effects on content creation in both treated groups.} 1. Treatment effect in one week ahead of the policy announcement is normalized to zero. 2. The 95\% point-wise confidence intervals are illustrated as the inner bars, and the standard errors are clustered at the streamer level.}
    \label{fig:es_supply}
\end{figure}

The first column of Table \ref{ATT-streaming-hrs} reports the impact of the banning policy on gambling livestreams. We find that the supply of gambling livestreams decreased across the platform after the policy implementation, with banned streamers reducing their gambling streams by 65\% ($= \exp{(-1.049)} - 1$) and unbanned streamers reducing theirs by 10.2\%. These estimates have two implications: first, in addition to the banned gambling content, banned streamers reduced livestreams of unbanned gambling content as well as banned gambling content. This is because the overall share of banned streams among all gambling streams for these streamers is 55.8\%, which is smaller than the 65.0\% reduction caused by the policy. Second, even though Twitch only banned four websites, unbanned streamers also reduced their creation of gambling content even though they were not directly affected by the policy. Therefore, the banning policy was successful in regulating the overall supply of online gambling content across the platform.

\begin{table}[h!]
\renewcommand{\arraystretch}{1.2}
\centering
\resizebox{\textwidth}{!}{%
\begin{threeparttable}
\begin{tabular}{lccc}
%\hline
\toprule
 & log(Gambling Hours +1) & log(Loot Box Hours +1) & log(Streaming Hours +1) \\ 
%\hline
\midrule
Banned $\times$ Post ($\beta_1$) & \begin{tabular}[c]{@{}c@{}}-1.049***\\ (0.088)\end{tabular} 
                   & \begin{tabular}[c]{@{}c@{}}0.013\\ (0.066)\end{tabular} 
                   & \begin{tabular}[c]{@{}c@{}}-0.606***\\ (0.083)\end{tabular} \\

$\Delta \%$        & -65.0\% & – & -45.4\% \\

Unbanned $\times$ Post ($\beta_3$) & \begin{tabular}[c]{@{}c@{}}-0.108**\\ (0.049)\end{tabular} 
                     & \begin{tabular}[c]{@{}c@{}}-0.044\\ (0.043)\end{tabular} 
                     & \begin{tabular}[c]{@{}c@{}}-0.172***\\ (0.045)\end{tabular} \\

$\Delta \%$         & -10.2\% & – & -15.8\% \\ \hline

Streamer FE          & $\checkmark$ & $\checkmark$ & $\checkmark$ \\
Week FE              & $\checkmark$ & $\checkmark$ & $\checkmark$ \\
Observations         & 112,222 & 112,222 & 112,222 \\
Mean Dependent Variable & 1.130 & 1.341 & 2.722 \\ \hline
\end{tabular}
\end{threeparttable}%
}
\caption{\small \textbf{The effect of Twitch’s banning policy on supply-side outcomes.} All standard errors are clustered at the streamer level. The mean dependent variable is calculated based on observations of all treated (banned and unbanned) streamers before the policy announcement. FE = fixed effects. Significance levels: *$p<0.1$; **$p<0.05$; ***$p<0.01$.}
\label{ATT-streaming-hrs}
\end{table}

To answer whether potential substitution between online gambling and gambling-like features in video games weakened the effect of banning policy on Twitch, we report the estimated effect on the supply of video games featuring gambling-like elements in the second column of Table \ref{ATT-streaming-hrs}. The overall non-significant estimates suggest that on the supply side, gambling streamers did not switch to these video games to maintain their viewership on average. This finding suggests that these two types of content are less substitutable than some policymakers have presumed. We will discuss how streamers differed in substituting gambling with those video games in the next section.

Finally, we turn to examine the policy's impact on total content creation on the treated streamers. Estimates in the last two columns of Table \ref{ATT-streaming-hrs} show that both banned and unbanned streamers reduced their total streaming hours after the policy implementation. Our estimate imply a reduction of 45.4\% for banned streamers and 15.8\% for unbanned streamers. Given that gambling content accounts for 27.5\% and 17.6\%  of the total content for banned and unbanned streamers, respectively, the estimates imply that streamers reduced the production of non-gambling content by 28.2\% and 16.2\%, respectively. 

% The above last 4 numbers have not been revised yet. (April 4, 2026)

\subsection{Heterogeneous Treatment Effects}\label{section-supply-side-het-results}

We now explore whether the policy impact on content creation varies across different types of streamers, as examining the impact on subgroups of streamers help us better understand the policy's impact on regulating gambling content and investigate the underlying mechanisms.

\noindent \textbf{Reliance on Gambling Content.} First, we examine heterogeneity by a streamer’s reliance on gambling content, measured as the pre-announcement share of gambling streaming hours in total streaming hours. Figure \ref{HTE_supply_reliance_1} and \ref{HTE_supply_reliance_2} show a similar pattern across the two treated groups: streamers that relied more heavily on gambling content responded more strongly by reducing gambling livestreams. One might expect that removing some gambling websites would lower competition and induce substitution toward other gambling websites, but Table \ref{ATT-streaming-hrs} shows little evidence of such substitution on average. Figure \ref{HTE_supply_reliance_2} suggests that this mechanism is relevant only for the least gambling-reliant streamers in the unbanned group. A plausible explanation is that these streamers were less dependent on gambling as a primary income source and faced a lower perceived risk of future platform sanctions, making them less likely to reduce gambling content strategically. Turning to the streaming hours of games with loot boxes, although there is no average substitution from gambling to loot box games at the platform level in Table \ref{ATT-streaming-hrs}, the most gambling-reliant streamers appeared to increase streaming of loot box games, likely to offset part of the demand loss associated with the sharp reduction in gambling content.

% One new concern here is that the high reduction in unbanned gambling livestreams may be attributed to misclassifications.

\begin{figure}[htb!]
     \centering
     \begin{subfigure}[b]{0.4\textwidth}
         \centering
         \includegraphics[width=\textwidth]{Figures/HTE/plot_subgrp_slotsshare_log_slots_hrs_ban.png}
         \caption{Gambling, Banned}
         \label{HTE_supply_reliance_1}
     \end{subfigure}
     \begin{subfigure}[b]{0.4\textwidth}
         \centering
         \includegraphics[width=\textwidth]{Figures/HTE/plot_subgrp_slotsshare_log_slots_hrs_unban.png}
         \caption{Gambling, Unbanned}
         \label{HTE_supply_reliance_2}
     \end{subfigure}
     \vspace{-0.3cm}
     \begin{subfigure}[b]{0.4\textwidth}
         \centering
         \includegraphics[width=\textwidth]{Figures/HTE/plot_subgrp_slotsshare_log_lootboxes_hrs_ban.png}
         \caption{Total Livestreams, Banned}
         \label{HTE_supply_reliance_3}
     \end{subfigure}
     \begin{subfigure}[b]{0.4\textwidth}
         \centering
         \includegraphics[width=\textwidth]{Figures/HTE/plot_subgrp_slotsshare_log_lootboxes_hrs_unban.png}
         \caption{Total Livestreams, Unbanned}
         \label{HTE_supply_reliance_4}
     \end{subfigure}
     \vspace{0.5cm}
     \caption{\small \textbf{Heterogeneous treatment effects on streamers with different pre-announcement shares of gambling livestreams}. 1. The coefficients in each subfigure show the point estimates of $\beta_1$ in specification \ref{twfe} for banned streamers and $\beta_3$ for unbanned streamers. 2. The streamer quartiles are derived based on the share of gambling livestreams over all livestreams before the policy announcement. 3. The 95\% point-wise confidence intervals are illustrated as the inner bars, and the standard errors are clustered at the streamer level.}
     \label{HTE_supply_reliance}
\end{figure}

\noindent \textbf{Streamer Popularity.} Next, we explore whether the policy impact on content creation varies across streamers with different levels of popularity. Figure \ref{HTE_supply_popularity} shows that banned streamers with higher popularity reduced their content creation of gambling livestreams slightly more than streamers with low popularity after the policy. More interestingly, only streamers with the highest level of popularity in the unbanned group significantly reduced their content creation of gambling livestreams, with a magnitude of approximately 34.8\%. The policy impact on total livestreaming hours also appears to increase with streamer popularity: while streamers with low popularity almost did not respond to the policy, streamers with the highest popularity contributed most to the overall reduction in content creation of total gambling livestreams. Since streamers with higher popularity had higher average streaming hours, our result suggests that the policy led to a large negative impact on the total content creation across the platform.

\begin{figure}[h!]
     \centering
     \begin{subfigure}[b]{0.4\textwidth}
         \centering
         \includegraphics[width=\textwidth]{Figures/HTE/plot_hdid_log_slots_hrs_ban.png}
         \caption{Gambling, Banned}
         \label{fig:token1}
     \end{subfigure}
     \begin{subfigure}[b]{0.4\textwidth}
         \centering
         \includegraphics[width=\textwidth]{Figures/HTE/plot_hdid_log_slots_hrs_unban.png}
         \caption{Gambling, Unbanned}
         \label{fig:token2}
     \end{subfigure}
     \vspace{-0.3cm}
     \begin{subfigure}[b]{0.4\textwidth}
         \centering
         \includegraphics[width=\textwidth]{Figures/HTE/plot_hdid_log_broadcast_hrs_ban.png}
         \caption{Total Livestreams, Banned}
         \label{fig:token3}
     \end{subfigure}
     \begin{subfigure}[b]{0.4\textwidth}
         \centering
         \includegraphics[width=\textwidth]{Figures/HTE/plot_hdid_log_broadcast_hrs_unban.png}
         \caption{Total Livestreams, Unbanned}
         \label{fig:token4}
     \end{subfigure}
     \vspace{0.5cm}
     \caption{\small \textbf{Heterogeneous treatment effects on streamers with different popularity}. 1. The coefficients in each subfigure show the point estimates of $\beta_1$ in specification \ref{twfe} for banned streamers and $\beta_3$ for unbanned streamers. 2. The streamer popularity quartiles are derived based on average concurrent viewership before the policy announcement. 3. The 95\% point-wise confidence intervals are illustrated as the inner bars, and the standard errors are clustered at the streamer level.}
     \label{HTE_supply_popularity}
\end{figure}

How can we explain the findings that more popular streamers were affected more by the policy? One potential hypothesis is that these streamers may care more about their reputation. On the one hand, Twitch regulated gambling livestreams because it was controversial and raised reputation concerns for the platform. If a highly popular streamer continues providing a large amount of gambling content but is later prohibited by the platform, they will lose more revenue due to their high number of subscriptions, in-stream donations and total hours watched. On the other hand, reducing gambling content and showing support to the platform's policy can enhance their reputation within the Twitch community. These factors may incentivize popular streamers to behave conservatively, resulting in a greater reduction of gambling streams.

To test this hypothesis, we proxy conservativeness of each streamer by the number of days its account was prohibited in the pre-announcement periods. In our sample, banned streamers were slightly more likely to be prohibited before the policy announcement: they averaged 0.634 prohibited days compared to 0.447 for unbanned streamers and 0.177 for untreated streamers. They also received slightly harsher punishments per prohibition, averaging 5.281 total banned days (8.330 days per prohibition) versus 3.130 days (7.002 days per prohibition) for unbanned streamers and 1.160 days (6.554 days per prohibition) for untreated streamers. If streamers value their reputation, they are more likely to carefully plan their streaming content to avoid temporary prohibitions from the platform. We run a DiD specification that includes interactions between treatment indicators and the number of days of account prohibition each streamer received before the policy announcement. Column 1 and 2 in Table \ref{mech} present our estimation results. In addition to the negative policy impact on gambling streams, we find that both banned and unbanned streamers reduced gambling livestreams less if they cared less about their reputation (reflected as having more days of prohibitions before the policy implementation). Specifically, a banned streamer with no prior account prohibitions reduced gambling streams by 64.6\%, and each additional day of prior prohibition resulted in a 0.4\% smaller reduction compared to the baseline. Moreover, unbanned streamers are less sensitive to streamer reputation, with each additional prior prohibition resulting in a 0.6\% smaller reduction, on top of a baseline decrease of -13.1\%. We observe similar effects when we replace the number of times prohibited by the total number of days prohibited before the policy implementation.

\begin{table}[h!]
\renewcommand{\arraystretch}{1.2}
\centering
\begin{threeparttable}
\begin{tabular}{lcc}
\hline
             & \multicolumn{2}{c}{log(Gambling Hours+1)} \\ \cline{2-3}
             & Main Effect       & Interaction Effect       \\ \hline
Banned $\times$ Post ($\beta_1$) $\times$ Banned Days       & \begin{tabular}[c]{@{}c@{}}-1.069***\\ (0.090)\end{tabular}  
                         & \begin{tabular}[c]{@{}c@{}}0.004* \\ (0.002)\end{tabular}      \\
Unbanned $\times$ Post ($\beta_3$) $\times$ Banned Days     & \begin{tabular}[c]{@{}c@{}}-0.126** \\  (0.050)\end{tabular}   
                         & \begin{tabular}[c]{@{}c@{}}0.006***\\  (0.002)\end{tabular}     \\ \hline
Streamer FE  & $\checkmark$               & $\checkmark$                \\
Week FE      & $\checkmark$               & $\checkmark$                \\
Observations & 112,222                    & 112,222                     \\
Mean Dependent Variable & 1.130           & 1.130                       \\
\hline 
\end{tabular}
\end{threeparttable}
\caption{\small \textbf{Underlying mechanisms of heterogeneous treatment effects on gambling livestreams}. All standard errors are clustered at the streamer level. The mean dependent variable is calculated based on observations of all treated streamers before the policy announcement. FE = fixed effects. Significance level: *$p<0.1$; **$p<0.05$; ***$p<0.01$}
\label{mech}
\end{table}

Another potential explanation for the positive correlation between streamer popularity and the regulatory effect is that more popular streamers may have been more likely to engage in multi-homing. This would have reduced their cost of transitioning gambling livestreams to other platforms and facilitated the migration of their gambling audiences. However, popular streamers were also more likely to be Twitch partners who held exclusive contracts that restricted them from streaming on competing platforms such as YouTube or Facebook. Therefore, we cannot draw a definitive conclusion about the role of multi-homing in driving the heterogeneity in treatment effects observed in our analysis.

\subsection{Demand-side Results}\label{section-demand-side-results}

Next, we present estimation results on key demand-side outcome variables from the full sample. We first focus on two sets of demand-side treatment effects: the effect on total hours watched, and the effect on the three tiers of subscriptions. Total hours watched is one of the most crucial metrics for assessing channel popularity and viewer engagement.\footnote{Twitch uses total hours watched to identify streamers eligible for its Streamer Achievement Program. See \url{https://help.twitch.tv/s/article/bleed-purple?language=en\_US}.} It is also closely linked to Twitch’s advertising revenue, as many Twitch streamers adopt mid-roll advertisements while they stream. Subscriptions, in contrast, represent a consistent and direct income source for streamers, who receive at least 50\% of the revenue from each subscription.

Table \ref{ATT_hrs_watched} presents estimated causal effects on total hours watched for all livestreaming categories. Banned streamers experienced an 83.2\% decrease in average weekly viewership, while unbanned streamers saw a 38.6\% decline. These estimates are largely comparable to Twitch's reported 75\% reduction in gambling viewership.\footnote{Source: \url{https://www.nbcnews.com/tech/internet/twitch-expands-ban-gambling-livestreams-also-says-viewership-content-7-rcna97943.}} Comparing these estimates with the 45.4\% and 15.8\% declines in total streaming hours reported in Table \ref{ATT-streaming-hrs} indicates that streamers experienced disproportionately larger losses in content consumption than in content creation, suggesting that their pre-policy audiences had stronger preferences for gambling content. Both treated groups exhibit a similar proportional decline in content consumption relative to content creation, implying consistent market responses.

\begin{table}[htb!]
\renewcommand{\arraystretch}{1.2}
\centering
\begin{threeparttable}
\begin{tabular}{lccccc}
\hline 
 & \multicolumn{5}{c}{log(Hours Watched + 1)} \\ \cline{2-6} 
 & Gambling & Overall & Loot Box & Other Games & Non-gaming \\ \hline
Banned $\times$ Post ($\beta_1$)       
& \begin{tabular}[c]{@{}c@{}}-4.207*** \\ (0.313)\end{tabular} 
& \begin{tabular}[c]{@{}c@{}}-1.783*** \\ (0.264)\end{tabular} 
& \begin{tabular}[c]{@{}c@{}}0.035 \\ (0.226)\end{tabular} 
& \begin{tabular}[c]{@{}c@{}}-0.317** \\ (0.153)\end{tabular} 
& \begin{tabular}[c]{@{}c@{}}-0.385* \\ (0.209)\end{tabular} \\
$\Delta \%$             
& -98.5\% 
& -83.2\% 
& – 
& -27.2\% 
& -32.0\% \\
Unbanned $\times$ Post ($\beta_3$)     
& \begin{tabular}[c]{@{}c@{}}-0.693*** \\ (0.180)\end{tabular} 
& \begin{tabular}[c]{@{}c@{}}-0.488*** \\ (0.136)\end{tabular} 
& \begin{tabular}[c]{@{}c@{}}-0.125 \\ (0.147)\end{tabular} 
& \begin{tabular}[c]{@{}c@{}}-0.410*** \\ (0.108)\end{tabular} 
& \begin{tabular}[c]{@{}c@{}}-0.318*** \\ (0.117)\end{tabular} \\
$\Delta \%$             
& -50.0\% 
& -38.6\% 
& – 
& -33.6\% 
& -27.2\% \\ \hline
Streamer FE  & $\checkmark$ & $\checkmark$ & $\checkmark$ & $\checkmark$ & $\checkmark$ \\
Week FE      & $\checkmark$ & $\checkmark$ & $\checkmark$ & $\checkmark$ & $\checkmark$ \\
Observations & 112,222 & 112,222 & 112,222 & 112,222 & 112,222 \\
Mean Dependent Variable 
& 4.450 & 9.133 & 4.837 & 2.265 & 4.436 \\ 
\hline 
\end{tabular}
\end{threeparttable}
\caption{\textbf{The effect of Twitch’s banning policy on hours watched.} All standard errors are clustered at the streamer level. The mean dependent variable is calculated based on observations of all treated (banned and unbanned) streamers before the policy announcement. FE = fixed effects. Significance level: *$p<0.1$; **$p<0.05$; ***$p<0.01$.}
\label{ATT_hrs_watched}
\end{table}

Next, we examine the source of decline in total hours watched by checking the policy's impacts on total hours watched for each livestream category. Banned streamers experienced a 98.5\% reduction in total hours watched of their gambling livestreams, whereas unbanned streamers had a reduction of approximately 50.0\%. Moreover, consumption of livestreams of other video games and non-gaming content also contributed to the decline in total hours watched, as they received around 27\%-33\% reduction after the policy implementation. 

While hours watched reflects viewer engagement and content popularity, it does not directly indicate the revenue received by each streamer. Therefore, we also examine the policy impact on three tiers of subscriptions on Twitch. These outcome variables are informative for two reasons: first, the three tiers of subscriptions directly reflect changes in streamers' total revenue, as a Tier 1 subscription is worth \$4.99, Tier 2 is worth \$9.99 and Tier 3 is worth \$24.99. Second, these subscriptions generally do not include exclusive livestreaming content\footnote{Viewers who purchase higher-tier subscriptions to a streamer enjoy benefits such as access to additional in-chat emotes, tenure-based chat badges, and subscriber-only chats.}. Therefore, a higher-tier subscription can be seen as a signal of viewer loyalty and engagement in the streamer's community. Table \ref{ATT_subs} shows that only Tier 1 subscriptions were significantly affected by the banning policy, with banned streamers experiencing a 45.4\% loss in weekly Tier 1 subscriptions after the policy and unbanned streamers experiencing a loss of -15.8\%, suggesting a huge loss in total revenue for the treated streamers due to the high number of Tier 1 subscriptions on average. In contrast, we find marginally significant to insignificant estimates of policy impact on both Tier 2 and Tier 3 subscriptions. Since viewers pay much more for higher tiers of subscriptions, only a small amount of loyal viewers are willing to support the streamer at these levels. This finding suggests that, on average, it was unlikely that either banned or unbanned streamers suffered significant losses in their loyal viewers or engagement from their core communities. By comparing the reduction in total hours watched and subscriptions, we also find that viewership is approximately twice as elastic as subscription consumption.

\begin{table}[h!]
\renewcommand{\arraystretch}{1.2}
\centering
\begin{threeparttable}
\begin{tabular}{lccc}
\hline 
             & log(Tier 1+1)           & log(Tier 2+1)        & log(Tier 3+1)         \\ \hline 
Banned $\times$ Post ($\beta_1$)       & \begin{tabular}[c]{@{}c@{}}-0.605***\\ (0.122)\end{tabular} & \begin{tabular}[c]{@{}c@{}}0.008 \\ (0.013)\end{tabular}  & \begin{tabular}[c]{@{}c@{}}-0.027* \\(0.016)\end{tabular} \\
$\Delta \%$             & -45.4\%          & -             & 2.4\%              \\
Unbanned $\times$ Post ($\beta_3$)     & \begin{tabular}[c]{@{}c@{}}-0.172*** \\  (0.057)\end{tabular} & \begin{tabular}[c]{@{}c@{}}0.013 \\ (0.011)\end{tabular} & \begin{tabular}[c]{@{}c@{}}0.014\\ (0.012)\end{tabular}  \\
$\Delta \%$             & -15.8\%          & -             & -              \\ \hline
Streamer FE  & $\checkmark$              & $\checkmark$           & $\checkmark$            \\
Week FE      & $\checkmark$              & $\checkmark$           & $\checkmark$            \\
Observations & 112,222          & 112,222       & 112,222        \\
Mean Dependent Variable & 2.876 & 0.181 & 0.178 \\ 
\hline 
\end{tabular}
\end{threeparttable}
\caption{\textbf{The effect of Twitch's banning policy on different tiers of subscriptions}. All standard errors are clustered at the streamer level. The mean dependent variable is calculated based on observations of all treated (banned and unbanned) streamers before the policy announcement. FE = fixed effects. Significance level: *$p<0.1$; **$p<0.05$; ***$p<0.01$.}
\label{ATT_subs}
\end{table}

\subsection{Viewer Engagement Results}\label{subsection-total-chat-engagement}
In addition to subscriptions, we also examine viewer engagement through chat activity, including messages and emotes. To study how the policy affected chat engagement, we construct daily measures at both the overall and category levels, including chat messages per hour, unique chatters per hour, average emotes per message, average message length in characters, and average token count per message.\footnote{To ensure that some of these measures are not affected by the total number of streaming hours, we divide them by hour to obtain per hour engagement in chats and use those for analysis.}  

Table \ref{table-summary-engagement} reports unconditional and group means of our chat engagement metrics. The first row shows that banned and unbanned streamers receive higher engagement than untreated streamers from their viewers in general, as they have more chatters (168.6 and 104.9 in levels) and receive more chat messages (3886.4 and 2470.0 in levels) per hour of livestream. Chat messages also differ in their composition. Viewers of treated streamers are more likely to include emotes in their chat messages (2.489 and 2.283 per message for banned and unbanned versus 1.867 for untreated). This may be because gambling content elicits excitement more frequently than other content, and thus prompt more (often duplicated) emotes from viewers. In contrast, we observe no major differences across treated and untreated groups in the average number of characters (around 24–25) or tokens (around 4–4.5) per chat message.

\begin{table}[!htb]
\renewcommand{\arraystretch}{1.2}
\centering
\begin{threeparttable}
\begin{tabular}{lcccc}
\hline
\multirow{2}{*}{}              & \multirow{2}{*}{Mean} & \multicolumn{3}{c}{Group Mean}                                                            \\ \cline{3-5} 
                               &                       & \multicolumn{1}{l}{Banned} & \multicolumn{1}{l}{Unbanned} & \multicolumn{1}{l}{Untreated} \\ \hline
\textbf{Engagement Volume Metrics} & \multicolumn{1}{l}{}  & \multicolumn{1}{l}{}       & \multicolumn{1}{l}{}         & \multicolumn{1}{l}{}          \\
Chats per Hour                & 5.659                 & 7.214                      & 6.416                        & 5.869                         \\
Chatters per Hour             & 2.975                 & 4.026                      & 3.336                        & 3.091                         \\
\textbf{Engagement Intensity Metrics} &                       &                            &                              &                               \\
Avg. Emotes            & 1.802                 & 2.489                      & 2.283                        & 1.867                         \\
Avg. Character Count            & 23.914                 & 24.916                      & 24.285                        & 24.746                         \\
Avg. Token Count            & 4.174                 & 4.482                      & 4.208                        & 4.309                         \\ \hline
\end{tabular}
\end{threeparttable}
\small
\vspace{0.5em}
\caption{\small \textbf{Summary statistics of engagement metrics.} 1. Summary statistics for the main sample of 5,101 Twitch streamers between the start date of the sample and the policy announcement date (August 1 - September 20, 2022). The unit of observation is at the streamer-weak level. 2. Engagement volume metrics are transformed as $\log(1+y)$, and engagement intensity metrics are measured per chat message.}
\label{table-summary-engagement}
\end{table}

Table \ref{ATT_engagement} summarizes the estimated policy effects on engagement metrics from the TWFE-DiD regressions. The first column shows large, significant reductions in the hourly number of chat messages in both banned streamers' livestreams (-69.3\%) and unbanned streamers' livestreams (-28.0\%). Although substantial, these reductions are smaller than the declines in total hours watched reported in Table \ref{ATT_hrs_watched}. Column 2 further shows that the total number of chatters also fell after the policy implementation for treated streamers, and the percentage change is smaller than that for chat messages. These patterns are consistent with our demand-side results, as remained viewers are expected to be more loyal and engaged, resulting in lower decline in chat messages. Beyond extensive-margin effects on viewership volume, the banning policy also reduced engagement intensity within each chat message (intensive margin). Columns 3 to 5 show that each chat message received by banned streamers in the post-policy period contained 0.443 fewer emotes, 2.767 fewer words and 0.504 fewer tokens on average. Messages to unbanned streamers were also less engaged, with reductions of 0.150 emotes, 0.890 characters, and 0.167 tokens per message. Together, these estimates provide valuable evidence that platform policies may affect both the volume of engagement and its intensity within messages.

\begin{table}[!h]
\renewcommand{\arraystretch}{1.2}
\centering
\begin{threeparttable}
\resizebox{\textwidth}{!}{%
\begin{tabular}{lccccc}
\hline 
& Chats per Hour & Chatters per Hour & Avg Emotes & Avg Character Count & Avg Token Count \\ \hline 
Banned $\times$ Post ($\beta_1$) & 
\begin{tabular}[c]{@{}c@{}}-1.182*** \\ (0.189)\end{tabular} &
\begin{tabular}[c]{@{}c@{}}-1.200*** \\ (0.175)\end{tabular} &
\begin{tabular}[c]{@{}c@{}}-0.443*** \\ (0.077)\end{tabular} &
\begin{tabular}[c]{@{}c@{}}-2.767*** \\ (0.761)\end{tabular} &
\begin{tabular}[c]{@{}c@{}}-0.504*** \\ (0.152)\end{tabular} \\
$\Delta \%$ & -69.3\% & -59.0\% & - & - & - \\
Unbanned $\times$ Post ($\beta_3$) &
\begin{tabular}[c]{@{}c@{}}-0.329*** \\ (0.102)\end{tabular} &
\begin{tabular}[c]{@{}c@{}}-0.323*** \\ (0.090)\end{tabular} &
\begin{tabular}[c]{@{}c@{}}-0.150*** \\ (0.051)\end{tabular} &
\begin{tabular}[c]{@{}c@{}}-0.890** \\ (0.431)\end{tabular} &
\begin{tabular}[c]{@{}c@{}}-0.167** \\ (0.077)\end{tabular} \\
$\Delta \%$ & -28.0\% & -20.4\% & - & - & - \\ \hline
Streamer FE  & $\checkmark$ & $\checkmark$ & $\checkmark$ & $\checkmark$ & $\checkmark$ \\
Week FE      & $\checkmark$ & $\checkmark$ & $\checkmark$ & $\checkmark$ & $\checkmark$ \\
Observations & 112,222 & 112,222 & 112,222 & 112,222 & 112,222 \\
Mean Dependent Variable & 6.677 & 3.561 & 2.350 & 24.491 & 4.298 \\ 
\hline 
\end{tabular}%
}
\end{threeparttable}
\caption{\textbf{The effect of Twitch's banning policy on engagement metrics}. All standard errors are clustered at the streamer level. The mean dependent variable is calculated based on observations of all treated (banned and unbanned) streamers before the policy announcement. FE = fixed effects. Significance level: *$p<0.1$; **$p<0.05$; ***$p<0.01$.}
\label{ATT_engagement}
\end{table}

\subsection{Results on Selected Groups}

Next, we present supply- and demand-side results on our subsample selected by the network analysis procedure. In Table \ref{ATT_net}, we summarize causal estimates for the same outcome variables as in Table \ref{ATT_hrs_watched} and \ref{ATT_subs} based on streamers selected from the network analysis, and compare them with previously reported results. We find that all effects are qualitatively similar to previously reported estimates, except that the DiD estimator at the aggregated level indicates a slight increase in Tier 3 subscriptions for unbanned streamers. However, this result does not change our conclusion that neither banned nor unbanned streamers suffered from losses in their core viewership, where both groups experienced significant loss in revenue from non-loyal viewers. Although we expect that the DiD estimates on our full streamer sample will overestimate the ATT on the total hours watched, the estimates in Table \ref{ATT_net} suggest that previous results indeed underestimate the negative impact on the affected streamers. One explanation to the results is that since the selected streamers share much fewer viewers with others on the platform, they are in general less popular and can be replaced by other streamers easily when they were forced not to change their streaming content, resulting in a higher effect on these streamers.

% Finalized on Sept 22, 2025.
\begin{table}[h!]
\renewcommand{\arraystretch}{1.2}
\centering
\resizebox{\textwidth}{!}{%
\begin{threeparttable}
\begin{tabular}{lcccc}
\hline 
             & log(Hours Watched + 1) & log(Tier 1 + 1) & log(Tier 2 + 1) & log(Tier 3 + 1) \\ \hline
Banned $\times$ Post ($\beta_1$) & \begin{tabular}[c]{@{}c@{}}-2.163***\\ (0.547)\end{tabular} 
                                 & \begin{tabular}[c]{@{}c@{}}-0.614**\\ (0.280)\end{tabular} 
                                 & \begin{tabular}[c]{@{}c@{}}0.018\\ (0.023)\end{tabular} 
                                 & \begin{tabular}[c]{@{}c@{}}-0.021\\ (0.035)\end{tabular} \\
$\Delta \%$                     & -88.5\% & -45.9\% & – & – \\[0.5em]
Unbanned $\times$ Post ($\beta_3$) & \begin{tabular}[c]{@{}c@{}}-0.764***\\ (0.244)\end{tabular} 
                                   & \begin{tabular}[c]{@{}c@{}}-0.214**\\ (0.102)\end{tabular} 
                                   & \begin{tabular}[c]{@{}c@{}}0.017\\ (0.020)\end{tabular} 
                                   & \begin{tabular}[c]{@{}c@{}}0.040**\\ (0.018)\end{tabular} \\
$\Delta \%$                     & -53.4\% & -19.3\% & – & 4.1\% \\ \hline
Streamer FE  & $\checkmark$ & $\checkmark$ & $\checkmark$ & $\checkmark$ \\
Week FE      & $\checkmark$ & $\checkmark$ & $\checkmark$ & $\checkmark$ \\
Observations & 10,494 & 10,494 & 10,494 & 10,494 \\ \hline
\end{tabular}
\end{threeparttable}%
}
\caption{\small \textbf{The impact of Twitch's banning policy on demand-side outcomes, based on selected groups of streamers.} All standard errors are clustered at the community level. FE = fixed effects. Significance level: *$p<0.1$; **$p<0.05$; ***$p<0.01$.}
\label{ATT_net}
\end{table}

In Table \ref{table:full-estimation-results-isolated-streamers} of Appendix \ref{appendix_full_results}, we present full estimation results for both demand-side and supply-side variables among the selected groups of streamers. This additional analysis addresses concerns regarding potential violations of SUTVA on the demand-side outcomes influencing streamers' decisions on content production. Overall, the findings corroborate those presented in Section \ref{section-supply-results}. Notably, the streaming hours for gambling content among unbanned streamers in the selected group were unaffected by the policy. This is expected because the selected streamers, who are less connected to those in different treated groups, typically have lower popularity, making them less responsive to the banning policy, as we displayed in Figure \ref{HTE_supply_popularity}. All other results remain qualitatively the same as those reported in Table \ref{table:full-supply-results} and Table \ref{table:full-demand-results}.

Our construction of streamer communities also allows us to fully internalize potential interference within each community by aggregating observations at the community level. Hence, we report the full estimation results based on community-level observations in Table \ref{table:full-estimation-results-isolated-streamers-community} of Appendix \ref{appendix_full_results}. Note that the number of observations (946) shrinks to the product of the number of distinct communities (43) and the number of weeks (22). We find qualitatively the same results both in terms of non-monetary content consumption (total hours watched) and monetary consumption (subscriptions) from the community-level regressions.

\section{Minor Viewer Analysis}\label{section-minor}
\subsection{Detection of Minor Content and Viewers}
To understand whether the gambling policy might be effective at protecting minors (viewers below 18), we need to first understand which content attracts a large audience of minors and then identify viewers who watch a lot of content targeted for minors. To answer this question, we first identify content more likely to be for minors using LLMs, given their strong performance for diverse downstream tasks. We start by running a pilot classification to test performances of  different model configurations from OpenAI on 50 randomly selected content (see Appendix \ref{appendix-minor-content-pilot} for details). We find that GPT o3 with web search, known for superior capability of logical tasks, offers the most conservative and accurate labels for content classification. Next, we run the model through the 10,970 content covered by our sample using a 2-step procedure. In step 1, we classify which content attracts large audience of minors using the same prompts as in pilot. We find that 4.1\% (451) content attracts minors, however, with some having minors as vast majority (e.g., Minecraft) and others appealing to much broader audience including minors (e.g., FIFA, NBA 2K). Thus, in step 2, we ask the LLM to evaluate the share of minor viewers for these 451 games on a scale of 0 to 100. We find that using a threshold of 50 yields 119 clearly more minor-targeted games while excluding popular games that attracts a more diverse viewer base. See Appendix \ref{appendix-minor-content-main} for more details.  

We use labeled games to identify viewers who are more likely to be minors from the 14,568,807 registered viewers in our sample. For this, we examine whether the user's viewership is concentrated on content classified as for minors. Appendix Table \ref{table:viewer-stats-pre-policy} shows the summary statistics on total sessions and sessions with minor content watched by viewers.\footnote{We define a session at the streamer-stream-content level, and a viewer is considered to have watched a content session if he appears at least once in that session. Using the tracked sessions watched by each viewer, we construct the total number of sessions watched and minor sessions watched (including the share) by matching content streamed in a session with our constructed minor content list.} On average 0.05\% sessions watched across all viewers featured minor content, and 15\% of viewers watched at least one minor session pre-policy. For these viewers, on average minor sessions cover 35\% of their viewing history. We show in Appendix Figures \ref{distribution-session-watched} and \ref{distribution-minor-share-watched} that given that some viewers may watch very few sessions in total, it is not ideal to use raw minor shares to proxy for minor viewers. For example, if a viewer only watches two content sessions and one session happens to be minor content, then the share would be 50\%. To account for this issue, instead, we classify viewers into likely minors, likely adults or uncertain by conducting two one-sided binomial tests (see Appendix \ref{appendix-minor-detection} for details). We refer to these groups as minors, adults, and uncertain-age viewers below. We show that the tests weight evidence by the number of watched sessions, preventing false minor and adult labels driven by noisy shares from low watching activity, while still flagging viewers whose behavior is statistically significant from the mean. With this, we obtain 5.99\% (872,672) viewers classified as minors, 4.42\% (643,941) as adults, and 89.59\% (13,052,194) as uncertain for subsequent analysis. 

% We need clarify that we use identified viewers and sessions to measure total viewership for each age group.

\subsection{Viewer Analysis Results}\label{subsection-minor-results}

We first examine whether users classified as minors were exposed to any gambling content on Twitch. Appendix Table \ref{table:summary-age-category-panels-all}(a) presents the average number of gambling and non-gambling views by each age group for all streamers on Twitch before the policy announcement.\footnote{Appendix Table \ref{table:summary-age-category-panels-treated} presents similar summary statistics for streamers treated by the policy. All comparisons are qualitatively the same.} Viewers classified as minors have lower exposure to gambling content than adults: on average, streamers had 5.8 minor viewers who watched their gambling sessions per week, compared to 58.1 adult viewers per week.\footnote{We use the number of views by each age group in this section instead of the total hours watched in Section \ref{section-demand-results}. This is because we cannot accurately identify a viewer's entry and exit of a streaming session from chat messages, making the measurement of each viewer's watching duration very imprecise.} This lower exposure to gambling content is partially due to minors' lower exposure to any content, with an average streamer having 1422.1 adults and 552.1 minors per week who watched their sessions. Conditional on watching gambling streams, minors are only slighty less engaged with gambling content (parts (b)-(e) of Appendix Table \ref{table:summary-age-category-panels-all}). These summary statistics show that minors were exposed to and consumed gambling content on Twitch less than adults, but still had some exposure in the pre-policy period.

%(110.1 minors for banned streamers and 42.7 minors for unbanned streamers)
%(1241.1 adults for banned streamers and 268.5 adults for unbanned streamers) 

We now turn to the question of whether the banning policy had a stronger or weaker effect on minor viewers compared to non-minors. For each age group, we estimate specification (\ref{twfe}) separately. We report the results for viewership of both gambling and overall content for banned streamers in Table \ref{ATT_log_num_views}. First, the banning policy led to substantial declines in viewership for banned streamers across all age groups. For gambling livestreams, the number of minor viewers reduced by approximately 78.9\%. However, this effect was smaller in magnitude compared to the reductions observed among identified adults (94.0\%) and uncertain-age viewers (94.9\%). These estimates are consistent with our findings reported in Table \ref{ATT_hrs_watched}. Declines in total viewership differed more prominently across age groups: while minor viewers only decreased their total weekly consumption of banned streamers' content by 25.1\%, adults and uncertain-age viewers reduced their consumption by 47.9\% and 40.3\%.

\begin{table}[!htb]
\renewcommand{\arraystretch}{1.2}
\centering
\resizebox{\textwidth}{!}{%
\begin{threeparttable}
\begin{tabular}{lcccccc}
\toprule
 & \multicolumn{3}{c}{log(Number of Gambling Views + 1)} & \multicolumn{3}{c}{log(Number of Total Views + 1)} \\
\cmidrule(lr){2-4} \cmidrule(lr){5-7}
 & Minor & Adult & Uncertain & Minor & Adult & Uncertain \\
\midrule
Banned $\times$ Post ($\beta_1$) & 
\begin{tabular}[c]{@{}c@{}}-1.557***\\ (0.183)\end{tabular} & 
\begin{tabular}[c]{@{}c@{}}-2.819***\\ (0.216)\end{tabular} & 
\begin{tabular}[c]{@{}c@{}}-2.981***\\ (0.237)\end{tabular} & 
\begin{tabular}[c]{@{}c@{}}-0.290***\\ (0.054)\end{tabular} & 
\begin{tabular}[c]{@{}c@{}}-0.652***\\ (0.073)\end{tabular} & 
\begin{tabular}[c]{@{}c@{}}-0.516***\\ (0.066)\end{tabular} \\
$\Delta \%$ & -78.9\% & -94.0\% & -94.9\% & -25.1\% & -47.9\% & -40.3\% \\
\addlinespace[0.3em]
Unbanned $\times$ Post ($\beta_3$) & 
\begin{tabular}[c]{@{}c@{}}-0.293***\\ (0.067)\end{tabular} & 
\begin{tabular}[c]{@{}c@{}}-0.434***\\ (0.118)\end{tabular} & 
\begin{tabular}[c]{@{}c@{}}-0.318**\\ (0.131)\end{tabular} & 
\begin{tabular}[c]{@{}c@{}}-0.074**\\ (0.036)\end{tabular} & 
\begin{tabular}[c]{@{}c@{}}-0.162***\\ (0.037)\end{tabular} & 
\begin{tabular}[c]{@{}c@{}}-0.030\\ (0.038)\end{tabular} \\
$\Delta \%$ & -25.4\% & -35.2\% & -27.2\% & -7.1\% & -14.9\% & - \\
\midrule
Streamer FE & $\checkmark$ & $\checkmark$ & $\checkmark$ & $\checkmark$ & $\checkmark$ & $\checkmark$ \\
Week FE & $\checkmark$ & $\checkmark$ & $\checkmark$ & $\checkmark$ & $\checkmark$ & $\checkmark$ \\
Observations & 80,762 & 92,753 & 92,861 & 80,762 & 92,753 & 92,861 \\
Mean Dependent Variable & 1.297 & 3.356 & 3.692 & 3.615 & 6.577 & 7.382 \\
\bottomrule
\end{tabular}
\end{threeparttable}}
\caption{\small \textbf{The effect of Twitch's banning policy on viewers' log number of views by age group.} 1. Number of views is defined as the average number of presence of viewers from each age group within each streamer-day. The log variables are transformed as $\log(1+y)$. 2. All standard errors are clustered at the streamer level. The mean dependent variable is calculated based on observations of all treated (banned and unbanned) streamers before the policy announcement. FE = fixed effects. Significance levels: *$p<0.1$; **$p<0.05$; ***$p<0.01$.}
\label{ATT_log_num_views}
\end{table}

Two behavioral mechanisms may help explain this weaker response among minors. First, minors may have been less exposed to scandals and social controversies surrounding gambling streamers, and thus less motivated to disengage completely from their livestreams. Second, minors might exhibit stronger user stickiness relative to adults, as adults are more likely to be multi-homing or have a broader interest in other adult-targeted content, thereby reallocating their attention more easily to other streamers or platforms. 

Unbanned streamers also experienced declines in both gambling and overall viewership, except that we find no significant reduction in overall content consumption from uncertain-age viewers. These estimates support our previous finding that unbanned gambling streamers lost viewership even if their content was not directly targeted by the policy, further suggesting a negative spillover effect from the banning policy. Moreover, note that while the relative declines in gambling content were largely comparable between Table \ref{ATT_log_num_views} and Table \ref{ATT_hrs_watched}, the estimated reductions in number of viewers of overall content was notably smaller compared to that measured by total hours watched reported in Table \ref{ATT_hrs_watched}. For example, total hours watched of all content fell by 83.2\% for banned streamers, yet the largest reduction in the number of total viewers was roughly one-half across age groups. This comparison suggests that the consumption reduction was more substantial at the intensive margin rather than at the extensive margin: viewers were less addicted to treated streamers' content in general, and these streamers lost viewing hours from each view.

In Appendix Tables \ref{ATT_log_chats_ps} -- \ref{ATT_viewer_token_count}, we further examine the policy's effect on engagement metrics, including the average number of chats sent in each view, the average number of emotes, characters, and tokens per chat message. On average, minor viewers sent 54.7\% fewer chat messages per viewing of banned streamers' gambling streams (declining from 2.14 to 0.97 messages). This reduction was smaller than that of identified adults ($66.8\%$) and almost identical to that of uncertain-age viewers ($54.3\%$). At the message level, minor viewers sent 0.107 fewer emotes, 8.256 fewer characters, and 1.502 fewer tokens per message in banned streamers' gambling livestreams, all of which represent smaller reductions than those observed among adult viewers (0.134 fewer emotes, 10.575 fewer characters and 1.787 fewer tokens) and among uncertain-age viewers (0.113 fewer emotes, 10.646 fewer characters and 1.958 fewer tokens). Finally, engagement metrics for unbanned streamers’ overall content remained largely unaffected, indicating that the spillover effect primarily affected content consumption rather than detailed interaction behaviors. Additionally, Appendix Figures \ref{es_log_num_views}–\ref{es_avg_token_count} present the time-varying treatment effects for all outcome variables. Across all age groups, the decline in gambling livestreams remained persistent two months after the policy implementation.

To summarize, we find that the banning policy substantially lowered both minors’ consumption of gambling livestreams and their engagement intensity. Considering that minors on average consumed far less gambling content on Twitch, these results suggest that self-regulation policies can be effective in reducing harmful exposure for minors, who often lack cognitive maturity to assess risk and are most susceptible to developing addictive behaviors. However, the policy was relatively less effective for minors compared with non-minor viewers, as their reduction in gambling viewership and engagement was notably smaller. This raises potential concerns, as minors may still be drawn to such content through social influence from friends or classmates who actively engage with gambling-related streams. We further discuss the implications and our recommendations in Section \ref{section-conclusion}. 

% \section{Website Traffic Analysis} \label{section-website}

% \textcolor{red}{This is for AE's comment on adding a brief 
% caveat when discussing null results for the website traffic section:!!!!!!!
% We note that Semrush provides estimated---not observed---traffic based on clickstream panels and search analytics, with potentially limited precision for niche gambling sites. We discuss the implications of this measurement error for our null results in Section~\ref{section-website-traffic}.}

\section{Implications and Conclusions}\label{section-conclusion}

In this paper, we empirically study the impact of Twitch's banning policy on content production and consumption. Based on a novel high-frequency panel dataset of livestreams from top Twitch influencers, we leverage video analysis, text analysis and a threshold-based classification approach to identify banned versus unbanned gambling content within streams, then partition streamers into groups which were directly or indirectly affected by the policy. Moreover, we employ multiple empirical approaches to overcome potential identification challenges.

Our findings yield several insights with managerial and policy relevance for digital platforms, content creators, and policymakers. First, we show that the banning policy led to a persistent decline in the supply of gambling livestreams and generated spillover effects on unbanned gambling content and gambling streams in other languages. Moreover, we find that the policy had more prominent impact on streamers with higher popularity, which can be explained by their concerns for personal reputation. These findings highlight the value of targeted self-regulation on digital platforms: rather than adopting blanket bans, platforms can design policies with narrower scopes but still induce broad behavioral changes by leveraging creator heterogeneity and motivation. Meanwhile, we find evidence of reduced production in non-targeted content categories, suggesting that even well-intentioned self-regulation may have unintended consequences on overall content volume and diversity.

Second, we document substantial policy effects on various viewer demand and engagement metrics, demonstrating the importance of using engagement measures extracted from unstructured data for policy evaluation. In particular, we find that the policy led to significant reductions in overall demand, while metrics reflecting loyal viewers’ activity remained largely unaffected. This heterogeneity implies that creators and platforms should anticipate non-uniform audience responses to regulatory interventions. Creators may need to strengthen engagement with peripheral viewers who are more prone to churn, by providing additional entertainment value. Platforms might consider incentive mechanisms such as rewards for longer viewing durations or gifted subscriptions to sustain engagement alongside self-regulation policies.

Finally, our analysis underscores the importance of protecting minors from potentially harmful content. Although viewers classified as minors are, on average, much less exposed to gambling streams on Twitch, those who do watch such content exhibit engagement levels comparable to non-minors. More crucially, we find that the policy produced smaller relative reductions in minors’ gambling viewership and engagement intensity, suggesting slightly lower effectiveness of the policy in protecting minors. This could reflect narrower interest of minors in alternative adult-targeted content, stronger user stickiness, or lower awareness of public debates over gambling scandals. Since minors are more susceptible to developing addictive behavior and peer influence, our estimates call for complementary actions that aim to protect minors, such as identifying potential minor users and downranking regulated content in their recommendations, to strengthen minor protection.

Our research is subjected to several limitations that point to potential directions for future work. First, it is plausible that streamers significantly affected by the policy may migrate to other platforms to continue streaming gambling content. While our dataset shows very few such cases among the top 6,000 streamers, it omits streamers with very low popularity and thus prevents us from fully discussing how the policy affected market exits among these streamers. Future research could benefit from datasets including an even broader range of streamers or streamer activities on other platforms, and examine how these policies influence competition with rival platforms. Second, we focus on how the policy changed overall supply and demand of streaming content on Twitch, without deeply examining the underlying mechanisms behind substitution and competition among different types of content. Future studies could explore how content creators make decisions when facing policy interventions and investigate the generalizability of our results to other regulation policies in content markets. Third, content regulation may incentivize creators to change their in-stream behaviors, such as encouraging more financial contributions from loyal fans or altering content tone. These potential changes are operated at the intensive rather than extensive margin, which requires deep learning or AI-driven methods to analyze within-stream behavorial shifts. A potential direction for future research would be to explore whether content regulations on digital platforms lead to qualitative changes in content creation, including greater reliance on loyal user interactions or AI-assisted production.


\section*{Funding and Competing Interests}\label{funding-competing-interests}
All author(s) certify that they have no affiliations with or involvement in any organization with any financial or nonfinancial interest in the materials discussed in this manuscript. Author(s) have no competing interests to declare.




\end{spacing}

\newpage

% \bibliographystyle{agsm}
% \bibliography{reference.bib} % if more than one, comma separated

{\singlespacing
\bibliographystyle{agsm}
\bibliography{reference.bib}
}

\newpage
%\begin{appendices}
\appendix
\renewcommand \thepart{}
\renewcommand \partname{}
\addcontentsline{toc}{section}{Online Appendix} % Add the appendix text to the document TOC
\part{Online Appendix} % Start the appendix part
\parttoc % Insert the appendix TOC

%\section*{Appendix}
\renewcommand{\thesection}{\Alph{section}}
\renewcommand{\thesubsection}{\Alph{section}.\arabic{subsection}}
\counterwithin{figure}{subsection}
\counterwithin{table}{subsection}
\counterwithin{algorithm}{subsection}

\renewcommand{\thealgorithm}{\thesubsection.\arabic{algorithm}}

\newpage

\section{Data: Additional Details}\label{data}
\subsection{Streamer Data Collection} \label{data-collection}
In our main data, we utilized a subset of high-frequency streaming activity data collected by \cite{yangwiptwitch}. This data tracks the streaming activities of 30,000 Twitch streamers (and individual registered viewers in streams), by sending requests to the Twitch API every 15 minutes from October 26, 2022 to August 20, 2023. Every 15 minutes, the authors retrieved from the API the stream information (e.g., started time, streamed game, viewer count, registered viewer list) of all tracked streamers. Table \ref{variable-def} describes variables retrieved from the API at the time of each request. 


\begin{table}[!htb]
\centering
\footnotesize
\renewcommand{\arraystretch}{1.2}
\begin{threeparttable}
\begin{tabular}{l l l}
\hline 
Variable & Description & Object \\
\hline
id & ID that identifies the stream & String \\
user\_id & ID of the streamer & String \\
user\_login & The login name of the streamer & String \\
game\_login & ID of the game or category being streamed & String \\
type & The type of stream (e.g., live) & String \\
language & The language that the stream uses & String \\
title & The title of the stream & String \\
started\_at& The UTC date and time of when the stream begins & Datetime object \\
viewer\_count & Number of users watching the stream & Integer \\
registered\_viewer\_list\tablefootnote{Registered viewers in chats were collected from the Twitch TMI API. The data collection on registered viewers stopped on March 31, 2023, after which Twitch permanently shut down the third-party Legacy Chatters endpoint on April 3, 2023.} & The list of registered users in chats & List of Strings \\
registered\_viewer\_count & Number of registered users in chats & Integer \\
thumbnail\_url& The URL to an image of a frame from the last 5 minutes of the stream & String \\
is\_mature& True if the stream is for mature audiences & Boolean \\
tags & The tags applied to the stream & String \\
timestamp & The ET date and time of when the data is requested from the API & Datetime object \\
\hline 
\end{tabular}
\end{threeparttable}
\caption{\small \textbf{Variable definitions}.}
\label{variable-def}
\end{table}

The streamers are pre-selected from a list of 3,606,607 streamers who appear in a 1-month pilot study (from September 21, 2022 to October 18, 2022). During the pilot study, the authors sent requests to the Twitch API, at an hourly rate, to retrieve the top 100,000 most-viewed streamers that were live on Twitch at the time of the request. In total, the pilot study covers activities of 3,606,607 unique streamers broadcasting 34,704 types of content (e.g., both games and non-games such as ``Just Chatting'', ``Pools, Hot Tubs, and Beaches", ``Sports"). From this list, the authors then selected a subset of 30,000 streamers to track using weighted sampling, with weights proportional to the average viewership over the pilot study. 
The sampling approach ensures that the tracked streamers are representative of streamers on the Twitch platform (i.e., covering both superstars and lesser-known streamers). The sample size was selected by taking into account: (1) the daily quota for Twitch API requests, (2) the need to balance data collection efficiency and coverage (since IDs of registered viewers in streams were also tracked). 

\noindent \textbf{Registered viewers in chats vs all viewers}. Since Twitch API does not provide the list of viewers, we use registered viewers in chats to proxy for all viewers in a stream. This enables us to use individuals' watching history to construct valid treated and untreated groups to study the policy's effect on the demand side (See Section \ref{section-demand-results} for details). Table \ref{distribution-comparison} shows the distribution of registered viewers in chats compared to the total viewers. We show that the distribution of registered viewers in chats is generally representative of the total viewership in stream, evidenced by similar distributions and medians (except that viewer count has a higher mean, as registered viewers in chats is a subset of viewers in streams) and high correlation (correlation = 0.988). 



\begin{figure}[!htb]
    \centering
    \begin{subfigure}[b]{0.499\textwidth}
        \centering
        \includegraphics[width=\textwidth]{Figures/Appendix/viewer_distribution.png}
        \label{fig:panelA}
    \end{subfigure}
    \hspace{-0.28cm}
    \begin{subfigure}[b]{0.499\textwidth}
        \centering
        \includegraphics[width=\textwidth]{Figures/Appendix/chatter_distribution.png}
        \label{fig:panelB}
    \end{subfigure}
    \caption{\small \textbf{Distribution of Registered Viewers vs Viewers}}.
    \label{distribution-comparison}
\end{figure}


\subsection{Emotes Data in Chat Logs}\label{appendix-chat-engagement-data}

Besides sending chat messages, Twitch viewers are usually heavily engaged in streams through sending emotes (including emojis). Emotes are small pictures that are used in place of a word or comment to express the emotion of the individual posting the emote. On Twitch, emotes come in two main types: global emotes, which are available to all users on the platform, and channel-specific emotes, which are unique to a particular channel. The channel-specific emotes are created by streamers and can be unlocked by viewers through several methods, such as following a channel (follower emotes), subscribing to a channel (subscriber emotes), cheering with Bits (bits tier emotes), using channel points (channel point emotes), or participating in celebrations or events within a channel (event emotes, such as like Hype Train or the Glitch emote Drops from GlitchCon 2020). Using Twitch API, we compiled a full list of 199,984 emotes covering our sample's streamers and their streams, including 308 global and 199,676 channel-specific emotes. See Figures \ref{global-emotes} and \ref{channel-specific-emotes} for examples.  


\begin{figure}[!htb]
    \centering
    \includegraphics[scale=0.4]{Figures/Appendix/global-emotes.png}
    \caption{\textbf{A screenshot of examples of the Twitch global emotes.} The figure illustrate examples about popular emotes that are accessible by any Twitch users. These emotes can be used in any channel without the need to follow, subscribe, etc to unlock the emotes. Source: \url{https://betterttv.com/emotes/global}.}
    \label{global-emotes}
\end{figure}




\begin{figure}[!htb]
    \centering
    \begin{subfigure}[b]{0.45\textwidth}
    \centering
    \includegraphics[width=\textwidth]{Figures/Appendix/Kaicenat-channel-emotes.png}
    \caption{KaiCenat}
    \label{fig:visit-us-desktop}
    \end{subfigure}
    \begin{subfigure}[b]{0.45\textwidth}
    \centering
    \includegraphics[width=\textwidth]{Figures/Appendix/kamet0-channel-emotes.png}
    \caption{Kamet0}
    \label{fig:visit-us-all}
    \end{subfigure}
    \caption{\textbf{Screenshots of Twitch channel-specific emotes}. The figures show examples of channel-specific emotes for Twitch streamers KaiCenat (left) and Kamet0 (right). These are custom emotes that can be unlocked by following, subscribing, cheering bits, or participating in certain celebrations or events within a channel that have them enabled.}
    \label{channel-specific-emotes}
\end{figure}


\newpage

\section{Banned Content and Streamer Detection: Details on Procedures} \label{appendix-banned-content-detection}

\subsection{Identifying Banned Content From Video Clips and VODs} \label{appendix-treatment-detection-clips}

For each treated streamer, we fetch URLs of all their past gambling video clips pre-policy using Twitch API.\footnote{Clipping (up to 60 seconds) is a functionality enabled by Twitch that allows both the streamer and any logged-in viewer to capture and share moments from a streamer's livestreams. Thus, we treat clips as random stream moments, for which we can use as one source to check any existence of banned content production by streamers. We restrict our analysis to content that was produced between August 1 and October 17, 2022 (inclusive) to align with the coverage of our stream-level data.} We also complement the video clip analysis with Video on Demand (VOD) - archives of video content - posted on other platforms if they are made available by any streamers (e.g., xQc, Adin Ross and ItsSliker).\footnote{VODs are archives of content previously streamed live on Twitch and can span several hours (if unedited). Although VODs on Twitch are only saved up to 2 months for Partners, streamers can enable VOD storage and export their videos to other platforms, such as YouTube, by linking with their Twitch accounts. For example, xQc's VODs can be found on YouTube via \url{https://www.youtube.com/@xqcgames3433/videos}. For long VODs, we manually check through the content to identify the presence (and time of presence) of the banned content.} 

\noindent \textbf{Detection Procedure}. We use a 3-step procedure as shown in Figure \ref{plot-banned-content-from-clips} to identify banned content and streamers. In step 1, we download all past gambling video clips of a streamer (over 4TB) based on clip URLs using \textit{youtube-dl} package (which also supports Twitch video downloads). In step 2, we decompose videos into frames and sample three equally spaced frames (start, middle and end) for content analysis.\footnote{This sampling approach is effective as clips are under 60 seconds with minimal scene changes, thus reducing the number of frames for further content processing without losing critical information.} Finally, in step 3, we convert sampled frames to grayscale\footnote{Grayscale conversion is a common pre-processing step in video analysis. It helps reduce distractions caused by color variance and enhances the focus on text detection.} and extract texts using Google's \textit{Tesseract-OCR}\footnote{\url{https://github.com/tesseract-ocr/tesseract}.}. We check the extracted texts against all banned websites' keywords. We use a restrictive criterion for classification by considering clips as containing banned content only if keywords from banned websites are detected in all sampled frames (or unbanned if absent in all sampled frames), to avoid misclassifications. 


\begin{figure}[h!]
    \centering
    \includegraphics[scale=0.5]{Figures/Identification/illustration-word-detection.png}
    \caption{\small \textbf{Banned content detection using videos}.}
    \label{plot-banned-content-from-clips}
\end{figure}



\noindent \textbf{Banned Stream and Streamer Construction}. To compile the final list of streams and streamers with banned content, we need to match all videos (i.e., video clips and VODs) with the original game streams. We experience a technical challenge here since (i) not all clips have a video offset time from the Twitch API, which tells us how many seconds into a video (and thus stream) the start of the clip occurs, and (ii) a clip's ``created time'' extracted from the Twitch API is not necessarily the time when the actual content is broadcasted but its upload time. This can create an issue if 
there is a severe lag between when a viewer clips an in-stream moment and when she uploads the clip.

For clips that suffer from these issues, our strategy of matching relies on the fact that if a clip features a streamer's desktop time (see Figure \ref{illustration-desktop-time}), the desktop time must reflect the \textit{local time} the content in that clip is streamed. In that case, we can use OCR to extract the timestamp\footnote{The timestamp will be extracted by OCR as part of the texts, in a format that contains the date, hour and minute component of a timestamp} from video frames and then use it to determine a clip's \textit{local} start and end time. There are several caveats: 

\begin{enumerate}
    \item Sometimes multiple timestamps may appear across all sampled frame of a video. For example, this could happen when a frame captures a moment of a streamer with a chat log shown on the side, where chats on Twitch naturally come with timestamps. In that case, it creates difficulty in determining which timestamp is the desktop time. 
    \item Streamers may operate at different time zones, which can alter both the date and the hour component of a timestamp. For example, since Twitch records timestamps in UTC, a clip capturing a stream moment at 08/13/2022 19:45 in UTC (as recorded by Twitch) may appear as 08/13/2022 15:45 in EDT (as reflected by the desktop time shown in a video frame). 
    \item Streamers may have their desktop time displayed in different formats (e.g., DD/MM/YYYY vs MM/DD/YYYY, or 09:30 PM vs 21:30), which creates confusions in matching to their streams' timestamps. 
\end{enumerate}

\begin{figure}[!htb]
    \centering
    \includegraphics[scale=0.5]{Figures/Appendix/illustration-desktop-time.png}
    \caption{\textbf{Desktop time during a broadcast}. This figure illustrates that OCR can detect a streamer's desktop time from video frames (if this information is displayed during the broadcast). In this example, it shows that XQC was broadcasting Stake at 9:30pm on August 11, 2022 at his local time.}
    \label{illustration-desktop-time}
\end{figure}

We use several methods to address all the above issues. First, we observe from OCR's extracted texts that if there are multiple timestamps extracted from a video frame, a streamer's desktop time always appears as the \textit{last} timestamp in the list. This is because most streamers use windows systems, where the desktop time always appears in the lower right corner of the screen (see the zoomed-in region in Figure \ref{illustration-desktop-time}). OCR detects text in reading order. Therefore, regardless of whether the frame is read by OCR from left to right or top to bottom, the desktop time always appears as the last timestamp in the list of extracted timestamps. This helps resolve the first concern. 

Next, we check whether the desktop time shown in a clip matches with the clip's upload time recorded by Twitch. To achieve this, we use a 2-step matching approach. In the first step, we start by matching the two timestamps based on only the \textit{minute} component of a timestamp, given potential concerns of time zone differences which could alter the date and hour components of a timestamp as discussed in the second concern. To facilitate matching, we convert both timestamps into the same 24-hour format (i.e., if AM/PM is present in the timestamp like 09:30PM, we convert it to 21:30). Then, we match the minute component of the two timestamps while allowing for discrepancy of up to 3 minutes between the timestamps. Since the desktop time may either capture the time at the start, middle or end of a clip, this small discrepancy is set to account for the maximum duration of a clip (i.e., 2 minutes) plus a minor time discrepancy due to transition of a minute (i.e., because desktop time does not reflect the second component, 9:30 PM can either be the start of 9:30 PM or close to 9:31PM). Our rationale here is that it is highly unlikely for the minute component of two timestamps to randomly fall within the same 3-minute range. Therefore, if there is a match, we are confident to proceed to the second step of the matching. 

In the second step, we continue by verifying the \textit{date and hour} component of the two timestamps. Twitch records timestamps in UTC. Given that there are 24 time zones, any local time can be up to 12 hours (in rare cases, 14 hours) ahead or behind UTC. We retain only clips where (i) there is a match on the date, (ii) the hour difference of the two timestamps does not exceed 12 hours, and (ii) the clips fall within a streamer's local time zone. We determine a streamer's local time zone by (i) obtaining information about a streamer's country and city of residence from \textit{Streams Charts}, and (ii) identifying the common hour difference of all clips of a streamer (based on clips' desktop time) relative to UTC and then infer a streamer's time zone based on this hour difference. All clips that satisfy the criteria and have clear start and end times are then matched to their original streams. In that case, if clips are found to contain banned content, the corresponding streams they are matched to are also indicated to contain banned content.

We note that, using the above procedure, some clips will be left out either if they do not contain a streamer's desktop time or cannot be matched based on timestamps. As discussed in Section \ref{section-treatment-detection}, our goal is to minimize false positives and avoid misclassifications. If a banned stream is not captured based on one data source, we will use alternative sources for detection. 


\subsection{Identifying Banned Content From Stream Titles} \label{appendix-treatment-detection-title}

As discussed in Section \ref{section-treatment-detection}, our strategy of detecting streams and streamers with banned content from stream titles relies on checking if both (i) a banned website's referral link is found in the stream title, and (ii) the in-stream game featured at that time is also of online \textit{gambling} content. To provide further support for this strategy, in Table \ref{table-detection-using-stream-titles}, we first provide side evidence by summarizing the number of streamers, with or without any history of streaming gambling content, who use banned \textit{referral links} in stream titles under different scenarios. We have several key findings:


First, we note that by definition, ``untreated'' streamers should have no usage of banned website referral links in \textit{gambling} streams, which is our used strategy. However, even for streams that do not contain gambling content, referral links of banned websites are also almost never used in stream titles by ``untreated'' streamers. Among 4,626 ``untreated'' streamers, only 2 streamers are found to include banned websites' referral links in their stream titles.\footnote{We manually checked the video clips of these two streamers (``rigoncs'' and ``xhocii'') on \textit{Stream Charts} for the streams which they were found to use banned websites' referral links and confirmed that they were not streaming any gambling content at that time. Thus, there is no issue with our treatment group classification. We then checked when these referral links were used. We found that they only appear in a few consecutive streams in a short time period. It could be that these two streamers have a short-term contract with banned websites such that they put referral links in their titles while not needed to stream actual gambling.} In contrast, among ``treated'' streamers, a lot more streamers (52 people) have such usages. 

Second, within ``treated'' streamers, almost no streamers include referral links to banned websites in the titles of their \textit{non-gambling} streams. Only 2 streamers (``pijack11'' and ``poderosobagual'') are exceptions, and they are also the two exceptions that have banned websites' referral links in stream titles in post-treatment period which should not happen. We manually checked the usages of commands in streams by these two streamers and argue that these two exceptions do not weaken the validity of our strategy. We find that they consistently used banned websites' referral links in both gambling and non-gambling streams in pre-treatment periods. Additionally, they only used the referral links in streaming titles within \textit{one day} after the policy is implemented, which could be because they did not immediately change their stream titles. 

Table \ref{table-detection-using-stream-titles} provides preliminary evidence that our strategy - aligned with the convention used by the Twitch community - could be reasonable at identifying banned streamers. To provide direct evidence, in Table \ref{table-detection-using-stream-titles-breakdown} Panel A, we validate our approach after classifying the ``treated'' streamers into banned and unbanned groups by showing that unbanned streamers do not include banned website referral links in their \textit{gambling} stream titles. In general, we also do not use such commands in \textit{any} stream titles. 


We note that our detection approach (``!" followed by banned website keywords) is more conservative compared to defining a stream (or streamer) as containing banned content solely based on the presence of banned website keywords in stream titles. In Table \ref{table-detection-using-stream-titles-breakdown} Panel B, we show that our proposed approach better minimizes concerns of false positives (i.e., misclassifying unbanned content as banned), though the less conservative definition also has reasonable good performance. Nevertheless, for readers concerned with how potential misclassification (if any) may impact the results, we conduct a formal sensitivity analysis to quantify the impact in Appendix \ref{appendix_misclassification}. We also emphasize that we use multiple data sources (video clips, titles, and chats which we discuss next) for banned detection, and all banned streamers who are detected from stream titles are verified to be in the banned group by other sources (video clips or chats), thus further validating our detection approach (see Figure \ref{figure:detection_banned_source} for details). 



% In general, the low usage of the referral link commands to banned gambling websites in non-gambling streams by ``untreated'' streamers is expected. Since these referral link commands are conventions used by gambling streamers to provide viewers quick access to gambling websites in streams, ``untreated'' streamers have low incentive to include these commands in their stream titles if they do not align with their content, as doing so could risk reducing engagement. Additionally, many streamers in our dataset are top streamers who treat streaming as a full-time job. Since these commands are a common convention on Twitch, most streamers have no reason to use them inappropriately unless they are actively streaming gambling content.


\begin{table}[!htb]
\centering
\small
\renewcommand{\arraystretch}{1.2}
\resizebox{\textwidth}{!}{%
\begin{threeparttable}
\begin{tabular}{lcc}
\hline 
& \makecell{``Treated" Streamers \\ (Streamers with \\ Gambling Content)} & \makecell{``Untreated" Streamers \\(Streamers without \\ Gambling Content)}  \\
\midrule
No. Streamers & 475 & 4,626  \\ \hline 
\multicolumn{3}{l}{\textit{\textbf{Using banned website referral links}}} \\
\makecell[l]{Banned Referrals in Gambling Stream Titles (\textbf{Our Approach})} & 50 & 0 \\ 
Banned Referrals in Any Stream Titles & 52 & 2 \\
Banned Referrals in Any Stream Titles Post-Treatment Only & 2 & 0 \\ \hline
% \multicolumn{3}{l}{\textit{\textbf{Panel B: Using banned website keywords}}} \\
% Banned Keywords in Gambling Stream Titles & 85 & 0 \\
% Banned Keywords in Any Stream Titles & 90 & 4 \\
% Banned Keywords in Any Stream Titles In Post-Treatment Period & 4 & 1 \\
% \hline 
\end{tabular}
\end{threeparttable}}
\caption{\textbf{Usage of banned referral links in stream titles for treated and untreated streamers}.}
\label{table-detection-using-stream-titles}
\end{table}

\begin{table}[!htb]
\centering
\small
\renewcommand{\arraystretch}{1.2}
\resizebox{\textwidth}{!}{%
\begin{threeparttable}
\begin{tabular}{lccc}
\hline 
& \makecell{Banned \\ Streamers} & \makecell{Unbanned \\ Streamers} & \makecell{Untreated \\ Streamers}  \\
\midrule
No. Streamers & 158 & 317 & 4,626 \\ \hline 
\multicolumn{4}{l}{\textit{\textbf{Panel A: Using banned website referral links}}} \\
Banned Referrals in Gambling Stream Titles (\textbf{Our Approach}) & 50 & 0 & 0 \\ 
Banned Referrals in Any Stream Titles & 52 & 0 & 2 \\
Banned Referrals in Any Stream Titles Post-Treatment Only  & 2 & 0 & 0 \\ \hline
\multicolumn{4}{l}{\textit{\textbf{Panel B: Using banned website keywords}}} \\
Banned Keywords in Gambling Stream Titles & 78 & 7 & 0 \\
Banned Keywords in Any Stream Titles & 83 & 7 & 4 \\
Banned Keywords in Any Stream Titles Post-Treatment Only & 4 & 0 & 1 \\
\hline 
\end{tabular}
\end{threeparttable}}
\caption{\textbf{Usage of banned referral links and banned website keywords in stream titles for banned, unbanned, and untreated streamers}. Treated streamers are separated into banned and unbanned groups, where the groups reflect those after we finish classifying the streamers using all sources.}
\label{table-detection-using-stream-titles-breakdown}
\end{table}



%%%%%%%%%%%%%%%%%%%%%%%%%%%%%%%%%%%%%%%%%%%%%%%%%%%%%%%%%%%%

\subsection{Identifying Banned Content From Chat Logs} \label{appendix-treatment-detection-chats}

\subsubsection{Justifications for Ground Truth Chat Sample Construction} \label{table-gt-sample-justification}

% To construct a ground truth chat sample for \textit{banned} gambling content, we use chats from gambling streams that both (i)
% have video clips with banned content detected on all three randomly sampled frames of a video using OCR, and (ii) have banned websites' referral links featured in stream titles. For unbanned chats sample, we use multiple criteria as follows: (i) the gambling streams must have video clips with \textbf{no} banned website names detected at all three randomly sampled frames using OCR, (ii) the gambling streams must come from streamers with no prior history of streaming content from banned websites, and (iii) the gambling streams must neither have banned referral links or keywords used in stream titles.

As discussed in Section \ref{treatment-detection-chats}, we use multiple criteria to construct ground truth chats samples separately for banned and unbanned gambling streams (see Figure \ref{table-gambling-stream-criteria}). 

\newtcolorbox{promptblockfixed}[1][]{
  enhanced, breakable=false,
  colback=darkred!4!white, colframe=darkred!80!black,
  colbacktitle=darkred!90!black, coltitle=white,
  boxrule=0.5pt, arc=1.5mm,
  left=6pt, right=6pt, top=4pt, bottom=4pt,
  boxsep=2pt, fonttitle=\bfseries, fontupper=\small, title=#1
}

\begin{figure}[h!]
\centering
\begin{minipage}[t]{0.49\textwidth}\vspace{0pt}% <-- forces top alignment
\begin{promptblockfixed}[Banned Gambling Streams]
\begin{enumerate}[label=(\roman*), leftmargin=*]
\item The gambling streams must have video clips with banned website names detected at all 3 randomly sampled frames using OCR.
\item The gambling stream titles must contain referral links to banned websites.
\end{enumerate}
\end{promptblockfixed}
\end{minipage}\hfill
\begin{minipage}[t]{0.49\textwidth}\vspace{0pt}% <-- forces top alignment
\begin{promptblockfixed}[Unbanned Gambling Streams]
\begin{enumerate}[label=(\roman*), leftmargin=*]
\item The gambling streams must have video clips with no banned website names detected at all three randomly sampled frames using OCR.
\item The gambling streams must result from a streamer with no banned content detected in any of his pre-policy gambling video clips using OCR.
\item The gambling stream titles must neither contain a referral link nor a keyword of a banned website.
\end{enumerate}
\end{promptblockfixed}
\end{minipage}
\caption{\textbf{Criteria used to construct ground-truth chat samples for banned and unbanned gambling streams}.}
\label{table-gambling-stream-criteria}
\end{figure}




While we use both sources (OCR and stream titles) to construct the ground truth chats sample for \textit{banned} gambling streams to ensure accuracy, for \textit{unbanned} gambling content, we emphasize that using the first source alone, i.e. OCR on gambling video clips for which criteria (i) and (ii) are based off, is not enough to guarantee that the stream has no banned content. OCR with no illegal website keywords detected in all three randomly sampled frames is only a necessary but insufficient condition for a stream to be unbanned. For example, a clip may still contain banned content, but OCR cannot detect it if a clip's frames do not feature banned website keywords. Thus, we need a second source, i.e., stream titles for which criteria (iii) is based off, to verify that our generated ground truth chat sample for unbanned is valid.

Having established previously that stream titles can be used to classify banned and unbanned content (see Appendix \ref{appendix-alternative-banned-detection-from-titles}), we now check stream titles for all gambling streams whose video clips already satisfy criteria (i) and (ii) for unbanned streams. We find that indeed almost no video clips (1/161584 =  0.0006$\%$ clips) contain referral links to banned websites (i.e., ``!" followed banned website keywords). This also serves as a cross-check that using referral links is a valid approach for banned and unbanned stream classification. To ensure our ground truth sample is of high quality, we enforce an even stricter restriction as in criteria (iii) to allow neither banned websites' referral link nor keywords to appear in stream titles. In general, the multiple criteria we imposed ensure that we have ground truth chats samples of high quality for both banned and unbanned streams, which serve next as our benchmarks for predicting which gambling streams contain banned content based on in-stream chats (if these streams can neither be classified based on video clips or stream titles). 


% \begin{table}[!htb]
% \centering
% \renewcommand{\arraystretch}{1.2}
% \begin{threeparttable}
% \begin{tabular}{lcc}
% \hline 
% & \makecell{``Treated'' Streamers}  \\
% \midrule
% Total Number & 475 \\ 
% Misclassified unbanned Title (OCR) & 2  \\ 
% Misclassified unbanned Title (OCR + Streamer History) & 2  \\ 
% \hline 
% \end{tabular}
% \end{threeparttable}
% \caption{\small \textbf{Detecting banned content from stream titles}. This table shows the detection of banned streamers by various source types for both the treated and untreated groups.}
% \label{tab-gt-justification}
% \end{table}

% \subsection{Inference and Prediction Using Chats} \label{prediction-chats}
% \begin{table}[!htb]
% \centering
% \small
% \renewcommand{\arraystretch}{1.1}
% \caption{Detecting Banned Content From Stream Title}
% \label{tab-gt-justification}
% \begin{threeparttable}
% \begin{subtable}[t]{\textwidth}\centering
%     \caption{Panel A}
%     \begin{tabular}{lcc}
%     \hline
%     \makecell{Measure}  & \makecell{Best Threshold}  & \makecell{In-Sample Results Across Folds}\\
%     \midrule
%     Total Error &  0.30 & 0.001 \\ 
%     Weighted Total Error  &  0.30 & 0.003 \\ 
%     Accuracy & 0.30 & 0.976 \\ 
%     AUC & 0.30 & 0.976 \\ 
%     Precision & 1.05 & 0.991 \\ 
%     Recall & 0.00 & 0.978 \\ 
%     F1-Score & 0.25 & 0.976 \\ 
%     \hline
%     \end{tabular}
% \end{subtable}
% \vskip\baselineskip

% \begin{subtable}[!ht]{\textwidth}\centering
%     \caption{Panel B}
%     \begin{tabular}{lcc}
%     \hline
%     \makecell{Measure}  & \makecell{Averaged Out-Of-Sample Results Across Folds} \\
%     \midrule
%     Total Error  & 0.002 \\ 
%     Weighted Total Error  &  0.003 \\ 
%     Type I Error &  0.003 \\ 
%     Type II Error &  0.024 \\ 
%     Accuracy & 0.871 \\ 
%     AUC & 0.975 \\ 
%     Precision & 0.964 \\ 
%     Recall & 0.974 \\ 
%     F1-Score & 0.853 \\ 
%     \hline
%     \end{tabular}
% \end{subtable}

% \begin{tablenotes}[para,flushleft]
% \footnotesize
% \textit{Notes}: Total Error=$\text{Type I}^2 + \text{Type II}^2$, Total Error Weighted=$2\times\text{Type I}^2 + \text{Type II}^2$, Precision=$\frac{\text{True Positive}}{\text{Predicted Positive}}$, Recall=$\frac{\text{True Positive}}{\text{Actual Positive}}$, F1=$\frac{2 \times \text{Precision} \times \text{Recall} }{\text{Precision+Recall}}$.
% \end{tablenotes}
% \end{threeparttable}
% \end{table}


\subsubsection{Banned Content Detection Using In-Stream Chats}\label{inference-step-chats}

To understand whether it is sufficient to use referral links to banned websites in chats for banned content classification, we now test our detection strategy based on in-stream chats, which is built on the hypothesis that banned gambling streams will contain a higher proportion of referral links to banned websites chats compared to unbanned gambling streams. 

First, we check the summary statistics of banned and unbanned ground truth chats samples (Table \ref{table:chat-gt-summary-stats-main}). We find that banned gambling streams contain 100 times more referral links of banned websites — despite having fewer total chats. In addition, banned gambling streams have 130 times more unique viewers who posted referral links to banned websites in chats. Both statistics suggest that there is a difference in referral links of banned websites in chats between banned and unbanned streams. 

\begin{table}[!htb]
\small
\centering
\renewcommand{\arraystretch}{1.2}
\begin{threeparttable}
\begin{tabular}{lcc}
\hline 
 & Banned Gambling Streams & Unbanned Gambling Streams \\
\hline
Total Chats & 4,031,759 & 4,305,627 \\
Unique Streams & 534 & 821 \\
Unique Streamers & 39 & 99 \\
Unique Viewers & 199,450 & 258,913 \\
Banned Referrals in Chats & 18,872 & 152 \\
Viewers Posting Banned Referrals in Chats & 12,155 & 92 \\
\hline 
\end{tabular}
\end{threeparttable}
\vspace{0.5em} 
\caption{\textbf{Summary statistics of ground truth chats samples for banned and unbanned gambling streams}.}
\label{table:chat-gt-summary-stats-main}
\end{table}


Given the difference, we next consider whether it is adequate to use a simple test statistic based on the number of referral links to banned websites per hour to classify gambling streams.\footnote{Using referral links per hour instead of total number of referral links in streams helps alleviate the concern that streams of longer duration tend to have more referral links.} In Figure \ref{plot-referral-links-distribution}, we first plot the probability density function and cumulative distribution function of referral links to banned websites used in chats per hour separately for banned and unbanned streams. We observe from the left panel that the distribution of referral links to banned websites in unbanned streams has a huge mass at zero, where banned streams have a significantly larger probability of containing more than one banned websites' referral links per hour. The difference in the distributions is further supported by the right panel, where we show that the distribution of referral links to banned websites in banned streams first-order statistically dominates that in unbanned streams. These visualizations support our hypothesis that using referral links of banned websites in chats could be a valid criteria for classifying banned and unbanned gambling streams. 


\begin{figure}[h!]
    \centering
    \begin{subfigure}[b]{0.49\textwidth}
        \centering
        \includegraphics[width=\linewidth]{Figures/Appendix/referrals-density.png}
        \caption{Density}
        \label{fig:referrals-density}
    \end{subfigure}
    \hfill
    \begin{subfigure}[b]{0.49\textwidth}
        \centering
        \includegraphics[width=\linewidth]{Figures/Appendix/referrals-cdf.png}
        \caption{Cumulative Distribution}
        \label{figure:referrals-cdf}
    \end{subfigure}
    \caption{\textbf{Referral links to banned websites per hour in chats for banned and unbanned gambling streams}. Panel (a) shows the density of banned referral links per hour in chats for banned (orange) and unbanned (blue) streams, whereas Panel (b) shows the cumulative distribution of banned referral links per hour. We right-censored the number of banned referral links per hour at one in both panels.}
    \label{plot-referral-links-distribution}
\end{figure}


Next, we evaluate whether there is a difference between the empirical distributions of the two samples formally by conducting a bootstrapped Kolmogorov-Smirnov (KS) test. In our case, the bootstrapped KS statistics is defined as the maximum absolute difference between two empirical distributions: 
\begin{align}
    D_{n,m} =  \sup_x |F_{\text{unbanned},n}(x) - F_{\text{banned},m}(x)|
\end{align}
$F_{\text{unbanned},n}$ and $F_{\text{banned},m}$ are the empirical distribution functions of the banned and unbanned chats samples respectively. We perform the test using the function from Bootstrapped KS2 Tester\footnote{Source: \url{https://github.com/swharden/Bootstrapped-KS2}} by bootstrapping 5000 times with $m=534$ observations for banned sample and $n=821$ for unbanned sample. Figure \ref{plot-ks-stats-inference-step} shows the distribution of bootstrapped KS statistics. Since at any common levels of $\alpha$, 
\begin{align*}
    D_{n,m} = 0.95 > c(\alpha)\sqrt{\frac{n+m}{nm}} = \sqrt{-\ln\left(\frac{\alpha}{2}\right) \times \frac{1}{2}}\sqrt{\frac{n+m}{nm}}
\end{align*}
we always reject the null that the numbers of referral links per hour follow the same distribution in banned and unbanned streams. 

% We use the bootstrapped version of the test to improve reliability of inference, particularly since our two samples are imbalance in size (534 banned versus 821 unbanned streams) and are also not extremely large.

\begin{figure}[h!]
    \centering
    \includegraphics[scale=0.7]{Figures/Identification/bootstrap-ks.png}
    \caption{\textbf{Distribution of bootstrapped Kolmogorov-Smirnov statistics.} The red curve shows the difference between the cumulative probability distributions of the banned and unbanned samples. The purple line shows the bootstrapped KS statistics, and the green line shows the value at which the distributional difference between the two samples is the largest.}
    \label{plot-ks-stats-inference-step}
\end{figure}

Finally, having established the distributional difference, we conduct a preliminary test to check whether it is adequate to classify banned and unbanned streams based on the KS statistic. We use the KS value (0.2) as the threshold value and find that a simple threshold-based classification rule already achieves good out-of-sample predictive performance (i.e., Type I error=0.115, Type II error=0.010, which are important measures in our context since we want to minimize misclassifications). We also experimented with different threshold values and find that all achieve good predictive performance; the only difference is that a higher threshold tends to be conservative in terms of Type I errors but allows for relatively more Type II errors. With this, we conclude that a threshold-based classification is adequate for detecting banned and unbanned gambling streams, where our target is just to find an optimal threshold value to minimize the total errors. 


\begin{table}[!htb]
\centering
\small
\renewcommand{\arraystretch}{1.2}
\begin{threeparttable}
% -------- Panel A --------
\begin{subtable}[t]{\textwidth}
\centering
\caption{Optimal thresholds and averaged performance across folds}
\label{tab-gt-justification}
\begin{tabular}{lccc}
\hline
Measure  & Best Threshold  & \makecell{Averaged In-Sample \\ Results Across Folds} & \makecell{Averaged Out-of-Sample \\ Results Across Folds} \\
\midrule
Total Error                    & 0.30 & 0.049 & 0.051 \\
Total Squared Error           & 0.30 & 0.001 & 0.002 \\
Weighted Total Squared Error  & 0.30 & 0.002 & 0.003 \\
Type I Error                  & 1.05 & 0.007 & 0.007 \\
Type II Error                 & 0.00 & 0.022 & 0.022 \\
Accuracy                      & 0.30 & 0.976 & 0.975 \\
AUC                           & 0.30 & 0.976 & 0.975 \\
Precision                     & 1.05 & 0.991 & 0.987 \\
Recall                        & 0.00 & 0.978 & 0.978 \\
F1-Score                      & 0.25 & 0.976 & 0.968 \\
\hline
\end{tabular}
\end{subtable}

\vspace{0.6em}

% -------- Panel B --------
\begin{subtable}[t]{0.7\textwidth}
\centering
\caption{Out-of-sample performance at threshold 0.3}
\label{out-of-sample-results}
\begin{tabular}{lc}
\hline
Measure  & \makecell{Averaged Out-of-Sample \\ Results Across Folds} \\
\midrule
Total Error                    & 0.051 \\
Total Squared Error           & 0.002 \\
Weighted Total Squared Error  & 0.003 \\
Type I Error                  & 0.024 \\
Type II Error                 & 0.026 \\
Accuracy                      & 0.975 \\
AUC                           & 0.975 \\
Precision                     & 0.964 \\
Recall                        & 0.974 \\
F1-Score                      & 0.968 \\
\hline
\end{tabular}
\end{subtable}
\caption{\textbf{Results for predictive performance}. Panel (a) summarizes the optimal threshold values for different measures and their corresponding
in-sample and out-of-sample performances. Panel (b) summarizes the predictive performance across measures using 0.3 as the optimal threshold. 
The total error is computed as the sum of Type I and Type II errors. The total squared error is the sum of the squared Type I and Type II errors. A weighted version of the total error gives twice the weight to the squared Type I error compared to the squared Type II error. Precision is computed as the ratio of true positives to predicted positives, while recall is the ratio of true positives to actual positives. The F1 score is computed as the weighted average of precision and recall, given by $2 \times \text{Precision} \times \text{Recall}$/$\text{Precision}+\text{Recall}$.}
\label{predictive-performance-results}
\end{threeparttable}
\end{table}


% \begin{figure}[!htb]
%     \centering
%     \begin{subfigure}{0.48\textwidth}
%         \centering
%         \includegraphics[scale=0.5]{Figures/Appendix/cm-out-of-sample.png}
%         \caption{\small Aggregated confusion matrix across folds}
%         \label{fig:confusion-matrix}
%     \end{subfigure}
%     \hfill
%     \begin{subfigure}{0.48\textwidth}
%         \centering
%         \includegraphics[scale=0.5]{Figures/Appendix/roc-out-of-sample.png}
%         \caption{\small Mean ROC curve across folds}
%         \label{figure:roc}
%     \end{subfigure}
%     \caption{\small \textbf{Performance Metrics: Confusion Matrix and ROC Curve}.}
% \end{figure}


\begin{figure}[h!]
    \centering
    \begin{subfigure}[b]{0.49\textwidth}
        \centering
        \includegraphics[width=\linewidth]{Figures/Appendix/cm-out-of-sample.png}
        \caption{Aggregated confusion matrix across folds}
        \label{fig:confusion-matrix}
    \end{subfigure}
    \hfill
    \begin{subfigure}[b]{0.49\textwidth}
        \centering
        \includegraphics[width=\linewidth]{Figures/Appendix/roc-out-of-sample.png}
        \caption{Average ROC curve across folds}
        \label{figure:roc}
    \end{subfigure}
    \caption{\textbf{Out-of-Sample Performances.} Panel (a) aggregates confusion matrices across folds; panel (b) shows the average receiver operating characteristic (ROC) curve across folds.}
    \label{fig:out_of_sample_perf}
\end{figure}


\begin{figure}[!htb]
    \centering
    \includegraphics[scale=0.5]{Figures/Appendix/banned-streamer-detection-three-sources-mod.png}
    \caption{\textbf{Detection of banned streamers by source type}. This figure shows the number of banned streamers detected by each source type alone and together. In total, we are able to detect 158 banned streamers using different sources.}
    %Sources such as video clips are subject to missing data, as not all streamers have gambling clips available for banned content detection. 
    \label{figure:detection_banned_source}
\end{figure}



\clearpage 
\newpage

\section{Robustness Check on Potential Misclassification in Treatment Assignments}\label{appendix_misclassification}

Although we have used multiple data sources and techniques to detect banned content and banned streamers, one may still concern that some streamers are not classified into the correct treated group due to missing data issues or measurement error from the data collection process. To alleviate this concern, we propose a sensitivity analysis to address potential misclassification issues in our causal estimates. Since we have chosen a conservative approach to detect banned streamers, our groups of banned streamers can be viewed as a subset of streamers who actually streamed banned gambling content. Therefore, we focus on potential misclassification in the group of unbanned streamers. The sensitivity parameter introduced in the sensitivity analysis is the fraction of streamers who actually streamed banned content but were classified into the unbanned group, i.e.
\begin{equation*}
    \frac{\#\; \{\text{streamers who streamed banned content} \mid \text{Classified as unbanned streamers}\}}{\#\; \{\text{streamers who streamed unbanned content} \mid \text{Classified as unbanned streamers}\}}
\end{equation*}

Note that streamer classification is based solely on their content during the pre-policy period. Thus, any potential misclassification arises only from exogenous missing data issues rather than from endogenous changes in response to the banning policy.

We use the sensitivity parameter to test our estimation results from  for the group of unbanned streamers. We find that the signs of all the estimated effects are robust as long as the sensitivity parameter is less than 0.999, i.e. less than half of all streamers in the unbanned group are misclassified. This condition is satisfied, because more than half of streamers in the current unbanned group have shown solid evidence of not streaming banned gambling content in collected video clips, stream titles and in-stream chats. We provide a formal proof of identification with misclassification in the DiD framework and the validity of the sensitivity parameter in the following sections.

\subsection{Technical Details of the Sensitivity Analysis}

To show the validity of our sensitivity analysis, we first summarize our problem with a more general econometric framework. Denote the outcome variable as $Y_t(D_1^\ast, D_2^\ast)$, whereas $t \in \{0,1\}$ denotes the pre- and post-treatment period respectively,\footnote{For simplicity, we present all identification results based on a two-period setup. All results can be extended directly into a panel model with multiple pre- and post-treatment periods.} and $D_1^\ast$, $D_2^\ast$ denote the binary treatments (``banned" group and ``unbanned" group) received by each streamer. However, $D_1^\ast$ and $D_2^\ast$ are latent variables which are not observable, and we are only able to see the revealed treatment assignment $D_1, D_2$ from data, where $D_1, D_2$ may not be identical to $D_1^\ast, D_2^\ast$ for all streamers. For example, if a streamer indeed received the first treatment, i.e. she should have been classified into the ``banned" group, but were misclassified into the ``unbanned" group in our dataset because we found no evidence of streaming banned content, then $D_1^\ast = D_2 = 1$ and $D_1 = D_2^\ast = 0$.

Using the above notations, we can rewrite our outcomes variable as follows:
\begin{align}\label{model}
\begin{split}
    Y_t &= \begin{cases}
        Y_t(1,0)D_1^\ast(1-D_2^\ast) + Y_t(0,1)(1-D_1^\ast)D_2^\ast + Y_t(0,0)(1-D_1^\ast)(1-D_2^\ast)\; & \text{, if } t=1\\
        Y_t(0,0) \; & \text{, if } t=0
    \end{cases}
\end{split}
\end{align}
In this framework, $(Y_0,Y_1,D_1,D_2)$ are \textit{observed} from the dataset, whereas $(D_1^\ast, D_2^\ast)$ are latent variables. We do not impose any latent structures on how the true treatment assignments are mapped to observed treatment assignments. Instead, we impose identification restrictions and bounds on conditional probabilities of (mis)classifications in this framework.

%$\epsilon_j = 1\{D_j^\ast \neq D_j\}$ is an unobserved variable indicating false recovery of the $j-$th treatment. $\epsilon_1 = 0$ means that we have correctly recovered the status of the first treatment for a streamer, and misclassified the streamer if $\epsilon_1 = 1$. Similarly, we fail to correctly recover the status of the second treatment for a streamer if $\epsilon_2 = 1$. 

As in all other DiD frameworks, the identification of ATT relies on imposing the parallel trend assumption. However, we require extra parallel trends because of potential misclassification. We state our assumptions need for identication as follows:

\begin{assumption}[Parallel Trends].\label{assumption-PT}
    \begin{enumerate}
        \item $\mathbb{E}[Y_1(0,0) - Y_0(0,0) \mid D_1^\ast = 1, D_2^\ast = 0] = \mathbb{E}[Y_1(0,0) - Y_0(0,0) \mid D_1^\ast = 0, D_2^\ast = 0]$
        \item 
        $\begin{aligned}[t]
            \mathbb{E}[Y_1(0,0)-Y_0(0,0) &\mid D_1^\ast=1, D_2^\ast=0, D_1=1,D_2=0] = \\
            &\mathbb{E}[Y_1(0,0)-Y_0(0,0) \mid D_1^\ast=1, D_2^\ast=0, D_1=0,D_2=1],
        \end{aligned}$\\[1ex]
        $\begin{aligned}[t]
            \mathbb{E}[Y_1(0,0)-Y_0(0,0) &\mid D_1^\ast=0, D_2^\ast=1, D_1=1,D_2=0] = \\
            &\mathbb{E}[Y_1(0,0)-Y_0(0,0) \mid D_1^\ast=0, D_2^\ast=1, D_1=0,D_2=1].
        \end{aligned}$
    \end{enumerate}
\end{assumption}

The first part of Assumption \ref{assumption-PT} is identical to the parallel trend assumption in standard DiD frameworks without misclassification. In addition, the second part of Assumption \ref{assumption-PT} posits that the outcome variables of both the correctly classified streamers and misclassified streamers should follow the same trend within each group  on the same true treatment assignments $D_1^\ast, D_2^\ast$, i.e. the expected outcome is mean independent of the treatment arm within each (true) treated group.

To show the relationship between our DiD estimates and the causal estimands of interest, we first state the following result which decomposes the DiD estimator under our setup:

\begin{proposition}\label{prop-1}
    \begin{align*}
        \begin{split}
            \beta_{1} &= \mathbb{E}[Y_1(1,0)-Y_1(0,0) \mid D_1^\ast=1, D_2^\ast=0,D_1=0,D_2=0]P(D_1^\ast=1,D_2^\ast=0 \mid D_1=1, D_2=0)\\
            &+ \mathbb{E}[Y_1(0,1)-Y_1(0,0) \mid D_1^\ast=0, D_2^\ast=1,D_1=1,D_2=0]P(D_1^\ast=0,D_2^\ast=1 \mid D_1=1, D_2=0)\\
            &-\mathbb{E}[Y_1(1,0)-Y_1(0,0) \mid D_1^\ast=1, D_2^\ast=0,D_1=0,D_2=0]P(D_1^\ast=1,D_2^\ast=0 \mid D_1=0, D_2=0)\\
            &-\mathbb{E}[Y_1(0,1)-Y_1(0,0) \mid D_1^\ast=0, D_2^\ast=1,D_1=0,D_2=0]P(D_1^\ast=0,D_2^\ast=1 \mid D_1=0, D_2=0)
        \end{split}
    \end{align*}
    and $\beta_{3}$ can be decomposed similarly.
\end{proposition}

As we have introduced in Section \ref{section-treatment-detection} of the paper, we have combined historical video clips, in-stream titles and chat logs to detect banned content as well as banned streamers in a conservative way. Therefore, we can rule out several possibilities of misclassification in our dataset. First, we have no misclassified streamers in the untreated group, as our untreated group only contains streamers with no record of streaming any gambling content in the ground-truth data extracted from Twitch API. Second, we have no misclassified streamers in the observed banned group. This is because all streamers in this group had shown sufficient evidence of streaming banned content. We summarize these restrictions as the following condition:

\begin{condition}[Conditional Probabilities of Classifications]\label{condition-1}
    \begin{align*}
    \begin{split}
        &P(D_1^\ast = 1, D_2^\ast = 0 \mid D_1 = 0, D_2 = 0) = 0\\
        &P(D_1^\ast = 0, D_2^\ast = 1 \mid D_1 = 0, D_2 = 0) = 0\\
        &P(D_1^\ast = 1, D_2^\ast = 0 \mid D_1 = 1, D_2 = 0) = 1\\
        0 \leq \;&P(D_1^\ast = 1, D_2^\ast = 0 \mid D_1 = 0, D_2 = 1) \leq 1
    \end{split}
    \end{align*}
\end{condition}


Based on these conditions, we can prove that we have correctly identified ATT for the banned streamers, and our causal estimate of the unbanned streamers can be decomposed into a weighted average of the actual ATT for the two treated groups. We summarize these results in the following proposition:

\begin{proposition}\label{prop-2}
Suppose that Assumption \ref{assumption-PT} and Condition \ref{condition-1} hold. Then,
    \begin{align}
    \begin{split}
    \beta_{1} = &\mathbb{E}[Y_1(1,0)-Y_1(0,0) \mid D_1^\ast=1, D_2^\ast=0,D_1=1,D_2=0]\\ \label{prop-decomposition}
    \beta_{3} = &\mathbb{E}[Y_1(1,0)-Y_1(0,0) \mid D_1^\ast=1, D_2^\ast=0,D_1=0,D_2=1]P(D_1^\ast=1,D_2^\ast=0 \mid D_1=0,D_2=1) +\\
    &\mathbb{E}[Y_1(0,1)-Y_1(0,0) \mid D_1^\ast=0, D_2^\ast=1,D_1=0,D_2=1]P(D_1^\ast=0,D_2^\ast=1 \mid D_1=0,D_2=1)\\
    = & ATT_1 \times P(D_1^\ast=1,D_2^\ast=0 \mid D_1=0, D_2=1) + ATT_2 \times P(D_1^\ast=0,D_2^\ast=1 \mid D_1=0, D_2=1)
    \end{split}
    \end{align}
\end{proposition}

The results stated in Proposition \ref{prop-2} is quite intuitive. On the one hand, since we do not have any misclassified streamer in the observed banned group, this group of treated streamers can be seen as a subpopulation of all streamers who indeed streamed banned content. Therefore, the second part of Assumption \ref{assumption-PT} ensures that the treatment effect is correctly identified \textit{in the banned group}, as in any classic DiD setup. On the other hand, the observed "unbanned" group is a potential mixture of both ``banned" and unbanned streamers, and the corresponding DiD estimate is therefore a weighted average of the two ATTs and is attenuated towards the direction of ATT for the ``banned" group. 

Since our DiD estimate for the banned group is correctly identified and is more negative compared to the DiD estimate of the unbanned group, suggesting that the actual ATT for the unbanned group might be nonnegative. Therefore, we adopt a sensitivity analysis based on the above identification results to address the concerns over potential misclassification. The sensitivity analysis is performed based on the following assumption:
\begin{assumption}[Bounds on Probability of Misclassification]\label{assumption-2}
    \begin{equation*}
        \frac{1}{\Gamma} \leq \frac{P(D_1^\ast=1, D_2^\ast = 0 \mid D_1=0,D_2=1)}{P(D_1^\ast=0, D_2^\ast = 1 \mid D_1=0,D_2=1)} \leq \Gamma
    \end{equation*}
\end{assumption}

We introduce sensitivity parameter $\Gamma$ in Assumption \ref{assumption-2}, which bound the fraction between misclassified streamers in the observed unbanned group compared to the correctly classified streamers. It is natural to adopt this sensitivity parameter since the we can choose the value of the sensitivity parameter based on how confident we are on the classification for each streamer in the observed unbanned group. We derive the bounds of ATT as in the following proposition:
\begin{proposition}\label{prop-3}
    Under Assumption \ref{assumption-PT}, \ref{assumption-2} and assume that Condition \ref{condition-1} holds, then
    \begin{align*}
        (1+\frac{1}{\Gamma})\beta_{3} - \Gamma \beta_{1} \leq ATT_2 \leq (1+\Gamma)\beta_{DiD,2} - \frac{1}{\Gamma} \beta_{DiD,1}\; \text{ if }\beta_{d} > 0\\
        (1+\Gamma)\beta_{3} - \frac{1}{\Gamma} \beta_{1} \leq ATT_2 \leq (1+\frac{1}{\Gamma})\beta_{3} - \Gamma \beta_{1}\; \text{ if }\beta_{d} < 0
    \end{align*}
\end{proposition}

\clearpage
\subsection{Usage of the Sensitivity Parameter}

We can then test the sensitivity of our causal estimates in Section \ref{section-supply-results} based on the derived bounds. For example, the estimated treatment effects on the log weekly streaming hours of gambling contents are -1.028 and -0.125 respectively for the two treated groups. Plugging these estimates into Proposition \ref{prop-3}, We get
\begin{align*}
    0.125(2.028\frac{1}{\Gamma}-1-\Gamma) < ATT_2 < 0.125(2.028\Gamma-1-\frac{1}{\Gamma})
\end{align*}
Based on the above result, we are confident with the result that the banning policy led to a reduction in the supply of gambling content of unbanned streamers as long as
\begin{equation*}
    0.125(2.028\Gamma-1-\frac{1}{\Gamma}) < 0 \Leftrightarrow \Gamma \in (0,0.9908)
\end{equation*}
i.e. the fraction of misclassified banned streamers in our observed unbanned group is less than approximately 50\%. Similarly, the range of sensitivity parameters for the estimates in Table \ref{ATT-streaming-hrs} are $(0, 2.427)$ and $(0, 1.465)$ respectively. Therefore, both of them are valid under the threat of misclassification as long as the fraction of misclassified banned streamers in our observed unbanned group is less than approximately 50\%.

\clearpage

\subsection{Proofs for Appendix \ref{appendix_misclassification}}

\subsubsection{Proof of Proposition \ref{prop-1}}

Note that
\begin{align*}
    \beta_{DiD,1} &= \mathbb{E}[Y_1 - Y_0 \mid D_1=1, D_2=0] - \mathbb{E}[Y_1 - Y_0 \mid D_1=0, D_2=0]\\
    &= (\mathbb{E}[Y_1 \mid D_1=1,D_2=0] - \mathbb{E}[Y_1 \mid D_1=0,D_2=0]) \\ &-  (\mathbb{E}[Y_0 \mid D_1=1,D_2=0] - \mathbb{E}[Y_0 \mid D_1=0,D_2=0])
\end{align*}

Plugging the expression of $Y_t$ from (\ref{model}), we have
\begin{align}
\begin{split}
    &\mathbb{E}[Y_1 \mid D_1=1,D_2=0] - \mathbb{E}[Y_1 \mid D_1=0,D_2=0]\\
    = \;&\mathbb{E}[Y_1(1,0)D_1^\ast(1-D_2^\ast) + Y_1(0,1)(1-D_1^\ast)D_2^\ast + Y_1(0,0)(1-D_1^\ast)(1-D_2^\ast) \mid D_1=1,D_2=0]\\
    &- \;\mathbb{E}[Y_1(1,0)D_1^\ast(1-D_2^\ast) + Y_1(0,1)(1-D_1^\ast)D_2^\ast + Y_1(0,0)(1-D_1^\ast)(1-D_2^\ast) \mid D_1=0,D_2=0]\\
    = \;&\mathbb{E}[(Y_1(1,0)-Y_1(0,0)D_1)D_1^\ast(1-D_2^\ast) + (Y_1(0,1)-Y_1(0,0))(1-D_1^\ast)D_2^\ast \\ &+Y_1(0,0)(1-D_1^\ast D_2^\ast) \mid D_1=1,D_2=0]\\
    &-\mathbb{E}[(Y_1(1,0)-Y_1(0,0)D_1)D_1^\ast(1-D_2^\ast) + (Y_1(0,1)-Y_1(0,0))(1-D_1^\ast)D_2^\ast \\ &+ Y_1(0,0)(1-D_1^\ast D_2^\ast) \mid D_1=0,D_2=0]\\
    = \;&\mathbb{E}[Y_1(1,0)-Y_1(0,0) \mid D_1^\ast=1, D_2^\ast=0,D_1=0,D_2=0]P(D_1^\ast=1,D_2^\ast=0 \mid D_1=1, D_2=0)\\
    &+ \;\mathbb{E}[Y_1(0,1)-Y_1(0,0) \mid D_1^\ast=0, D_2^\ast=1,D_1=0,D_2=0]P(D_1^\ast=0,D_2^\ast=1 \mid D_1=1, D_2=0)\\
    &+ \;\mathbb{E}[Y_1(0,0) \mid D_1=0,D_2=0)] - \mathbb{E}[Y_1(0,0)D_1^\ast D_2^\ast \mid D_1=1,D_2=0]\\\
    &- \;\{\mathbb{E}[Y_1(1,0)-Y_1(0,0) \mid D_1^\ast=1, D_2^\ast=0,D_1=1,D_2=0]P(D_1^\ast=1,D_2^\ast=0 \mid D_1=0, D_2=0)\\
    &+ \;\mathbb{E}[Y_1(0,1)-Y_1(0,0) \mid D_1^\ast=0, D_2^\ast=1,D_1=0,D_2=0]P(D_1^\ast=0,D_2^\ast=1 \mid D_1=0, D_2=0)\\
    & + \;\mathbb{E}[Y_1(0,0) \mid D_1=0,D_2=0)] - \mathbb{E}[Y_1(0,0)D_1^\ast D_2^\ast \mid D_1=1,D_2=0]\}\\
    = \;&\mathbb{E}[Y_1(1,0)-Y_1(0,0) \mid D_1^\ast=1, D_2^\ast=0,D_1=0,D_2=0]P(D_1^\ast=1,D_2^\ast=0 \mid D_1=1, D_2=0)\\
    &+ \;\mathbb{E}[Y_1(0,1)-Y_1(0,0) \mid D_1^\ast=0, D_2^\ast=1,D_1=1,D_2=0]P(D_1^\ast=0,D_2^\ast=1 \mid D_1=1, D_2=0)\\
    &- \;\mathbb{E}[Y_1(1,0)-Y_1(0,0) \mid D_1^\ast=1, D_2^\ast=0,D_1=0,D_2=0]P(D_1^\ast=1,D_2^\ast=0 \mid D_1=0, D_2=0)\\
    &- \;\mathbb{E}[Y_1(0,1)-Y_1(0,0) \mid D_1^\ast=0, D_2^\ast=1,D_1=0,D_2=0]P(D_1^\ast=0,D_2^\ast=1 \mid D_1=0, D_2=0)
\end{split}
\end{align}
whereas the last equality holds under Assumption \ref{assumption-PT} and the fact that
\begin{equation*}
    \mathbb{E}[Y_1(0,0)D_1^\ast D_2^\ast \mid D_1=1,D_2=0] \equiv 0
\end{equation*}
under Assumption \ref{assumption-2}. The decomposition of $\beta_{DiD,2}$ can be obtained similarly. $\square$

\subsubsection{Proof of Proposition \ref{prop-2}}

The results are derived directly by plugging the bounds of conditional probabilities into results of Proposition \ref{prop-2}. $\square$

\subsubsection{Proof of Proposition \ref{prop-3}}

Note that with model (\ref{model}), there will be no misclassification in the untreated group, i.e.
\begin{equation*}
    P(D_1^\ast=1, D_2^\ast = 0 \mid D_1=0,D_2=1)+P(D_1^\ast=0, D_2^\ast = 1 \mid D_1=0,D_2=1)=1
\end{equation*}
Combining this equality with Assumption \ref{assumption-2} yields the following bounds on the probability of correct classification conditional on actual second treatment:
\begin{align}\label{proof-prop-2-bound}
    \frac{1}{1+\Gamma} \leq P(D_1^\ast=0,D_2^\ast=1 \mid D_1=0,D_2=1) \leq \frac{1}{1+1/\Gamma}
\end{align}
Therefore,
\begin{align*}
    ATT_2 &= \frac{1}{P(D_1^\ast=0,D_2^\ast=1 \mid D_1=0,D_2=1)}\beta_{DiD,2}-ATT_1 \times\frac{P(D_1^\ast=1,D_2^\ast=0 \mid D_1=0,D_2=1)}{P(D_1^\ast=0,D_2^\ast=1 \mid D_1=0,D_2=1)}\\
    &= \frac{1}{P(D_1^\ast=0,D_2^\ast=1 \mid D_1=0,D_2=1)}\beta_{DiD,2}-\beta_{DiD,1} \times\frac{P(D_1^\ast=1,D_2^\ast=0 \mid D_1=0,D_2=1)}{P(D_1^\ast=0,D_2^\ast=1 \mid D_1=0,D_2=1)}\\
    &\leq \frac{1}{P(D_1^\ast=0,D_2^\ast=1 \mid D_1=0,D_2=1)}\beta_{DiD,2} - 1/\Gamma\beta_{DiD_1}\\
    &\leq (1+\Gamma)\beta_{DiD,2} - 1/\Gamma\beta_{DiD,1}
\end{align*}
when $\beta_{DiD,d} >0$ for $d=1,2$. The lower bound of $ATT_2$ when $\beta_{DiD,d} > 0$ and the bounds when $\beta_{DiD,d} < 0$ can be obtained similarly. $\square$


\clearpage
\newpage

\section{Additional Robustness Checks}\label{appendix-robustness}

% \subsection{Event Studies}\label{appendix-event-study}

% Given that the banning policy reduced the supply of gambling streams and caused spillover effects on unbanned gambling streams, we now investigate the persistence of the impact. The effect of the policy might diminish if treated streamers temporarily ceased gambling livestreams immediately after the policy's implementation, but resumed them later when demand increased, such as during holidays. We use an event-study to analyze potential changes in the treatment effects over time, and to provide visual evidence on the parallel trend assumption. The event study adopts the following specification:
% \begin{equation}
%     y_{it} = \alpha_i + \gamma_t + \sum_{t = -7}^{14}\tilde{\beta}_{1t} \text{Banned}_i + \sum_{t = -7}^{14}\tilde{\beta}_{2t} \text{Unbanned}_i + \varepsilon_{it} \label{equation-es-specification}
% \end{equation}

% The time indices in equation (\ref{equation-es-specification}) are determined as we have 7 weeks before the policy announcement and 14 weeks after. We normalize the point estimate one week before the policy announcement to zero, and report the point-wise confidence interval of each estimate based on standard errors clustered at the streamer level.

% Figure \ref{es_supply_ban} and Figure \ref{es_supply_unban} show the estimates of time-varying treatment effects on the four types of content creation among treated streamers discussed in Section \ref{section-supply-results}. First, we observe a small negative pre-trend in the supply of gambling content, but the magnitude is significantly smaller compared to the post-treatment estimates. There was also a small pre-trend in the supply of gambling content among unbanned streamers, but there was a clear upward trend in the estimates which was opposite to the post-period effects. Therefore, these event study plots do not invalidate our main findings. Consistent with our estimates reported in Table \ref{ATT-streaming-hrs}, we find that the policy led to a persistent reduction in gambling livestreams after the policy implementation. Interestingly, we observe that the policy's negative impact increased in two weeks after the policy implementation and remained persistent afterwards. This result provides an additional evidence that the reduction in the supply of gambling content among banned streamers were not completely from the banned websites (because they had been removed from the platform from the first week after the policy implementation). Streamers complete their adjustment of gambling content creation after the second week, as the estimates are not statistically different thereafter. Second, panel (c) in Figure \ref{es_supply_ban} shows that the reduction in total livestreaming hours were persistent for banned streamers, while this reduction diminished for unbanned streamers in the last couple of periods in our panel. Finally, we find that the reductions in the supply of other games for both treated groups were estimated due to the disappearance of a positive pre-trend. We provide additional causal inference methods to test the robustness of these results.

% \begin{figure}[!htb]
%      \centering
%      \begin{subfigure}[b]{0.45\textwidth}
%          \centering
%          \includegraphics[width=\textwidth]{Figures/event/Event_ban_log_slots_hrs.png}
%          \caption{Gambling Content}
%          \label{fig:token1}
%      \end{subfigure}
%      \begin{subfigure}[b]{0.45\textwidth}
%          \centering
%          \includegraphics[width=\textwidth]{Figures/event/Event_ban_log_lootboxes_hrs.png}
%          \caption{Loot Box Games}
%          \label{fig:token2}
%      \end{subfigure}
%      \vspace{-0.3cm}
%      \begin{subfigure}[b]{0.45\textwidth}
%          \centering
%          \includegraphics[width=\textwidth]{Figures/event/Event_ban_log_broadcast_hrs.png}
%          \caption{Total Livestreams}
%          \label{fig:token3}
%      \end{subfigure}
%      \begin{subfigure}[b]{0.45\textwidth}
%          \centering
%          \includegraphics[width=\textwidth]{Figures/event/Event_ban_log_others_hrs.png}
%          \caption{Other Games}
%          \label{fig:token4}
%      \end{subfigure}
%     \vspace{0.5cm}
%     \caption{\small \textbf{Time-varying treatment effects on content creation in group of banned streamers}. 1. Each subfigure shows the point estimates of $\tilde{\beta}_{1t}$. 2. $\tilde{\beta}_{1t}$ in one week ahead of the policy announcement is normalized to zero. 3. The 95\% point-wise confidence intervals are illustrated as the inner bars, and the standard errors are clustered at the streamer level.}
%      \label{es_supply_ban}
% \end{figure}

% \begin{figure}[!htb]
%      \centering
%      \begin{subfigure}[b]{0.45\textwidth}
%          \centering
%          \includegraphics[width=\textwidth]{Figures/event/Event_unban_log_slots_hrs.png}
%          \caption{Gambling Content}
%          \label{fig:token1}
%      \end{subfigure}
%      \begin{subfigure}[b]{0.45\textwidth}
%          \centering
%          \includegraphics[width=\textwidth]{Figures/event/Event_unban_log_lootboxes_hrs.png}
%          \caption{Loot Box Games}
%          \label{fig:token2}
%      \end{subfigure}
%      \vspace{-0.3cm}
%      \begin{subfigure}[b]{0.45\textwidth}
%          \centering
%          \includegraphics[width=\textwidth]{Figures/event/Event_unban_log_broadcast_hrs.png}
%          \caption{Total Livestreams}
%          \label{fig:token3}
%      \end{subfigure}
%      \begin{subfigure}[b]{0.45\textwidth}
%          \centering
%          \includegraphics[width=\textwidth]{Figures/event/Event_unban_log_others_hrs.png}
%          \caption{Other Games}
%          \label{fig:token4}
%      \end{subfigure}
%      \vspace{0.5cm}
%      \caption{\small \textbf{Time-varying treatment effects on content creation in group of unbanned streamers}. Each subfigure shows the point estimates of $\tilde{\beta}_{2t}$. 2. $\tilde{\beta}_{2t}$ in one week ahead of the policy announcement is normalized to zero. 3. The 95\% point-wise confidence intervals are illustrated as the inner bars, and the standard errors are clustered at the streamer level.}
%      \label{es_supply_unban}
% \end{figure}

% \begin{figure}[!htb]
%      \centering
%      \begin{subfigure}[b]{0.45\textwidth}
%          \centering
%          \includegraphics[width=\textwidth]{Figures/event/Event_ban_log_hrs_watched.png}
%          \caption{Total Hours Watched}
%          \label{fig:token1}
%      \end{subfigure}
%      \begin{subfigure}[b]{0.45\textwidth}
%          \centering
%          \includegraphics[width=\textwidth]{Figures/event/Event_ban_log_Tier1.png}
%          \caption{Tier 1 Subscription}
%          \label{fig:token2}
%      \end{subfigure}
%      \vspace{-0.3cm}
%      \begin{subfigure}[b]{0.45\textwidth}
%          \centering
%          \includegraphics[width=\textwidth]{Figures/event/Event_ban_log_Tier2.png}
%          \caption{Tier 2 Subscription}
%          \label{fig:token3}
%      \end{subfigure}
%      \begin{subfigure}[b]{0.45\textwidth}
%          \centering
%          \includegraphics[width=\textwidth]{Figures/event/Event_ban_log_Tier3.png}
%          \caption{Tier 3 Subscription}
%          \label{fig:token4}
%      \end{subfigure}
%     \vspace{0.5cm}
%     \caption{\small \textbf{Time-varying treatment effects on content consumption in group of banned streamers}. 1. Each subfigure shows the point estimates of $\tilde{\beta}_{1t}$. 2. $\tilde{\beta}_{1t}$ in one period ahead of the policy announcement is normalized to zero. 3. The 95\% point-wise confidence intervals are illustrated as the inner bars, and the standard errors are clustered at the streamer level.}
%      \label{es_demand_ban}
% \end{figure}

% \begin{figure}[!htb]
%      \centering
%      \begin{subfigure}[b]{0.45\textwidth}
%          \centering
%          \includegraphics[width=\textwidth]{Figures/event/Event_unban_log_hrs_watched.png}
%          \caption{Total Hours Watched}
%          \label{fig:token1}
%      \end{subfigure}
%      \begin{subfigure}[b]{0.45\textwidth}
%          \centering
%          \includegraphics[width=\textwidth]{Figures/event/Event_unban_log_Tier1.png}
%          \caption{Tier 1 Subscription}
%          \label{fig:token2}
%      \end{subfigure}
%      \vspace{-0.3cm}
%      \begin{subfigure}[b]{0.45\textwidth}
%          \centering
%          \includegraphics[width=\textwidth]{Figures/event/Event_unban_log_Tier2.png}
%          \caption{Tier 2 Subscription}
%          \label{fig:token3}
%      \end{subfigure}
%      \begin{subfigure}[b]{0.45\textwidth}
%          \centering
%          \includegraphics[width=\textwidth]{Figures/event/Event_unban_log_Tier3.png}
%          \caption{Tier 3 Subscription}
%          \label{fig:token4}
%      \end{subfigure}
%      \vspace{0.5cm}
%      \caption{\small \textbf{Time-varying treatment effects on content consumption in group of unbanned streamers}. Each subfigure shows the point estimates of $\tilde{\beta}_{2t}$. 2. $\tilde{\beta}_{2t}$ in one period ahead of the policy announcement is normalized to zero. 3. The 95\% point-wise confidence intervals are illustrated as the inner bars, and the standard errors are clustered at the streamer level.}
%      \label{es_demand_unban}
% \end{figure}

\subsection{DiD with Matched Sample}\label{appendix-match}

To test the robustness of our empirical results, we perform several additional checks. First, our treated and control groups differ significantly in size, and concerns may arise about selection issues due to the nonrandom choice of streaming gambling content. To mitigate these concerns, we use a matched sample to reconduct our main analysis. In marketing research, the matched samples are often generated through Propensity Score Matching \citep{angrist2009mostly}, which allows researchers to match treated and control units based on the estimated probability of receiving the treatment conditional on a given set of covariates. We begin by applying nearest-neighbor matching with a logit model. However, one may also argue that if a streamer had never streamed any gambling content before the policy announcement, her probability of being treated was always zero. To handle this, we also apply Coarsened Exact Matching (CEM), which matches units into stratas based on a set of coarsened covariates. In both methods, we match units using pre-treatment observations of weekly hours watched by viewers, streaming hours of games with and without loot boxes, and of non-gaming content to ensure that matched streamers have both similar streaming schedules and popularity. Table \ref{table-summary-matching} shows the balance check after applying each matching algorithm, and we find that each covariate is well balanced in terms of both mean and distribution.

\begin{table}[!htb]
\renewcommand{\arraystretch}{1.2}
\centering
\resizebox{\textwidth}{!}{
\begin{threeparttable}
\begin{tabular}{lcccccccc}
\hline 
                    & \multicolumn{4}{c}{PSM}                                                                                                       & \multicolumn{4}{c}{CEM}                                                                                                       \\ \cline{2-9} 
                    & \multicolumn{2}{c}{Banned}                                    & \multicolumn{2}{c}{Unbanned}                                  & \multicolumn{2}{c}{Banned}                                    & \multicolumn{2}{c}{Unbanned}                                  \\ \hline
                    & \multicolumn{1}{l}{Mean Diff} & \multicolumn{1}{l}{Var Ratio} & \multicolumn{1}{l}{Mean Diff} & \multicolumn{1}{l}{Var Ratio} & \multicolumn{1}{l}{Mean Diff} & \multicolumn{1}{l}{Var Ratio} & \multicolumn{1}{l}{Mean Diff} & \multicolumn{1}{l}{Var Ratio} \\ \cline{2-9} 
Hours Watched     & 0.093                        & 1.000                        & 0.017                        & 0.981                        & 0.006                        & 1.005                        & 0.003                        & 1.000                        \\
Loot Box Hours    & 0.141                        & 1.166                        & 0.089                        & 1.025                        & -0.001                       & 1.005                        & -0.000                       & 1.001                        \\
Other Games Hours & -0.051                       & 0.845                        & -0.061                       & 0.841                        & -0.005                       & 0.998                        & -0.002                       & 0.998                        \\
Nongaming Hours   & -0.015                       & 0.937                        & -0.091                       & 0.876                        & -0.005                       & 0.998                        & 0.000                        & 1.000                        \\ \hline 
\end{tabular}
\end{threeparttable}}
\caption{\textbf{Summary statistics of matched samples}. The table reports standardized mean differences and variance ratios of covariates between treated and control streamers based on each matching method. The mean differences are calculated by subtracting the mean of the matched control streamers from that of the matched treated streamers. All covariates are log-transformed as $\log(1+y)$.}
\label{table-summary-matching}
\end{table}

Table \ref{table:full-results-matched} reports the TWFE-DiD estimates based on matched units for each algorithm, using the same outcome variables as our main analyses. All findings are qualitatively unchanged, except that we now observe a small negative effect (less than 3\%) on the most expensive subscription choice for banned streamers after the policy. However, given that this impact is significantly smaller than the effect on the cheapest subscription, these results still support our finding that the policy had minimal or no impact on the more loyal viewers of the treated streamers.
\begin{table}[!htb]
\renewcommand{\arraystretch}{1.2}
\centering
\begin{subtable}{\textwidth}
\caption{PSM Sample}
\label{subtable:psm_results}
\centering
\resizebox{\textwidth}{!}{%
\begin{tabular}{lccccccc}
\hline 
            & \multicolumn{3}{c}{Supply} & \multicolumn{4}{c}{Demand} \\ \cline{2-8}
            & Gambling & Loot Box & Streaming & Hrs Watched & Tier 1 & Tier 2 & Tier 3 \\ \hline
Banned $\times$ Post ($\beta_1$)    
            & \begin{tabular}[c]{@{}c@{}}-1.050***\\ (0.088)\end{tabular} 
            & \begin{tabular}[c]{@{}c@{}}-0.001\\ (0.068)\end{tabular}  
            & \begin{tabular}[c]{@{}c@{}}-0.606***\\ (0.085)\end{tabular} 
            & \begin{tabular}[c]{@{}c@{}}-1.800***\\ (0.271)\end{tabular} 
            & \begin{tabular}[c]{@{}c@{}}-0.623***\\ (0.126)\end{tabular}  
            & \begin{tabular}[c]{@{}c@{}}0.008\\ (0.015)\end{tabular} 
            & \begin{tabular}[c]{@{}c@{}}-0.026\\ (0.017)\end{tabular} \\
Banned $\times$ Between ($\beta_2$)    
            & \begin{tabular}[c]{@{}c@{}}0.166**\\ (0.074)\end{tabular}   
            & \begin{tabular}[c]{@{}c@{}}-0.051\\ (0.056)\end{tabular} 
            & \begin{tabular}[c]{@{}c@{}}-0.066\\ (0.059)\end{tabular}    
            & \begin{tabular}[c]{@{}c@{}}-0.333*\\ (0.181)\end{tabular}   
            & \begin{tabular}[c]{@{}c@{}}0.023\\ (0.092)\end{tabular}     
            & \begin{tabular}[c]{@{}c@{}}-0.008\\ (0.015)\end{tabular} 
            & \begin{tabular}[c]{@{}c@{}}-0.023*\\ (0.013)\end{tabular}  \\
Unbanned $\times$ Post ($\beta_3$)    
            & \begin{tabular}[c]{@{}c@{}}-0.109**\\ (0.049)\end{tabular}  
            & \begin{tabular}[c]{@{}c@{}}-0.048\\ (0.045)\end{tabular} 
            & \begin{tabular}[c]{@{}c@{}}-0.173***\\ (0.046)\end{tabular} 
            & \begin{tabular}[c]{@{}c@{}}-0.477***\\ (0.141)\end{tabular} 
            & \begin{tabular}[c]{@{}c@{}}-0.156***\\ (0.060)\end{tabular} 
            & \begin{tabular}[c]{@{}c@{}}0.015\\ (0.011)\end{tabular}   
            & \begin{tabular}[c]{@{}c@{}}0.015\\ (0.012)\end{tabular}    \\
Unbanned $\times$ Between ($\beta_4$)     
            & \begin{tabular}[c]{@{}c@{}}0.133***\\ (0.039)\end{tabular}  
            & \begin{tabular}[c]{@{}c@{}}-0.031\\ (0.041)\end{tabular} 
            & \begin{tabular}[c]{@{}c@{}}-0.156***\\ (0.03)\end{tabular}  
            & \begin{tabular}[c]{@{}c@{}}-0.142\\ (0.123)\end{tabular}    
            & \begin{tabular}[c]{@{}c@{}}-0.151***\\ (0.056)\end{tabular} 
            & \begin{tabular}[c]{@{}c@{}}0.03***\\ (0.01)\end{tabular}  
            & \begin{tabular}[c]{@{}c@{}}0.01\\ (0.009)\end{tabular}     \\ \hline
Streamer FE & $\checkmark$ & $\checkmark$ & $\checkmark$ & $\checkmark$ & $\checkmark$ & $\checkmark$ & $\checkmark$ \\
Week FE     & $\checkmark$ & $\checkmark$ & $\checkmark$ & $\checkmark$ & $\checkmark$ & $\checkmark$ & $\checkmark$ \\
Observations & 102,872 & 102,872 & 102,872 & 102,872 & 102,872 & 102,872 & 102,872 \\
\hline 
\end{tabular}
} 
\end{subtable}

\vspace{0.5cm}

\begin{subtable}{\textwidth}
\caption{CEM Sample}
\label{subtable:cem_results}
\centering
\resizebox{\textwidth}{!}{%
\begin{tabular}{lccccccc}
\hline 
            & \multicolumn{3}{c}{Supply} & \multicolumn{4}{c}{Demand} \\ \cline{2-8}
            & Gambling & Loot Box & Streaming & Hrs Watched & Tier 1 & Tier 2 & Tier 3 \\ \hline
Banned $\times$ Post ($\beta_1$)    
            & \begin{tabular}[c]{@{}c@{}}-1.049***\\ (0.088)\end{tabular} 
            & \begin{tabular}[c]{@{}c@{}}0.013\\ (0.066)\end{tabular}  
            & \begin{tabular}[c]{@{}c@{}}-0.605***\\ (0.083)\end{tabular} 
            & \begin{tabular}[c]{@{}c@{}}-1.781***\\ (0.264)\end{tabular} 
            & \begin{tabular}[c]{@{}c@{}}-0.609***\\ (0.122)\end{tabular} 
            & \begin{tabular}[c]{@{}c@{}}0.006\\ (0.014)\end{tabular}  
            & \begin{tabular}[c]{@{}c@{}}-0.028*\\ (0.016)\end{tabular} \\
Banned $\times$ Between ($\beta_2$)    
            & \begin{tabular}[c]{@{}c@{}}0.166**\\ (0.074)\end{tabular}   
            & \begin{tabular}[c]{@{}c@{}}-0.007\\ (0.054)\end{tabular} 
            & \begin{tabular}[c]{@{}c@{}}-0.129***\\ (0.037)\end{tabular} 
            & \begin{tabular}[c]{@{}c@{}}-0.273\\ (0.173)\end{tabular}    
            & \begin{tabular}[c]{@{}c@{}}0.042\\ (0.087)\end{tabular}     
            & \begin{tabular}[c]{@{}c@{}}0.004\\ (0.013)\end{tabular}  
            & \begin{tabular}[c]{@{}c@{}}-0.018\\ (0.012)\end{tabular} \\
Unbanned $\times$ Post ($\beta_3$)    
            & \begin{tabular}[c]{@{}c@{}}-0.108**\\ (0.049)\end{tabular}  
            & \begin{tabular}[c]{@{}c@{}}-0.046\\ (0.043)\end{tabular} 
            & \begin{tabular}[c]{@{}c@{}}-0.172***\\ (0.045)\end{tabular} 
            & \begin{tabular}[c]{@{}c@{}}-0.487***\\ (0.136)\end{tabular} 
            & \begin{tabular}[c]{@{}c@{}}-0.174***\\ (0.058)\end{tabular} 
            & \begin{tabular}[c]{@{}c@{}}0.012\\ (0.011)\end{tabular}  
            & \begin{tabular}[c]{@{}c@{}}0.013\\ (0.012)\end{tabular}  \\
Unbanned $\times$ Between ($\beta_4$)     
            & \begin{tabular}[c]{@{}c@{}}0.133***\\ (0.039)\end{tabular}  
            & \begin{tabular}[c]{@{}c@{}}-0.046\\ (0.043)\end{tabular}  
            & \begin{tabular}[c]{@{}c@{}}-0.060\\ (0.039)\end{tabular} 
            & \begin{tabular}[c]{@{}c@{}}-0.195\\ (0.119)\end{tabular}   
            & \begin{tabular}[c]{@{}c@{}}-0.178***\\ (0.053)\end{tabular} 
            & \begin{tabular}[c]{@{}c@{}}0.023**\\ (0.009)\end{tabular} 
            & \begin{tabular}[c]{@{}c@{}}0.003\\ (0.008)\end{tabular}  \\ \hline
Streamer FE & $\checkmark$ & $\checkmark$ & $\checkmark$ & $\checkmark$ & $\checkmark$ & $\checkmark$ & $\checkmark$ \\
Week FE     & $\checkmark$ & $\checkmark$ & $\checkmark$ & $\checkmark$ & $\checkmark$ & $\checkmark$ & $\checkmark$ \\
Observations & 111,474 & 111,474 & 111,474 & 111,474 & 111,474 & 111,474 & 111,474 \\
\hline 
\end{tabular}
} 
\end{subtable}
\caption{\textbf{Full estimation results with matched sample}. All standard errors are clustered at the streamer level. FE = fixed effects. Significance level: *$p<0.1$; **$p<0.05$; ***$p<0.01$.}
\label{table:full-results-matched}
\end{table}

\subsection{IV-DiD}\label{appendix-iv-did}

Second, we implement a DiD specification using an instrumental variables approach. Specifically, we instrument treatment assignment with the shares of banned gambling content for banned streamers and the shares of gambling livestreams for unbanned streamers. This strategy mitigates concerns regarding selection bias stemming from systematic, unobserved differences between banned and unassigned streamers. By leveraging these shares as instruments, we isolate exogenous variation in treatment exposure driven by the streamers' pre-existing content compositions, rather than endogenous shifts in behavior or platform-wide trends that might otherwise confound the estimates in the full sample. Results are presented in Table \ref{table:full-results-ivdid}.

\begin{table}[htbp]
\centering
\resizebox{\textwidth}{!}{%
\begin{tabular}{lccccccc}
\hline
 & \multicolumn{3}{c}{Supply} & \multicolumn{4}{c}{Demand} \\ \cline{2-8}
 & Gambling & Loot Box & Streaming & Hrs Watched & Tier 1 & Tier 2 & Tier 3 \\ \hline
Banned $\times$ Post ($\hat{\beta}_1$)
    & \begin{tabular}[c]{@{}c@{}}-1.260***\\ (0.109)\end{tabular}
    & \begin{tabular}[c]{@{}c@{}}-0.093\\ (0.083)\end{tabular}
    & \begin{tabular}[c]{@{}c@{}}-0.862***\\ (0.111)\end{tabular}
    & \begin{tabular}[c]{@{}c@{}}-2.695***\\ (0.346)\end{tabular}
    & \begin{tabular}[c]{@{}c@{}}-1.161***\\ (0.169)\end{tabular}
    & \begin{tabular}[c]{@{}c@{}}-0.007\\ (0.018)\end{tabular}
    & \begin{tabular}[c]{@{}c@{}}-0.072***\\ (0.024)\end{tabular} \\
Banned $\times$ Between ($\hat{\beta}_2$)
    & \begin{tabular}[c]{@{}c@{}}0.073\\ (0.094)\end{tabular}
    & \begin{tabular}[c]{@{}c@{}}-0.047\\ (0.071)\end{tabular}
    & \begin{tabular}[c]{@{}c@{}}-0.078\\ (0.074)\end{tabular}
    & \begin{tabular}[c]{@{}c@{}}-0.430**\\ (0.209)\end{tabular}
    & \begin{tabular}[c]{@{}c@{}}-0.013\\ (0.106)\end{tabular}
    & \begin{tabular}[c]{@{}c@{}}0.010\\ (0.016)\end{tabular}
    & \begin{tabular}[c]{@{}c@{}}-0.021\\ (0.018)\end{tabular} \\
Unbanned $\times$ Post ($\hat{\beta}_3$)
    & \begin{tabular}[c]{@{}c@{}}-0.563***\\ (0.167)\end{tabular}
    & \begin{tabular}[c]{@{}c@{}}0.391***\\ (0.098)\end{tabular}
    & \begin{tabular}[c]{@{}c@{}}-0.200\\ (0.146)\end{tabular}
    & \begin{tabular}[c]{@{}c@{}}-0.810*\\ (0.449)\end{tabular}
    & \begin{tabular}[c]{@{}c@{}}-0.084\\ (0.199)\end{tabular}
    & \begin{tabular}[c]{@{}c@{}}-0.001\\ (0.026)\end{tabular}
    & \begin{tabular}[c]{@{}c@{}}0.021\\ (0.032)\end{tabular} \\
Unbanned $\times$ Between ($\hat{\beta}_4$)
    & \begin{tabular}[c]{@{}c@{}}0.066\\ (0.124)\end{tabular}
    & \begin{tabular}[c]{@{}c@{}}0.196**\\ (0.082)\end{tabular}
    & \begin{tabular}[c]{@{}c@{}}0.093\\ (0.113)\end{tabular}
    & \begin{tabular}[c]{@{}c@{}}0.226\\ (0.332)\end{tabular}
    & \begin{tabular}[c]{@{}c@{}}-0.079\\ (0.123)\end{tabular}
    & \begin{tabular}[c]{@{}c@{}}0.003\\ (0.019)\end{tabular}
    & \begin{tabular}[c]{@{}c@{}}0.008\\ (0.020)\end{tabular} \\ \hline
Streamer FE & $\checkmark$ & $\checkmark$ & $\checkmark$ & $\checkmark$ & $\checkmark$ & $\checkmark$ & $\checkmark$ \\
Week FE     & $\checkmark$ & $\checkmark$ & $\checkmark$ & $\checkmark$ & $\checkmark$ & $\checkmark$ & $\checkmark$ \\
Observations & 112,222 & 112,222 & 112,222 & 112,222 & 112,222 & 112,222 & 112,222 \\
\hline
\end{tabular}%
}
\caption{\textbf{Full estimation results with IV-DiD specification}. All standard errors are clustered at the streamer level. FE = fixed effects. Significance level: *$p<0.1$; **$p<0.05$; ***$p<0.01$.}
\label{table:full-results-ivdid}
\end{table}

\subsection{PPML}\label{appendix-ppml}

Our primary specifications employ log-transformed outcomes estimated via OLS with two-way fixed effects. This approach implicitly assumes constant error variance across units; however, our outcome variables are non-negative, right-skewed, and heteroskedastic. Because of Jensen’s inequality, log-linearizing a multiplicative model under heteroskedasticity may yield inconsistent estimates. To address this, we estimate a multiplicative model using Poisson Pseudo-Maximum Likelihood (PPML). PPML estimates the conditional mean of the level outcome directly, ensuring consistency regardless of the error distribution, provided the conditional mean is correctly specified. It also naturally accommodates zero values and uses moment conditions that down-weight extreme observations relative to log-OLS, providing greater robustness when units exhibit significantly different scales of activity. Results are presented in Table \ref{table:full-results-ppml}. Note that the causal effect on gambling hours are substantially large because PPML fits on level outcomes with extreme zero-inflation.

\begin{table}[htbp]
\centering
\resizebox{\textwidth}{!}{%
\begin{tabular}{lccccccc}
\hline
 & \multicolumn{3}{c}{Supply} & \multicolumn{4}{c}{Demand} \\ \cline{2-8}
 & Gambling & Loot Box & Streaming & Hrs Watched & Tier 1 & Tier 2 & Tier 3 \\ \hline
Banned $\times$ Post ($\hat{\beta}_1$)
    & \begin{tabular}[c]{@{}c@{}}-17.013***\\ (0.149)\end{tabular}
    & \begin{tabular}[c]{@{}c@{}}-0.044\\ (0.075)\end{tabular}
    & \begin{tabular}[c]{@{}c@{}}-0.394***\\ (0.053)\end{tabular}
    & \begin{tabular}[c]{@{}c@{}}-0.422***\\ (0.098)\end{tabular}
    & \begin{tabular}[c]{@{}c@{}}-0.127\\ (0.089)\end{tabular}
    & \begin{tabular}[c]{@{}c@{}}-0.099\\ (0.096)\end{tabular}
    & \begin{tabular}[c]{@{}c@{}}-0.382***\\ (0.133)\end{tabular} \\
Banned $\times$ Between ($\hat{\beta}_2$)
    & \begin{tabular}[c]{@{}c@{}}0.091\\ (0.065)\end{tabular}
    & \begin{tabular}[c]{@{}c@{}}-0.084\\ (0.064)\end{tabular}
    & \begin{tabular}[c]{@{}c@{}}-0.052\\ (0.035)\end{tabular}
    & \begin{tabular}[c]{@{}c@{}}-0.217**\\ (0.086)\end{tabular}
    & \begin{tabular}[c]{@{}c@{}}-0.139\\ (0.092)\end{tabular}
    & \begin{tabular}[c]{@{}c@{}}-0.064\\ (0.083)\end{tabular}
    & \begin{tabular}[c]{@{}c@{}}-0.333**\\ (0.159)\end{tabular} \\
Unbanned $\times$ Post ($\hat{\beta}_3$)
    & \begin{tabular}[c]{@{}c@{}}-15.858***\\ (0.129)\end{tabular}
    & \begin{tabular}[c]{@{}c@{}}-0.032\\ (0.062)\end{tabular}
    & \begin{tabular}[c]{@{}c@{}}-0.067\\ (0.043)\end{tabular}
    & \begin{tabular}[c]{@{}c@{}}-0.410**\\ (0.183)\end{tabular}
    & \begin{tabular}[c]{@{}c@{}}-0.163\\ (0.128)\end{tabular}
    & \begin{tabular}[c]{@{}c@{}}-0.240***\\ (0.090)\end{tabular}
    & \begin{tabular}[c]{@{}c@{}}-0.460\\ (0.303)\end{tabular} \\
Unbanned $\times$ Between ($\hat{\beta}_4$)
    & \begin{tabular}[c]{@{}c@{}}0.155**\\ (0.061)\end{tabular}
    & \begin{tabular}[c]{@{}c@{}}-0.013\\ (0.042)\end{tabular}
    & \begin{tabular}[c]{@{}c@{}}-0.023\\ (0.025)\end{tabular}
    & \begin{tabular}[c]{@{}c@{}}-0.214*\\ (0.110)\end{tabular}
    & \begin{tabular}[c]{@{}c@{}}-0.099\\ (0.075)\end{tabular}
    & \begin{tabular}[c]{@{}c@{}}0.108\\ (0.069)\end{tabular}
    & \begin{tabular}[c]{@{}c@{}}-0.243\\ (0.271)\end{tabular} \\ \hline
Streamer FE & $\checkmark$ & $\checkmark$ & $\checkmark$ & $\checkmark$ & $\checkmark$ & $\checkmark$ & $\checkmark$ \\
Week FE     & $\checkmark$ & $\checkmark$ & $\checkmark$ & $\checkmark$ & $\checkmark$ & $\checkmark$ & $\checkmark$ \\
Observations & 112,222 & 112,222 & 112,222 & 112,222 & 112,222 & 112,222 & 112,222 \\
\hline
\end{tabular}%
}
\caption{\textbf{Full estimation results with PPML}. All standard errors are clustered at the streamer level. FE = fixed effects. Significance level: *$p<0.1$; **$p<0.05$; ***$p<0.01$.}
\label{table:full-results-ppml}
\end{table}

\subsection{Size-Weighted DiD}\label{appendix-weighted-did}

Finally, we conduct a DiD specification that weights each streamer by their popularity, measured by average concurrent viewers (ACV). This estimator serves as a robustness check against potential biases inherent in log-linear specifications. By reweighting observations based on pre-period demand, this approach accounts for systematic differences in baseline scale and mitigates the concern that estimates may be disproportionately driven by low-demand units with high fractions of zero observations. This ensures that the regression’s influence is aligned with the actual economic importance of the streamers in our sample. Results are presented in Table \ref{table:full-results-wdid}. Note that the causal effect on total hours watched and Tier 1 subscriptions are amplified relative to baseline TWFE, suggesting the baseline estimates may be attenuated by measurement error in treatment assignment. 

\begin{table}[htbp]
\centering
\resizebox{\textwidth}{!}{%
\begin{tabular}{lccccccc}
\hline
 & \multicolumn{3}{c}{Supply} & \multicolumn{4}{c}{Demand} \\ \cline{2-8}
 & Gambling & Loot Box & Streaming & Hrs Watched & Tier 1 & Tier 2 & Tier 3 \\ \hline
Banned $\times$ Post ($\hat{\beta}_1$)
    & \begin{tabular}[c]{@{}c@{}}-1.170***\\ (0.284)\end{tabular}
    & \begin{tabular}[c]{@{}c@{}}0.196\\ (0.129)\end{tabular}
    & \begin{tabular}[c]{@{}c@{}}-0.255***\\ (0.080)\end{tabular}
    & \begin{tabular}[c]{@{}c@{}}-0.386***\\ (0.130)\end{tabular}
    & \begin{tabular}[c]{@{}c@{}}0.140\\ (0.147)\end{tabular}
    & \begin{tabular}[c]{@{}c@{}}0.043\\ (0.051)\end{tabular}
    & \begin{tabular}[c]{@{}c@{}}0.062\\ (0.052)\end{tabular} \\
Banned $\times$ Between ($\hat{\beta}_2$)
    & \begin{tabular}[c]{@{}c@{}}-0.128\\ (0.320)\end{tabular}
    & \begin{tabular}[c]{@{}c@{}}-0.132\\ (0.134)\end{tabular}
    & \begin{tabular}[c]{@{}c@{}}-0.045\\ (0.071)\end{tabular}
    & \begin{tabular}[c]{@{}c@{}}-0.025\\ (0.075)\end{tabular}
    & \begin{tabular}[c]{@{}c@{}}0.263***\\ (0.092)\end{tabular}
    & \begin{tabular}[c]{@{}c@{}}0.064\\ (0.047)\end{tabular}
    & \begin{tabular}[c]{@{}c@{}}-0.006\\ (0.078)\end{tabular} \\
Unbanned $\times$ Post ($\hat{\beta}_3$)
    & \begin{tabular}[c]{@{}c@{}}-0.392***\\ (0.115)\end{tabular}
    & \begin{tabular}[c]{@{}c@{}}-0.160\\ (0.212)\end{tabular}
    & \begin{tabular}[c]{@{}c@{}}-0.235**\\ (0.104)\end{tabular}
    & \begin{tabular}[c]{@{}c@{}}-0.408**\\ (0.187)\end{tabular}
    & \begin{tabular}[c]{@{}c@{}}-0.101\\ (0.078)\end{tabular}
    & \begin{tabular}[c]{@{}c@{}}-0.067\\ (0.094)\end{tabular}
    & \begin{tabular}[c]{@{}c@{}}-0.034\\ (0.073)\end{tabular} \\
Unbanned $\times$ Between ($\hat{\beta}_4$)
    & \begin{tabular}[c]{@{}c@{}}-0.075\\ (0.103)\end{tabular}
    & \begin{tabular}[c]{@{}c@{}}0.037\\ (0.106)\end{tabular}
    & \begin{tabular}[c]{@{}c@{}}-0.066\\ (0.060)\end{tabular}
    & \begin{tabular}[c]{@{}c@{}}-0.124\\ (0.086)\end{tabular}
    & \begin{tabular}[c]{@{}c@{}}0.057\\ (0.075)\end{tabular}
    & \begin{tabular}[c]{@{}c@{}}0.078**\\ (0.034)\end{tabular}
    & \begin{tabular}[c]{@{}c@{}}0.060**\\ (0.030)\end{tabular} \\ \hline
Streamer FE & $\checkmark$ & $\checkmark$ & $\checkmark$ & $\checkmark$ & $\checkmark$ & $\checkmark$ & $\checkmark$ \\
Week FE     & $\checkmark$ & $\checkmark$ & $\checkmark$ & $\checkmark$ & $\checkmark$ & $\checkmark$ & $\checkmark$ \\
Observations & 112,222 & 112,222 & 112,222 & 112,222 & 112,222 & 112,222 & 112,222 \\
\hline
\end{tabular}%
}
\caption{\textbf{Full estimation results with weighted DiD specification}. All standard errors are clustered at the streamer level. FE = fixed effects. Significance level: *$p<0.1$; **$p<0.05$; ***$p<0.01$.}
\label{table:full-results-wdid}
\end{table}

\clearpage
\newpage

\section{Minor Content and Viewer Detection}\label{appendix-minor-detection}

\subsection{Minor Content Detection}\label{appendix-minor-content-detection}

\subsubsection{Pilot Study}\label{appendix-minor-content-pilot}
We first conducted a pilot round of classification on 50 randomly selected content from the list of 10,970 unique content covered by our sample using LLMs, given their strong performance for diverse downstream tasks. We labeled whether each content attracts minors, with explanations and citations whenever possible, using four different model configurations from OpenAI: Model 1 - GPT-4o without web search, Model 2 - GPT-4o with web search, Model 3 - GPT o3 with web search (medium reasoning), and Model 4 - GPT o3 with web search (high reasoning).\footnote{For all models with web search enabled, we used high search context size. For all models with a ``temperature'' parameter, we explored setting temperature to either 0 (which ensures output consistency) or 0.2 (which allows light stochastic responses over different iterations).
Note that GPT-4o does not have a reasoning effort parameter.} Below shows the prompts we used based on the type of content streamed on Twitch, which can either be games or non-gaming categories such as ``Just Chatting'' and ``Pool, Hot Tubs, and Beaches''. When the content is a video game, we supply key game attributes (if available) to the prompts to help guide the web search to output accurate results: 

\begin{promptblock}[Case 1: Prompt for detecting minor content in video games]
You are an experienced, highly qualified video game professional. You are familiar with the target audience of each video game.  

You are provided with the name of a video game: \{game\_name\}. The video game is developed by \{developer\}, belongs to the \{genres\} genre(s), 
is available on gaming platforms such as \{gaming\_platforms\}, and carries the following age ratings(s) information: \{age\_ratings\}. 

Your task is to determine whether the game attracts a large and active player or viewer base under the age of 18, with explanations. Here is a list of criteria for your evaluation: 

\begin{enumerate}[leftmargin=*, nosep, topsep=2pt, itemsep=1pt, parsep=0pt]
  \item A significant proportion of its frequent or daily players or viewers are teenagers or children.
  \item The game is well-known and regularly discussed among youth in surveys, social media, or gaming communities.
  \item It maintains high average playtime and repeated logins by young players.
  \item It is available on widely used devices (e.g., mobile, tablets, consoles) and often free-to-play.
  \item The game influences youth trends, online content (e.g., YouTube, TikTok), or is reflected in schoolyard discussions and digital communities.
  \item The game often appears in top rankings for children’s or teen gaming behavior, as reported by analytics firms, parental control services, or platforms like Roblox, Twitch, or YouTube.
\end{enumerate}

Respond using JSON format only: 

\begin{lstlisting}[language=json]
{{
  "game_name": "{game_name}"
  "target_people_under_18": "[Yes or No]"
  "explanation":"[explanation of your evaluation in 2-3 sentences, with citations or links to back up your explanation whenever possible]"}}
\end{lstlisting}

Don’t start with any introduction and provide only json output. Start with ```json.
\end{promptblock}

\begin{promptblock}[Case 2: Prompt for detecting minor content in non-gaming categories]
You are an experienced, highly qualified platform analyst. You are familiar with the target audience of each streaming category on Twitch. 

You are provided with the name of a non-gaming related stream category: \{name\}. Your task is to determine whether it attracts a large and active viewer base under the age of 18, with explanations. 

Respond using JSON format only: 

\begin{lstlisting}[language=json]
{{
  "name": "{name}"
  "target_people_under_18": "[Yes or No]"
  "explanation":"[explanation of your evaluation in 2-3 sentences, with citations or links to back up your explanation whenever possible]"}}
\end{lstlisting}

Don’t start with any introduction and provide only json output. Start with ```json.
\end{promptblock}

We manually inspect through the 50 responses across the different model configurations. We find that, in general, enabling web search provides sufficiently accurate labels and explanations that reflect content’s appealed audience. Among the models with web search enabled, ``GPT-o3'' - known for its superior logical reasoning - offers the most accurate responses with the best explanations. Using either high or medium reasoning for ``GPT-o3'' offers comparable results, with the medium reasoning one providing slightly more most conservative response in terms of classifying content whose target audience is not minors. Thus, we proceed the main classification task with GPT o3 with web search (high search context size, medium reasoning) as the main model. 



\subsubsection{Main Study}\label{appendix-minor-content-main}

For our main classification, we follow a 2-step procedure to detect minor-targeted games by first obtaining a list of minor games (which may include popular games that attract broad audience not just minors) and then evaluate the share of minors attracted by each game to obtain a final list of games that mainly target minors. For step 1, we used the same prompts as in the pilot study. For step 2, the prompt we used for step 2 is shown below, with key game attributes information supplemented if available. For this step, we evaluate minor shares using ``GPT-o3'' both with and without web search, and take the average of the two for subsequent analysis.\footnote{The correlation in the two versions are 0.62.} This helps alleviate the concern that publicly available stats (which web search model uses) may undercount minors due to privacy rules.

\begin{promptblock}[Prompt: Detect the share of minors for each game]
You are an experienced, highly qualified video game professional. You are familiar with the audience composition of each video game.  

You are provided with the name of a video game: \{video\_game\}. The video game is developed by \{developer\}, belongs to the \{genres\} genre(s), 
is available on gaming platforms such as \{gaming\_platforms\}, and carries the following age ratings(s) information: \{age\_ratings\}. 

Your task is to determine the share of its player base that is under the age of 18 based on the hours played, with explanations. 

The share you output needs to fall in between 0 and 100. 

Respond using JSON format only: 
\begin{lstlisting}[language=json]
{{
  "video_game": "{video_game}"
  "share_people_under_18": "your evaluation"
  "explanation":"[explanation of your evaluation in 2-3 sentences, with citations or links to back up your explanation whenever possible]"}}
\end{lstlisting}

Don’t start with any introduction and provide only json output. Start with ```json.
\end{promptblock}

We tested different cutoff thresholds and use the list at the threshold of 50 as our final minor game list, ensuring the games satisfy the following criteria: first, the list must capture popular games which are known for highly targeting minors. This include, for example, ``Roblox", ``Minecraft" and ``Brawl Stars", which also contribute to the majority of consumption of minor-targeted content in our sample. These games help provide high identifying power on detecting minor viewers in our sample. Second, the list must also contain other games, which are less influential but are also clearly minor-targeted, providing coverage without diluting accuracy. Finally, at the same time, the list must exclude popular games that attract broad audience, not just minors (e.g., FIFA, NBA 2K), which reduces noise. 

\newpage
\subsection{Minor Viewer Detection}\label{appendix-minor-viewer-detection}

To identity likely minor and adult viewers, we conduct two one-sided binomial tests to reject viewers who are not minors or adults. To achieve this, for each viewer $i$, we examine their pre-policy minor sessions viewing behavior. Specifically, let \(n_i\) denote the total number of sessions watched and \(k_i\) the number of minor sessions watched, and thus \(p_i\) is the share of minor sessions watched by viewer $i$ on Twitch prior to the policy. Let the benchmark population proportion be  \(p_0=0.054\), which is the mean share of minor sessions watched across viewers in the pre-policy period as shown in Table \ref{table:viewer-stats-pre-policy}. We model \(X_i\sim\mathrm{Binomial}(n_i,p_i)\) under the null \(p_i=p_0\) and perform two one–sided tests at the significance level of $\alpha=0.05$ with the hypotheses specified as follows.

For right-tail test: 
\[
\left\{
\begin{aligned}
H_0 &: \; p_i \le p_0 = 0.054, \text{i.e., the viewer is not a minor} \\
H_1 &: \; p_i >  p_0 = 0.054.
\end{aligned}
\right.
\]

For left-tail test: 
\[
\left\{
\begin{aligned}
H_0 &: \; p_i \ge p_0 = 0.054, \text{i.e., the viewer is not an adult}\\
H_1 &: \; p_i <  p_0 = 0.054.
\end{aligned}
\right.
\]

Intuitively, this means that (1) if we can reject $p_i \le p_0$ in favor of 
$p_i >  p_0$,  the viewer's observed share of minor sessions watched is unusually high and thus we label her as minor; (2) if we can reject that $p_i \ge p_0$ in favor of $p_i <  p_0$, the viewer's observed share of minor sessions watched is unusually high and thus we label her as adult; if neither test rejects, we are uncertain about the identity of the viewer and thus label her into the uncertain category. 

To find the rejection regions, we note that we only need to compute the two critical values for every unique $n$ that is observed in our data. Specifically, for every unique $n$, we find
\[
x_{\mathrm{right}}(n)\;=\;\min\big\{x\in\{0,\ldots,n\}:\Pr[X\ge x\,|\,n,p_0]\le 0.05\big\}
\]
\[
x_{\mathrm{left}}(n)\;=\;\max\big\{x\in\{0,\ldots,n\}:\Pr[X\le x\,|\,n,p_0]\le 0.05\big\}
\]
These values determine the smallest (respectively, largest) number of minor sessions watched that would reject the null on the right (respectively, left) side. Note that because the binomial distribution is discrete, a one–sided rejection region may not exist for some \((n,p_0,0.05)\). This happens when $n<2$ for the right-tail test and $n<55$ for the left-tail test. In such cases the corresponding critical value is undefined and that tail provides no basis for rejection, and so we classify viewers falling under that scenario also into uncertain. 

Finally, given the critical values for the relevant \(n_i\), we classify each viewer by comparing their observed minor sessions watched \(k_i\) to the critical values: 
\[
\text{label}(i)=
\begin{cases}
\text{minor}, & \text{if } x_{\mathrm{right}}(n_i)\text{ exists and } k_i \ge x_{\mathrm{right}}(n_i),\\[2pt]
\text{adult}, & \text{if } x_{\mathrm{left}}(n_i)\text{ exists and }  k_i \le x_{\mathrm{left}}(n_i),\\[2pt]
\text{uncertain}, & \text{otherwise}.
\end{cases}
\]

In total, out of the 14,568,807 viewers observed in our data, we have 5.99\% (872,672) classified as minors, 4.42\% (643,941) as adults, and the remaining 89.59\% (13,052,194) falling into uncertain. 



\begin{table}[!h]
\small
\centering
\renewcommand{\arraystretch}{1.2}
\begin{threeparttable}
\begin{tabular}{lrrr}
\hline
Statistics &
\makecell{Total sessions watched} &
\makecell{Total minor sessions watched} &
\makecell{Share of minor \\ sessions watched} \\
\midrule
Mean                & 15.42 & 0.77 & 0.05 \\
Std                 & 69.38 & 4.60 & 0.18 \\
Min                 & 1.00  & 0.00 & 0.00 \\
%25\textsubscript{quantile} & 2.00  & 0.00 & 0.00 \\
Medium & 6.00  & 0.00 & 0.00 \\
%75\textsubscript{quantile} & 16.00 & 0.00 & 0.00 \\
Max                 & 168,519.00 & 5,261.00 & 1.00 \\
%Nonzero viewer share  & 1.00 & 0.15 & 0.15 \\
Nonzero mean        & 15.42 & 5.05 & 0.35 \\
%Nonzero median      & 6.00  & 2.00 & 0.25 \\
\hline
\end{tabular}
\end{threeparttable}
\caption{\textbf{Summary statistics on viewer watched sessions (pre-policy announcement)} On average, 0.05\% of the sessions watched by minors are about minor content. Conditional on viewers who watched minor sessions (around 15\% of viewers), 35\% of the sessions watched about minor content.}
\label{table:viewer-stats-pre-policy}
\end{table}

\newpage
\begin{figure}[!htb]
    \centering
    \includegraphics[width=0.75\linewidth]{Figures/Descriptives/distribution-num-sessions.png}
    \caption{\textbf{Distribution of the number of sessions watched across viewers (pre-policy)}. Sessions are capped at 100 sessions, with the bar at 100 summarizing the right tail mass. The mean number of sessions watched across viewers is 15.42.}
    \label{distribution-session-watched}
\end{figure}
\begin{figure}[!htb]
    \centering
    \includegraphics[width=0.75\linewidth]{Figures/Descriptives/distribution-minor-share.png}
    \caption{\textbf{Distribution of the share of minor sessions watched across viewers (pre-policy; zero share excluded)}. The distribution excludes the share of minor sessions that are zero, which appear for 84.7\% of viewers. Spikes at shares such as 25\%,50\%,100\% arise from small denominator - i.e., the total number of sessions watched for some viewers is very small. For example, if a viewer only watches one content session and that session happens to be minor content, then the share would be 100\%. We discuss our strategy in regards to this issue for minor viewer classification in Section \ref{section-minor}.}
    \label{distribution-minor-share-watched}
\end{figure}








\newpage

\section{Additional Estimation Results}\label{appendix_full_results}

In this section, we present full estimation results of supply-side and demand-side outcome variables.

% Table F.0.1 has been revised on Sept 21.
\begin{table}[!htb]
\renewcommand{\arraystretch}{1.1}
\centering
\begin{subtable}{\textwidth}
\centering 
\caption{Estimates for Gambling and Loot Box Games}
\label{subtable:a}
\begin{tabular}{lcc@{\hspace{0.5cm}}cc}
\hline 
                            & \multicolumn{2}{c}{log(Gambling hours + 1)}                                                                                 & \multicolumn{2}{c}{log(Loot Box Games+ 1)}                                                                                      \\ 
\cline{2-5}
                            & TWFE              & SynthDiD          & TWFE               & SynthDiD          \\ 
\hline
Banned $\times$ Post ($\beta_1$)                   & \begin{tabular}[c]{@{}c@{}}-1.049***\\ (0.088)\end{tabular} & \begin{tabular}[c]{@{}c@{}}-1.001***\\ (0.082)\end{tabular} & \begin{tabular}[c]{@{}c@{}}0.013\\ (0.066)\end{tabular}  & \begin{tabular}[c]{@{}c@{}}0.021\\ (0.054)\end{tabular}  \\
Banned $\times$ Between ($\beta_2$)                   & \begin{tabular}[c]{@{}c@{}}0.171*\\ (0.074)\end{tabular}    & \begin{tabular}[c]{@{}c@{}}0.083\\ (0.070)\end{tabular}     & \begin{tabular}[c]{@{}c@{}}-0.002\\ (0.054)\end{tabular} & \begin{tabular}[c]{@{}c@{}}0.011\\ (0.050)\end{tabular}  \\
Unbanned $\times$ Post ($\beta_3$)                   & \begin{tabular}[c]{@{}c@{}}-0.108**\\ (0.049)\end{tabular}   & \begin{tabular}[c]{@{}c@{}}-0.130***\\ (0.043)\end{tabular} & \begin{tabular}[c]{@{}c@{}}-0.044\\ (0.043)\end{tabular} & \begin{tabular}[c]{@{}c@{}}-0.054\\ (0.036)\end{tabular} \\
Unbanned $\times$ Between ($\beta_4$)                   & \begin{tabular}[c]{@{}c@{}}0.133***\\ (0.039)\end{tabular}  & \begin{tabular}[c]{@{}c@{}}0.102\\ (0.037)***\end{tabular}     & \begin{tabular}[c]{@{}c@{}}-0.026\\ (0.039)\end{tabular} & \begin{tabular}[c]{@{}c@{}}-0.009\\ (0.036)\end{tabular} \\
\hline
Streamer FE             & $\checkmark$ & $\checkmark$ & $\checkmark$ & $\checkmark$ \\
Week FE                 & $\checkmark$ & $\checkmark$ & $\checkmark$ & $\checkmark$ \\
Observations            & 112,222      & 112,222      & 112,222      & 112,222      \\
Mean Dependent Variable & 1.130        & 1.130        & 1.341        & 1.341        \\ 
\hline 
\end{tabular}
\end{subtable}
\vspace{0.15cm}
\begin{subtable}{\textwidth}
\centering 
\caption{Estimates for Streaming Hours and Other Games}
\label{subtable:b}
\begin{tabular}{lcc@{\hspace{0.5cm}}cc} % Reduced space to 0.5cm for better alignment
\hline 
                            & \multicolumn{2}{c}{log(Streaming hours + 1)}                                                                                & \multicolumn{2}{c}{log(Other Games+ 1)}                                                                                          \\ 
\cline{2-5}
                            & TWFE              & SynthDiD          & TWFE               & SynthDiD          \\ 
\midrule
Banned $\times$ Post ($\beta_1$)                   & \begin{tabular}[c]{@{}c@{}}-0.606***\\ (0.083)\end{tabular} & \begin{tabular}[c]{@{}c@{}}-0.585***\\ (0.075)\end{tabular} & \begin{tabular}[c]{@{}c@{}}-0.145***\\ (0.04)\end{tabular}  & \begin{tabular}[c]{@{}c@{}}-0.149***\\ (0.032)\end{tabular} \\
Banned $\times$ Between ($\beta_2$)                   & \begin{tabular}[c]{@{}c@{}}-0.047\\ (0.056)\end{tabular}    & \begin{tabular}[c]{@{}c@{}}-0.041\\ (0.052)\end{tabular}    & \begin{tabular}[c]{@{}c@{}}-0.128***\\ (0.037)\end{tabular}  & \begin{tabular}[c]{@{}c@{}}-0.098***\\ (0.033)\end{tabular} \\
Unbanned $\times$ Post ($\beta_3$)                   & \begin{tabular}[c]{@{}c@{}}-0.172***\\ (0.045)\end{tabular} & \begin{tabular}[c]{@{}c@{}}-0.194***\\ (0.040)\end{tabular} & \begin{tabular}[c]{@{}c@{}}-0.133***\\ (0.032)\end{tabular}  & \begin{tabular}[c]{@{}c@{}}-0.105***\\ (0.034)\end{tabular} \\
Unbanned $\times$ Between ($\beta_4$)                   & \begin{tabular}[c]{@{}c@{}}-0.059\\ (0.039)\end{tabular}    & \begin{tabular}[c]{@{}c@{}}-0.062\\ (0.037)\end{tabular}    & \begin{tabular}[c]{@{}c@{}}-0.156***\\ (0.028)\end{tabular} & \begin{tabular}[c]{@{}c@{}}-0.141***\\ (0.026)\end{tabular} \\
\hline
Streamer FE             & $\checkmark$ & $\checkmark$ & $\checkmark$ & $\checkmark$ \\
Week FE                 & $\checkmark$ & $\checkmark$ & $\checkmark$ & $\checkmark$ \\
Observations            & 112,222      & 112,222      & 112,222      & 112,222      \\
Mean Dependent Variable & 2.722        & 2.722        & 0.540        & 0.540        \\ 
\hline 
\end{tabular}
\end{subtable}
\vspace{-0.3cm}
\caption{\small \textbf{Full estimation results of supply-side outcomes}. All standard errors are clustered at the streamer level. The mean dependent variable is calculated based on observations of all streamers before the policy announcement. FE = fixed effects. Significance level: *$p<0.1$; **$p<0.05$; ***$p<0.01$.}
\label{table:full-supply-results}
\end{table}

\newpage

\begin{table}[!htb]
\renewcommand{\arraystretch}{1.2}
\centering
\begin{subtable}{\textwidth}
\centering
\caption{Estimates for Hours Watched and Tier 1 Subscriptions}
\label{subtable:hours-watched-tier1}
\begin{tabular}{lcc@{\hspace{1cm}}cc}
\hline 
                            & \multicolumn{2}{c}{log(Hours Watched + 1)} & \multicolumn{2}{c}{log(Tier 1+1)} \\ \cline{2-3} \cline{4-5} 
                            & TWFE       & SynthDiD  & TWFE       & SynthDiD  \\ \hline
Banned $\times$ Post ($\beta_1$)                   & \begin{tabular}[c]{@{}c@{}}-1.694***\\ (0.268)\end{tabular} & \begin{tabular}[c]{@{}c@{}}-1.696***\\ (0.241)\end{tabular} & \begin{tabular}[c]{@{}c@{}}-0.59***\\ (0.121)\end{tabular}  & \begin{tabular}[c]{@{}c@{}}-0.584***\\ (0.110)\end{tabular} \\
Banned $\times$ Between ($\beta_2$)                   & \begin{tabular}[c]{@{}c@{}}-0.268\\ (0.173)\end{tabular}    & \begin{tabular}[c]{@{}c@{}}-0.272\\ (0.160)\end{tabular}    & \begin{tabular}[c]{@{}c@{}}0.046\\ (0.087)\end{tabular}     & \begin{tabular}[c]{@{}c@{}}0.040\\ (0.078)\end{tabular}     \\
Unbanned $\times$ Post ($\beta_3$)                   & \begin{tabular}[c]{@{}c@{}}-0.523***\\ (0.134)\end{tabular} & \begin{tabular}[c]{@{}c@{}}-0.554***\\ (0.122)\end{tabular} & \begin{tabular}[c]{@{}c@{}}-0.178**\\ (0.058)\end{tabular}  & \begin{tabular}[c]{@{}c@{}}-0.185***\\ (0.054)\end{tabular}  \\
Unbanned $\times$ Between ($\beta_4$)                   & \begin{tabular}[c]{@{}c@{}}-0.218*\\ (0.117)\end{tabular}    & \begin{tabular}[c]{@{}c@{}}-0.223\\ (0.113)\end{tabular}    & \begin{tabular}[c]{@{}c@{}}-0.181***\\ (0.053)\end{tabular} & \begin{tabular}[c]{@{}c@{}}-0.171***\\ (0.046)\end{tabular}  \\ \hline
Streamer FE             & $\checkmark$ & $\checkmark$ & $\checkmark$ & $\checkmark$ \\
Week FE                 & $\checkmark$ & $\checkmark$ & $\checkmark$ & $\checkmark$ \\
Observations            & 112,222      & 112,222      & 112,222      & 112,222      \\
Mean Dependent Variable & 9.133        & 9.133        & 2.875        & 2.875        \\ \hline 
\end{tabular}
\end{subtable}
\vspace{0.5cm}
\begin{subtable}{\textwidth}
\centering
\caption{Estimates for Tier 2 and Tier 3 Subscriptions}
\label{subtable:tier2-tier3}
\begin{tabular}{lcc@{\hspace{1cm}}cc}
\hline 
                            & \multicolumn{2}{c}{log(Tier 2+1)} & \multicolumn{2}{c}{log(Tier 3+1)} \\ \cline{2-3} \cline{4-5} 
                            & TWFE       & SynthDiD  & TWFE       & SynthDiD  \\ \hline
Banned $\times$ Post ($\beta_1$)                   & \begin{tabular}[c]{@{}c@{}}0.01\\ (0.013)\end{tabular}    & \begin{tabular}[c]{@{}c@{}}0.017\\ (0.011)\end{tabular}   & \begin{tabular}[c]{@{}c@{}}-0.024\\ (0.015)\end{tabular} & \begin{tabular}[c]{@{}c@{}}-0.024\\ (0.014)\end{tabular} \\
Banned $\times$ Between ($\beta_2$)                   & \begin{tabular}[c]{@{}c@{}}0.004\\ (0.013)\end{tabular}   & \begin{tabular}[c]{@{}c@{}}0.010\\ (0.013)\end{tabular}   & \begin{tabular}[c]{@{}c@{}}-0.018\\ (0.012)\end{tabular} & \begin{tabular}[c]{@{}c@{}}-0.015\\ (0.012)\end{tabular} \\
Unbanned $\times$ Post ($\beta_3$)                   & \begin{tabular}[c]{@{}c@{}}0.012\\ (0.011)\end{tabular}   & \begin{tabular}[c]{@{}c@{}}0.013\\ (0.009)\end{tabular}   & \begin{tabular}[c]{@{}c@{}}0.013\\ (0.012)\end{tabular}  & \begin{tabular}[c]{@{}c@{}}0.011\\ (0.011)\end{tabular}  \\
Unbanned $\times$ Between ($\beta_4$)                   & \begin{tabular}[c]{@{}c@{}}0.023**\\ (0.009)\end{tabular} & \begin{tabular}[c]{@{}c@{}}0.027**\\ (0.009)\end{tabular} & \begin{tabular}[c]{@{}c@{}}0.003\\ (0.008)\end{tabular}  & \begin{tabular}[c]{@{}c@{}}0.003\\ (0.009)\end{tabular}  \\ \hline
Streamer FE             & $\checkmark$ & $\checkmark$ & $\checkmark$ & $\checkmark$ \\
Week FE                 & $\checkmark$ & $\checkmark$ & $\checkmark$ & $\checkmark$ \\
Observations            & 112,222      & 112,222      & 112,222      & 112,222      \\
Mean Dependent Variable & 0.181        & 0.181        & 0.178        & 0.178        \\ \hline 
\end{tabular}
\end{subtable}
\caption{\small \textbf{Full estimation results of demand-side outcomes}. All standard errors are clustered at the streamer level. The mean dependent variable is calculated based on observations of all streamers before the policy announcement. FE = fixed effects. Significance level: *$p<0.1$; **$p<0.05$; ***$p<0.01$.}
\label{table:full-demand-results}
\end{table}


\newpage


\begin{table}[!htb]
\renewcommand{\arraystretch}{1.2}
\centering
\begin{subtable}{\textwidth}\centering
\caption{Supply-Side Outcomes}
\label{subtable:supply-side-outcomes}
\begin{tabular}{lcccc}
\hline 
            & Gambling                                                    & Loot Box                                                  & Streaming Hours                                             & Other Games                                                 \\ \cline{1-5} 
Banned $\times$ Post ($\beta_1$)     & \begin{tabular}[c]{@{}c@{}}-1.347***\\ (0.164)\end{tabular} & \begin{tabular}[c]{@{}c@{}}-0.127\\ (0.145)\end{tabular} & \begin{tabular}[c]{@{}c@{}}-0.769***\\ (0.180)\end{tabular} & \begin{tabular}[c]{@{}c@{}}-0.147*\\ (0.0797)\end{tabular}   \\
Banned $\times$ Between ($\beta_2$)     & \begin{tabular}[c]{@{}c@{}}0.027\\ (0.132)\end{tabular}     & \begin{tabular}[c]{@{}c@{}}0.011\\ (0.123)\end{tabular}  & \begin{tabular}[c]{@{}c@{}}-0.036\\ (0.135)\end{tabular}    & \begin{tabular}[c]{@{}c@{}}-0.221***\\ (0.075)\end{tabular}  \\
Unbanned $\times$ Post ($\beta_3$)     & \begin{tabular}[c]{@{}c@{}}-0.065\\ (0.07)\end{tabular}    & \begin{tabular}[c]{@{}c@{}}0.008\\ (0.081)\end{tabular}  & \begin{tabular}[c]{@{}c@{}}-0.217***\\ (0.072)\end{tabular}  & \begin{tabular}[c]{@{}c@{}}-0.138**\\ (0.055)\end{tabular}   \\
Unbanned $\times$ Between ($\beta_4$)      & \begin{tabular}[c]{@{}c@{}}0.176***\\ (0.063)\end{tabular}     & \begin{tabular}[c]{@{}c@{}}0\\ (0.081)\end{tabular}  & \begin{tabular}[c]{@{}c@{}}-0.093\\ (0.069)\end{tabular}    & \begin{tabular}[c]{@{}c@{}}-0.265***\\ (0.065)\end{tabular} \\ \hline 
Streamer FE & $\checkmark$                                                         & $\checkmark$                                                      & $\checkmark$                                                         & $\checkmark$                                                         \\
Week FE     & $\checkmark$                                                         & $\checkmark$                                                      & $\checkmark$                                                         & $\checkmark$                                                         \\
Observations & 10,494                                                       & 10,494                                                    & 10,494                                                   
& 10,494                                                       \\
Mean Dependent Variable & 1.470 & 1.210 & 2.700 & 0.437 \\
\hline 
\end{tabular}
\end{subtable}
\vspace{0.5cm}
\begin{subtable}{\textwidth}\centering
\caption{Demand-Side Outcomes}
\label{subtable:demand-side-outcomes}
\resizebox{\textwidth}{!}{%
\begin{tabular}{lcccc}
\hline 
            & \multicolumn{1}{c}{log(Hours Watched+1)} & \multicolumn{1}{c}{log(Tier 1+1)} & \multicolumn{1}{c}{log(Tier 2+1)} & \multicolumn{1}{c}{log(Tier 3+1)} \\ \hline
Banned $\times$ Post ($\beta_1$)     & \begin{tabular}[c]{@{}c@{}}-2.163***\\ (0.547)\end{tabular} & \begin{tabular}[c]{@{}c@{}}-0.614**\\ (0.280)\end{tabular} & \begin{tabular}[c]{@{}c@{}}-0.018\\ (0.023)\end{tabular} & \begin{tabular}[c]{@{}c@{}}-0.021\\ (0.035)\end{tabular} \\
Banned $\times$ Between ($\beta_2$)     & \begin{tabular}[c]{@{}c@{}}-0.021\\ (0.401)\end{tabular} & \begin{tabular}[c]{@{}c@{}}0.318\\ (0.206)\end{tabular} & \begin{tabular}[c]{@{}c@{}}0.010\\ (0.024)\end{tabular} & \begin{tabular}[c]{@{}c@{}}-0.016\\ (0.026)\end{tabular} \\
Unbanned $\times$ Post ($\beta_3$)     & \begin{tabular}[c]{@{}c@{}}-0.764***\\ (0.244)\end{tabular} & \begin{tabular}[c]{@{}c@{}}-0.214**\\ (0.102)\end{tabular} & \begin{tabular}[c]{@{}c@{}}0.017\\ (0.02)\end{tabular} & \begin{tabular}[c]{@{}c@{}}0.04**\\ (0.018)\end{tabular} \\
Unbanned $\times$ Between ($\beta_4$)     & \begin{tabular}[c]{@{}c@{}}-0.328\\ (0.209)\end{tabular} & \begin{tabular}[c]{@{}c@{}}-0.271***\\ (0.103)\end{tabular} & \begin{tabular}[c]{@{}c@{}}0.043**\\ (0.021)\end{tabular} & \begin{tabular}[c]{@{}c@{}}0.026\\ (0.02)\end{tabular} \\ \hline
Streamer FE & $\checkmark$ & $\checkmark$ & $\checkmark$ & $\checkmark$ \\
Week FE     & $\checkmark$ & $\checkmark$ & $\checkmark$ & $\checkmark$ \\
Observations & 10,494 & 10,494 & 10,494 & 10,494 \\
Mean Dependent Variable & 9.002 & 3.279 & 0.158 & 0.150 \\ 
\hline 
\end{tabular}}
\end{subtable}
\caption{\small \textbf{Full estimation results with selected groups of streamers}. All standard errors are clustered at the streamer level. The mean dependent variable is calculated based on observations of all streamers before the policy announcement. FE = fixed effects. Significance level: *$p<0.1$; **$p<0.05$; ***$p<0.01$.}
\label{table:full-estimation-results-isolated-streamers}
\end{table}

\newpage

% Finalized on Sept 22, 2025.
\begin{table}[!htb]
\renewcommand{\arraystretch}{1.2}
\centering
\begin{subtable}{\textwidth}\centering
\caption{Supply-Side Outcomes}
\label{subtable:supply-side-outcomes}
\begin{tabular}{lcccc}
\hline 
            & Gambling                                                    & Loot Box                                                  & Streaming Hours                                             & Other Games                                                 \\ \cline{1-5} 
Banned $\times$ Post ($\beta_1$)     & \begin{tabular}[c]{@{}c@{}}-1.508***\\ (0.123)\end{tabular} & \begin{tabular}[c]{@{}c@{}}0.099\\ (0.197)\end{tabular} & \begin{tabular}[c]{@{}c@{}}-0.677***\\ (0.200)\end{tabular} & \begin{tabular}[c]{@{}c@{}}-0.015\\ (0.195)\end{tabular}   \\
Banned $\times$ Between ($\beta_2$)     & \begin{tabular}[c]{@{}c@{}}-0.078\\ (0.137)\end{tabular}     & \begin{tabular}[c]{@{}c@{}}0.398\\ (0.375)\end{tabular}  & \begin{tabular}[c]{@{}c@{}}0.306\\ (0.228)\end{tabular}    & \begin{tabular}[c]{@{}c@{}}0.071\\ (0.204)\end{tabular}  \\
Unbanned $\times$ Post ($\beta_3$)     & \begin{tabular}[c]{@{}c@{}}-0.139\\ (0.18)\end{tabular}    & \begin{tabular}[c]{@{}c@{}}0.092\\ (0.103)\end{tabular}  & \begin{tabular}[c]{@{}c@{}}-0.227*\\ (0.136)\end{tabular}  & \begin{tabular}[c]{@{}c@{}}-0.128\\ (0.091)\end{tabular}   \\
Unbanned $\times$ Between ($\beta_4$)      & \begin{tabular}[c]{@{}c@{}}0.156\\ (0.133)\end{tabular}     & \begin{tabular}[c]{@{}c@{}}0.102\\ (0.105)\end{tabular}  & \begin{tabular}[c]{@{}c@{}}-0.044\\ (0.09)\end{tabular}    & \begin{tabular}[c]{@{}c@{}}-0.13\\ (0.102)\end{tabular} \\ \hline 
Streamer FE & $\checkmark$                                                         & $\checkmark$                                                      & $\checkmark$                                                         & $\checkmark$                                                         \\
Week FE     & $\checkmark$                                                         & $\checkmark$                                                      & $\checkmark$                                                         & $\checkmark$                                                         \\
Observations & 946 & 946 & 946 & 946    \\
Mean Dependent Variable & 1.470 & 1.210 & 2.700 & 0.437\\ \hline 
\end{tabular}
\end{subtable}
\vspace{0.5cm}
\begin{subtable}{\textwidth}\centering
\caption{Demand-Side Outcomes}
\label{subtable:demand-side-outcomes}
\resizebox{\textwidth}{!}{%
\begin{tabular}{lcccc}
\hline 
            & \multicolumn{1}{c}{log(Hours Watched+1)} 
            & \multicolumn{1}{c}{log(Tier 1+1)} 
            & \multicolumn{1}{c}{log(Tier 2+1)} 
            & \multicolumn{1}{c}{log(Tier 3+1)} \\ \hline
Banned $\times$ Post ($\beta_1$)     
    & \begin{tabular}[c]{@{}c@{}}-1.935***\\ (0.867)\end{tabular} 
    & \begin{tabular}[c]{@{}c@{}}-0.754***\\ (0.298)\end{tabular} 
    & \begin{tabular}[c]{@{}c@{}}-0.004\\ (0.026)\end{tabular} 
    & \begin{tabular}[c]{@{}c@{}}-0.059\\ (0.064)\end{tabular} \\
Banned $\times$ Between ($\beta_2$)     
    & \begin{tabular}[c]{@{}c@{}}0.560\\ (0.503)\end{tabular} 
    & \begin{tabular}[c]{@{}c@{}}0.155\\ (0.188)\end{tabular} 
    & \begin{tabular}[c]{@{}c@{}}-0.006\\ (0.027)\end{tabular} 
    & \begin{tabular}[c]{@{}c@{}}-0.015\\ (0.021)\end{tabular} \\
Unbanned $\times$ Post ($\beta_3$)     
    & \begin{tabular}[c]{@{}c@{}}-0.961*\\ (0.535)\end{tabular} 
    & \begin{tabular}[c]{@{}c@{}}-0.485*\\ (0.282)\end{tabular} 
    & \begin{tabular}[c]{@{}c@{}}0.019\\ (0.057)\end{tabular} 
    & \begin{tabular}[c]{@{}c@{}}0.045*\\ (0.027)\end{tabular} \\
Unbanned $\times$ Between ($\beta_4$)     
    & \begin{tabular}[c]{@{}c@{}}-0.670*\\ (0.343)\end{tabular} 
    & \begin{tabular}[c]{@{}c@{}}-0.689**\\ (0.273)\end{tabular} 
    & \begin{tabular}[c]{@{}c@{}}0.060\\ (0.059)\end{tabular} 
    & \begin{tabular}[c]{@{}c@{}}0.048\\ (0.042)\end{tabular} \\ \hline
Streamer FE & $\checkmark$ & $\checkmark$ & $\checkmark$ & $\checkmark$ \\
Week FE     & $\checkmark$ & $\checkmark$ & $\checkmark$ & $\checkmark$ \\
Observations & 946 & 946 & 946 & 946 \\
Mean Dependent Variable & 9.002 & 3.279 & 0.158 & 0.150 \\ 
\hline 
\end{tabular}}
\end{subtable}
\caption{\small \textbf{Full estimation results with selected streamer communities}. All standard errors are clustered at the community level. The mean dependent variable is calculated based on observations of all selected communities before the policy announcement. FE = fixed effects. Significance level: *$p<0.1$; **$p<0.05$; ***$p<0.01$.}
\label{table:full-estimation-results-isolated-streamers-community}
\end{table}

\clearpage
\newpage

\section{Network Analysis and Community Selection}\label{appendix_network}
As explained in Section \ref{subsection-proposed-solution}, we use the Louvain method \citep{blondel2008fast} to identify 55 communities from the network as listed in Table \ref{table:detected-communities}. For each treatment group, we want to select a subset of streamers such that there is no significant across-group interference. However, it is not straightforward to select streamers for the banned and unbanned groups since most of them are within communities that are mixed in nature in terms of treatment status. These are the communities which we need the automated approach as described in Section \ref{subsection-proposed-solution} to select subgroups of streamers of the same treatment status from. 

To determine which mixed-type communities to use for each treatment group, we use a set of simple criteria. First, among these mixed-type communities, we exclude very small communities such as community 32 and 51 (i.e., communities that do not have at least 10 streamers) from consideration, since they lack the size necessary to extract meaningful subgroups of homogeneous streamers from. We also exclude communities where streamers of different treatment status are highly interwoven such as community 8 and 16 (i.e., variance of the shares of banned, unbanned and untreated less than a threshold of 0.05), as it is operationally difficult to disentangle streamers of the same status from others to ensure low interference. 

Then, from the remaining mixed-type communities, we prioritize the selection of banned and unbanned streamers over untreated streamers when it comes to determine streamers of which treatment status to select from each community. More specifically, since we can only select streamers of one treatment status from each community to avoid interference across treatment status, we prioritize the selection of banned streamers from a community (and thus determine the starting node as a banned streamer and iteratively select banned streamers using the algorithm proposed in Section \ref{subsection-proposed-solution}) if the community consists of a substantial share of banned streamers (at least 20\%). After addressing communities with a large share of banned streamers, we apply the same logic to select unbanned streamers in the remaining communities. Finally, the rest of mixed-type communities are assigned for untreated streamers. 

Table \ref{table:used-communities} shows the final corresponding communities used for each treatment status. Since there are naturally fewer banned and unbanned streamers compared to untreated in our sample, this selection approach is necessary to ensure that our final subgroups contain enough banned and unbanned streamers to main a good representation of treatment statues as in our original sample. 



% While \textit{homogeneous communities} are not subject to SUTVA violation, we begin by identifying \textit{dominantly homogeneous or highly mixed communities} where there could be potential across-group interference. These are communities which we need an automated approach to select subgroups of streamers of the same treatment status from. For these mixed-treatment communities, we use a set of criteria to retain only those that are operationally easy and meaningful to disentangle streamers of 
% one treatment status from others. To ensure that the subgroups contain enough banned and unbanned streamers, we prioritize the selection of banned and then unbanned streamers first from these mixed-treatment communities. In Appendix \ref{community-selection}, we provide detailed discussions and display the communities used for each treatment status. 



% For any given community $j$, the algorithm works by tranversing the network $G_j$ from the start node $s_{c,j}$ in an interative manner. A queue $Q$ is used to keep track of the nodes that are encountered which we need to explore next in the iterative implementation of BFS. Beginning from the start node, we check whether a node has been explored before enqueueing it. We dequeue a node as we explore and process it (in our case, remove if different from target treatment status) and then enqueueing all of its unvisited neighboring nodes next. Essentially, the queue ensures that nodes closer to the start node will be visited earlier than nodes that are further away. The algorithm continues layer by layer until the queue is empty, i.e, all nodes have been processed, rendering a final subset of streamers of homogeneous treatment status. Algorithm \ref{algorithm-1} outlines the pseudo-code of the procedure.   
\clearpage

\begin{tabularx}{}{@{}lllll@{}}
\hline
\makecell[l]{Community ID} & \makecell[l]{Community Size} & \makecell[l]{Untreated \\ Share (\%)} & \makecell[l]{Unbanned \\ Share (\%)} & \makecell[l]{Banned \\Share (\%)} \\ \hline
\endfirsthead
\hline
\makecell[l]{Community ID} & \makecell[l]{Community Size} & \makecell[l]{Untreated \\ Share (\%)} & \makecell[l]{Unbanned \\ Share (\%)} & \makecell[l]{Banned \\Share (\%)}  \\
\hline
\endhead
\hline
\multicolumn{5}{r}{} \\
\endfoot
\hline
\caption{\textbf{Community clusters and treatment status}. This table presents the 55 communities detected by the Louvain Community Detection algorithm. For each community cluster, we show the number of streamers within the cluster and the share of streamers of different treatment status. Note that some clusters consist of streamers of one treatment status, while others are mixed. We discuss the streamers and communities used in our constructed treated and untreated groups for SUTVA alleviation in Section \ref{network-analysis}.}
\label{table:detected-communities}
\endlastfoot
0 & 3 & 0.00 & 66.67 & 33.33 \\
1 & 204 & 74.51 & 9.31 & 16.18 \\
2 & 7 & 100.00 & 0.00 & 0.00 \\
3 & 4 & 25.00 & 75.00 & 0.00 \\
4 & 3 & 0.00 & 100.00 & 0.00 \\
5 & 2 & 50.00 & 50.00 & 0.00 \\
6 & 39 & 89.74 & 2.56 & 7.69 \\
7 & 3 & 100.00 & 0.00 & 0.00 \\
8 & 118 & 33.05 & 38.14 & 28.81 \\
9 & 2 & 50.00 & 50.00 & 0.00 \\
10 & 42 & 73.81 & 0.00 & 26.19 \\
11 & 4 & 100.00 & 0.00 & 0.00 \\
12 & 8 & 25.00 & 75.00 & 0.00 \\
13 & 75 & 8.00 & 42.67 & 49.33 \\
14 & 46 & 100.00 & 0.00 & 0.00 \\
15 & 13 & 100.00 & 0.00 & 0.00 \\
16 & 39 & 33.33 & 43.59 & 23.08 \\
17 & 3 & 100.00 & 0.00 & 0.00 \\
18 & 116 & 25.86 & 63.79 & 10.34 \\
19 & 17 & 100.00 & 0.00 & 0.00 \\
20 & 2 & 100.00 & 0.00 & 0.00 \\
21 & 20 & 100.00 & 0.00 & 0.00 \\
22 & 5 & 100.00 & 0.00 & 0.00 \\
23 & 2 & 100.00 & 0.00 & 0.00 \\
24 & 2 & 50.00 & 0.00 & 50.00 \\
25 & 3 & 100.00 & 0.00 & 0.00 \\
26 & 16 & 93.75 & 0.00 & 6.25 \\
27 & 2 & 0.00 & 50.00 & 50.00 \\
28 & 2 & 100.00 & 0.00 & 0.00 \\
29 & 19 & 100.00 & 0.00 & 0.00 \\
30 & 2 & 0.00 & 0.00 & 100.00 \\
31 & 2 & 100.00 & 0.00 & 0.00 \\
32 & 4 & 25.00 & 50.00 & 25.00 \\
33 & 2 & 100.00 & 0.00 & 0.00 \\
34 & 4 & 100.00 & 0.00 & 0.00 \\
35 & 2 & 100.00 & 0.00 & 0.00 \\
36 & 2 & 100.00 & 0.00 & 0.00 \\
37 & 2 & 100.00 & 0.00 & 0.00 \\
38 & 2 & 100.00 & 0.00 & 0.00 \\
39 & 2 & 100.00 & 0.00 & 0.00 \\
40 & 2 & 100.00 & 0.00 & 0.00 \\
41 & 7 & 100.00 & 0.00 & 0.00 \\
42 & 3 & 100.00 & 0.00 & 0.00 \\
43 & 2 & 100.00 & 0.00 & 0.00 \\
44 & 2 & 100.00 & 0.00 & 0.00 \\
45 & 2 & 100.00 & 0.00 & 0.00 \\
46 & 3 & 100.00 & 0.00 & 0.00 \\
47 & 2 & 50.00 & 50.00 & 0.00 \\
48 & 2 & 100.00 & 0.00 & 0.00 \\
49 & 19 & 42.11 & 57.89 & 0.00 \\
50 & 2 & 0.00 & 100.00 & 0.00 \\
51 & 4 & 50.00 & 50.00 & 0.00 \\
52 & 2 & 0.00 & 100.00 & 0.00 \\
53 & 2 & 0.00 & 100.00 & 0.00 \\
54 & 2 & 50.00 & 50.00 & 0.00 \\
\end{tabularx}




% \begin{table}[!htb]
% \small
% \centering
% \renewcommand{\arraystretch}{1}
% \begin{threeparttable}
% \begin{tabular}{llllll}
% \hline 
% Community ID & Community Size & Untreated Share (\%) & Unbanned Share (\%) & Banned Share (\%) \\
% \hline
% 0 & 3 & 0.00 & 66.67 & 33.33 \\
% 1 & 6 & 100.00 & 0.00 & 0.00 \\
% 2 & 2 & 0.00 & 100.00 & 0.00 \\
% 3 & 209 & 75.12 & 9.09 & 15.79 \\
% 4 & 2 & 0.00 & 100.00 & 0.00 \\
% 5 & 2 & 50.00 & 50.00 & 0.00 \\
% 6 & 4 & 25.00 & 75.00 & 0.00 \\
% 7 & 3 & 0.00 & 100.00 & 0.00 \\
% 8 & 2 & 50.00 & 50.00 & 0.00 \\
% 9 & 2 & 50.00 & 50.00 & 0.00 \\
% 10 & 122 & 33.61 & 38.52 & 27.87 \\
% 11 & 4 & 100.00 & 0.00 & 0.00 \\
% 12 & 46 & 100.00 & 0.00 & 0.00 \\
% 13 & 13 & 100.00 & 0.00 & 0.00 \\
% 14 & 8 & 25.00 & 75.00 & 0.00 \\
% 15 & 77 & 7.79 & 42.86 & 49.35 \\
% 16 & 3 & 100.00 & 0.00 & 0.00 \\
% 17 & 39 & 33.33 & 43.59 & 23.08 \\
% 18 & 17 & 100.00 & 0.00 & 0.00 \\
% 19 & 115 & 25.22 & 64.35 & 10.43 \\
% 20 & 42 & 73.81 & 0.00 & 26.19 \\
% 21 & 2 & 100.00 & 0.00 & 0.00 \\
% 22 & 20 & 100.00 & 0.00 & 0.00 \\
% 23 & 5 & 100.00 & 0.00 & 0.00 \\
% 24 & 2 & 100.00 & 0.00 & 0.00 \\
% 25 & 2 & 50.00 & 0.00 & 50.00 \\
% 26 & 3 & 100.00 & 0.00 & 0.00 \\
% 27 & 16 & 93.75 & 0.00 & 6.25 \\
% 28 & 2 & 100.00 & 0.00 & 0.00 \\
% 29 & 39 & 89.74 & 2.56 & 7.69 \\
% 30 & 19 & 100.00 & 0.00 & 0.00 \\
% 31 & 2 & 0.00 & 0.00 & 100.00 \\
% 32 & 2 & 100.00 & 0.00 & 0.00 \\
% 33 & 4 & 25.00 & 50.00 & 25.00 \\
% 34 & 2 & 100.00 & 0.00 & 0.00 \\
% 35 & 4 & 100.00 & 0.00 & 0.00 \\
% 36 & 2 & 100.00 & 0.00 & 0.00 \\
% 37 & 2 & 100.00 & 0.00 & 0.00 \\
% 38 & 2 & 100.00 & 0.00 & 0.00 \\
% 39 & 2 & 100.00 & 0.00 & 0.00 \\
% 40 & 2 & 100.00 & 0.00 & 0.00 \\
% \hline 
% \end{tabular}
% \end{threeparttable}
% \caption{\small \textbf{Community clusters and treatment status (Part 1)}.}
% \label{table:detected-communities}
% \end{table}

% \begin{table}[!htb]
% \small
% \centering
% \renewcommand{\arraystretch}{1}
% \begin{threeparttable}
% \begin{tabular}{llllll}
% \hline 
% Community ID & Community Size & Untreated Share (\%) & Unbanned Share (\%) & Banned Share (\%) \\
% \hline
% 41 & 2 & 100.00 & 0.00 & 0.00 \\
% 42 & 7 & 100.00 & 0.00 & 0.00 \\
% 43 & 3 & 100.00 & 0.00 & 0.00 \\
% 44 & 3 & 100.00 & 0.00 & 0.00 \\
% 45 & 2 & 100.00 & 0.00 & 0.00 \\
% 46 & 2 & 100.00 & 0.00 & 0.00 \\
% 47 & 2 & 100.00 & 0.00 & 0.00 \\
% 48 & 2 & 50.00 & 50.00 & 0.00 \\
% 49 & 2 & 100.00 & 0.00 & 0.00 \\
% 50 & 19 & 42.11 & 57.89 & 0.00 \\
% 51 & 2 & 0.00 & 100.00 & 0.00 \\
% \hline 
% \end{tabular}
% \end{threeparttable}
% \caption{\small \textbf{Community clusters and treatment status (Part 2)}. This table presents the 52 communities detected by the Louvain Community Detection algorithm. For each community cluster, we show the number of streamers within the cluster and the share of streamers of different treatment status. Note that some clusters consist of streamers of one treatment status, while others are mixed. We discuss the streamers and communities used in our constructed treated and untreated groups for SUTVA alleviation in Section \ref{network-analysis}.}
% \label{table:detected-communities-part2}
% \end{table}


\begin{table}[!htb]
\centering
\renewcommand{\arraystretch}{1.2}
\begin{threeparttable}

% Panel A: Communities with solely unbanned or banned status
\begin{tabular}{lccccc}
\hline 
\multicolumn{5}{c}{Panel A: Communities for Banned Group} \\
\hline
Community ID & Center & Community Size & Streamers Kept & Banned Share (\%) \\
\midrule
13 & jonvlogs  & 75 & 37 & 49.33 \\ %hudsonamorim1 is also okay 
10 & aker\_ & 42 & 11 & 26.19 \\
Others (1 cluster) & - & 2 & 2  & 100.00 \\
\hline 
\label{table_cluster_panel_a}
\end{tabular}

% Panel B: Communities with solely unbanned or banned status
\begin{tabular}{lccccc}
\hline 
\multicolumn{5}{c}{Panel B: Communities for Unbanned Group} \\
\hline
Community ID & Center & Community Size & Streamers Kept & Unbanned Share (\%) \\
\midrule
18 & zloyn & 116 & 73 &  62.93 \\
49 & ilgabbrone & 19 & 10 & 52.63 \\
Others (4 clusters) & - & 9 & 9  & 100.00 \\
\hline 
\end{tabular}

\vspace{0.5cm}
% Panel C: Communities with solely untreated status
\begin{tabular}{lccccc}
\hline 
\multicolumn{5}{c}{Panel C: Communities for Untreated Group} \\ \hline
Community ID & Center & Community Size & Streamers Kept & Untreated Share (\%) \\
\midrule
1 & nyanners & 204 & 98 & 48.04 \\
6 & trymacs & 39 & 35 & 89.74\\
26 & meduska & 16 & 15 & 93.75 \\
Others (30 clusters) & - & 187 & 187  & 100.00 \\
\hline 
\end{tabular}
\end{threeparttable}
\caption{\textbf{Community clusters by treatment status}. This table presents the communities clusters used for each group. Panel A shows clusters used for banned streamers, Panel B shows clusters used for unbanend streamers, and Panel C shows clusters used for untreated streamers. The ``others'' communities in the column are communities that contain only a treatment status of only one type. Specifically, communities 2,  7, 11, 14, 15, 17, 19, 20, 21, 22, 23, 25, 28, 29, 31, 33, 34, 35, 36, 37, 38, 39, 40, 41, 42, 43, 44, 45, 46, 48 are solely of untreated streamers; 4,50,52,53 are solely unbanned streamers; 30 is solely banned streamers.}
\label{table:used-communities}
\end{table}





% \begin{table}[!htb]
% \centering
% \renewcommand{\arraystretch}{1.2}
% \begin{threeparttable}

% % Panel A: Communities with solely unbanned or banned status
% \begin{tabular}{lccccc}
% \hline 
% \multicolumn{5}{c}{Panel A: Communities for Banned Group} \\
% \hline
% Community ID & Center & Community Size & Streamers Kept & Banned Share (\%) \\
% \midrule
% 15 & jonvlogs & 77 & 38 & 49.35 \\
% 20 & aker\_ & 42 & 11 & 26.19 \\
% Others (1 cluster) & - & 2 & 2  & 100.00 \\
% \hline 
% \label{table_cluster_panel_a}
% \end{tabular}

% % Panel B: Communities with solely unbanned or banned status
% \begin{tabular}{lccccc}
% \hline 
% \multicolumn{5}{c}{Panel B: Communities for Unbanned Group} \\
% \hline
% Community ID & Center & Community Size & Streamers Kept & Unbanned Share (\%) \\
% \midrule
% 19 & zloyn & 115 & 82 & 64.35 \\
% Others (4 clusters) & - & 9 & 9  & 100.00 \\
% \hline 
% \end{tabular}

% \vspace{0.5cm}
% % Panel C: Communities with solely untreated status
% \begin{tabular}{lccccc}
% \hline 
% \multicolumn{5}{c}{Panel C: Communities for Untreated Group} \\ \hline
% Community ID & Center & Community Size & Streamers Kept & Untreated Share (\%) \\
% \midrule
% 3 & nyanners & 209 & 98 & 75.12 \\
% Others (29 clusters) & - & 183 & 183  & 100.00 \\
% \hline 
% \end{tabular}
% \end{threeparttable}
% \caption{\small \textbf{Community clusters by treatment status}. This table presents the communities clusters used for each group. Panel A shows clusters used for banned streamers, Panel B shows clusters used for unbanend streamers, and Panel C shows clusters used for untreated streamers. The ``others'' communities in the column are communities that contain only a treatment status of only one type. Specifically, communities 1, 11,12,13, 16, 18, 21, 22, 23, 24, 26, 28, 30, 32, 34, 35, 36, 37, 38, 39, 40, 41, 42, 43, 44, 45, 46, 47, 49 are solely of untreated streamers; 2,4,7, 51 are solely unbanned streamers; 31 is solely banned streamers.}
% \label{table:used-communities}
% \end{table}



\begin{figure}[h!]
  \captionsetup{type=algorithm} % treat this figure as an Algorithm for caption/counters
  \centering
  \begin{minipage}{0.98\linewidth}
  \noindent\rule{\linewidth}{0.6pt}
    %\vspace{0.1\baselineskip}
    \begin{algorithmic}[1]
    \Require A graph $G_j$ and start node (streamer) $s_{center,j}$ of $G_j$ for any community $j$
    % \Ensure $T$
    \State Initialize $V \coloneq $ a set to store all visited streamers
    \State Initialize $T \coloneq $ a set to store all streamers of the same treatment status as $s_{c,j}$
    \State Initialize $Q \coloneq $ a queue to store all streamers encountered but not yet visited 
    \State $Q$.enqueue($s_{center,j}$)
    \State $V$.add($s_{center,j}$)
    \State Set target\_treatment $\coloneq$ treatment status of $s_{center,j}$
    \While{$Q \neq \emptyset$}
        \State current\_streamer $\coloneq$ $Q$.dequeue()
        \If{current\_streamer $\notin$ $V$}
            \State $V$.add(current\_streamer)
            \If{current\_streamer's treatment $==$ target\_treatment}
                \State $T$.add(current\_streamer)
                \For{neighbor $\in$ $G_j$.neighbors(current\_streamer)}
                    \If{neighbor $\notin$ $V$}
                        \State $Q$.enqueue(neighbor)
                    \EndIf
                \EndFor
            \EndIf
        \EndIf
    \EndWhile
    \State \Return A subset of streamers $T$ with homogeneous treatment status
    \end{algorithmic}
  %\vspace{0.1\baselineskip}
  \noindent\rule{\linewidth}{0.6pt}
  \end{minipage}
  \caption{\textbf{Breadth-First Search Based Filtering Algorithm}. We use the algorithm in Section \ref{subsection-proposed-solution} to automate the selection of subgroups of streamers 
  who share the same treatment status but have minimal to no overlap in viewers to reduce SUTVA.}
  \label{algorithm-1}
\end{figure}


\begin{figure}[ht!]
     \centering
     % First row
     \begin{subfigure}[b]{0.24\textwidth}
         \centering
         \includegraphics[width=\textwidth]{Figures/network/community14.png}
         \caption{Community 14}
         \label{fig:token1}
     \end{subfigure}
    \begin{subfigure}[b]{0.24\textwidth}
         \centering
         \includegraphics[width=\textwidth]{Figures/network/community10.png}
         \caption{Community 10}
         \label{fig:token3}
     \end{subfigure}
     \begin{subfigure}[b]{0.24\textwidth}
         \centering
         \includegraphics[width=\textwidth]{Figures/network/community13.png}
         \caption{Community 13}
         \label{fig:token2}
     \end{subfigure}
     \begin{subfigure}[b]{0.24\textwidth}
         \centering
         \includegraphics[width=\textwidth]{Figures/network/community16.png}
         \caption{Community 16}
         \label{fig:token4}
     \end{subfigure}
     \vspace{0.5cm}
     % Second row
     \begin{subfigure}[b]{0.24\textwidth}
         \centering
         \includegraphics[width=\textwidth]{Figures/network/community8.png}
         \caption{Community 8}
         \label{fig:token5}
     \end{subfigure}
     \begin{subfigure}[b]{0.24\textwidth}
         \centering
         \includegraphics[width=\textwidth]{Figures/network/community18.png}
         \caption{Community 18}
         \label{fig:token6}
     \end{subfigure}
     \begin{subfigure}[b]{0.24\textwidth}
         \centering
         \includegraphics[width=\textwidth]{Figures/network/community32.png}
         \caption{Community 32}
         \label{fig:token7}
     \end{subfigure}
     \begin{subfigure}[b]{0.24\textwidth}
         \centering
         \includegraphics[width=\textwidth]{Figures/network/community51.png}
         \caption{Community 51}
         \label{fig:token8}
     \end{subfigure}
    \vspace{0.5cm}
        \caption{\small \textbf{Examples of Communities}. Communities are detected using the Louvain Method \citep{blondel2008fast}. Red denotes untreated, blue unbanned (treated) and green banned (treated) group.}
     \label{local-network}
\end{figure}



\begin{figure}[h!]
    \centering
    \includegraphics[scale=0.8]{Figures/Descriptives/bfs.png} 
    \caption{\textbf{Filtering streamers of different treatment status iteratively using BFS}. For a given community, the algorithm works by traversing the network from the start node iteratively. A queue is used to keep track of the nodes that are encountered which we need to explore next in the iterative implementation of BFS. We begin from the start streamer ``jonvlogs" (green node), search the network outward from him, reaching to his closest streamers which can be of different treatment statuses (blue nodes) (Panel a). We declare his direct neighbors to be at distance 1, and 
    filtering out streamers of different statuses connected to him (Panel b). Next, we find all neighbors of ``jonvlogs'''s direct streamers (not counting streamers who are already direct neighbors of the central streamer) and declare those to be at distance 2 (Panel c). For example, ``pluzinho'' (green node) is one such streamer at distance 1 who is connected to other streamers of different status (blue node). The BFS process then includes all direct neighbors of ``pluzinho'', again filtering out any streamers of different treatment status from him (Panel d). The process continues successively through the network until all streamers are processed.}
    \label{bfs-streamer-filter}
\end{figure}



\begin{table}[h!]
\small
\centering
\renewcommand{\arraystretch}{1.2}
\begin{threeparttable}
\begin{tabular}{llllllll}
\hline 
Group & Mean & Std & 5\textsubscript{quantile} & 25\textsubscript{quantile} & 50\textsubscript{quantile} & 75\textsubscript{quantile} & 95\textsubscript{quantile} \\
\midrule
Untreated & 15,395.00 & 29,541.50 & 1,338.60 & 3,957.00 & 7,933.50 & 16,248.75 & 49,197.95 \\
Unbanned & 18,184.54 & 43,827.04 & 666.10 & 3,792.50 & 8,698.00 & 19,213.50 & 52,443.20 \\
Banned & 37,134.57 & 54,788.90 & 3,782.20 & 10,825.00 & 22,637.00 & 35,796.50 & 124,309.20 \\
\hline 
\end{tabular}
\end{threeparttable}
\caption{\textbf{Viewer count before the policy announcement}.}
\label{table:network-group-statistics}
\end{table}








\newpage
\clearpage
\section{Additional Figures and Tables}\label{appendix_figures}


\begin{figure}[h!]
    \centering
    \begin{subfigure}[b]{0.9\textwidth}
        \centering
        \includegraphics[width=\linewidth]{Figures/Appendix/PEGI-rating.png}
        \caption{PEGI}
        \label{fig:pegi}
    \end{subfigure}

    \medskip % vertical space between the two
    \begin{subfigure}[b]{0.8\textwidth}
        \centering
        \includegraphics[width=\linewidth]{Figures/Appendix/ESRB-rating.jpg}
        \caption{ESRB}
        \label{fig:esrb}
    \end{subfigure}
    \caption{\textbf{In-game purchases labels from PEGI and ESRB.} The figures show the labels applied to games that include loot boxes or randomized in-game purchases. In PEGI, games are labeled with ``Includes Paid Random Items'' (top); source: \url{https://pegi.info/what-do-the-labels-mean}. In ESRB, games are labeled with ``In-Game Purchases (Includes Random Items'' (bottom); source: \url{https://www.esrb.org/blog/in-game-purchases-includes-random-items/} }
    \label{lootbox_ratings}
\end{figure}

\begin{figure}[H]
    \centering
    \begin{subfigure}[b]{0.45\textwidth}
    \centering
    \includegraphics[width=\textwidth]{Figures/Others/semrush_weekly_effect_on_Paid Social United States Desktop.jpg}
    \caption{U.S., Desktop}
    \label{fig:visit-us-desktop}
    \end{subfigure}
    \begin{subfigure}[b]{0.45\textwidth}
    \centering
    \includegraphics[width=\textwidth]{Figures/Others/semrush_weekly_effect_on_Paid Social United States All Devices.jpg}
    \caption{U.S., All Devices}
    \label{fig:visit-us-all}
    \end{subfigure}
    \begin{subfigure}[b]{0.45\textwidth}
    \centering
    \includegraphics[width=\textwidth]{Figures/Others/semrush_weekly_effect_on_Paid Social Worldwide Desktop.jpg}
    \caption{Worldwide, Desktop}
    \label{fig:visit-ww-desktop}
    \end{subfigure}
    \begin{subfigure}[b]{0.45\textwidth}
    \centering
    \includegraphics[width=\textwidth]{Figures/Others/semrush_weekly_effect_on_Paid Social Worldwide All Devices.jpg}
    \caption{Worldwide, All Devices}
    \label{fig:visit-ww-all}
    \end{subfigure}
    \caption{\textbf{Time-varying policy effect on paid social media search}. 1. This figure presents the time-varying treatment effects on total paid search to banned gambling websites, using unbanned gambling websites described in Table \ref{table:traffic-websites} as the control group. 2. We include linear time trends for each website and cluster the standard errors at the website level.}
    \label{es_traffic_paid}
\end{figure}

\begin{figure}[H]
    \centering
    \begin{subfigure}[b]{0.45\textwidth}
    \centering
    \includegraphics[width=\textwidth]{Figures/Others/semrush_weekly_effect_on_Organic Social United States Desktop.jpg}
    \caption{U.S., Desktop}
    \label{fig:visit-us-desktop}
    \end{subfigure}
    \begin{subfigure}[b]{0.45\textwidth}
    \centering
    \includegraphics[width=\textwidth]{Figures/Others/semrush_weekly_effect_on_Organic Social United States All Devices.jpg}
    \caption{U.S., All Devices}
    \label{fig:visit-us-all}
    \end{subfigure}
    \begin{subfigure}[b]{0.45\textwidth}
    \centering
    \includegraphics[width=\textwidth]{Figures/Others/semrush_weekly_effect_on_Organic Social Worldwide Desktop.jpg}
    \caption{Worldwide, Desktop}
    \label{fig:visit-ww-desktop}
    \end{subfigure}
    \begin{subfigure}[b]{0.45\textwidth}
    \centering
    \includegraphics[width=\textwidth]{Figures/Others/semrush_weekly_effect_on_Organic Social Worldwide All Devices.jpg}
    \caption{Worldwide, All Devices}
    \label{fig:visit-ww-all}
    \end{subfigure}
    \caption{\textbf{Time-varying policy effect on organic social media search}. 1. This figure presents the time-varying treatment effects on total organic search to banned gambling websites, using unbanned gambling websites described in Table \ref{table:traffic-websites} as the control group. 2. We include linear time trends for each website and cluster the standard errors at the website level.}
    \label{es_traffic_organic}
\end{figure}

% ===============================================================


% =========================
% Table 1: All Streamers
% =========================

\begin{table}[!htb]
\renewcommand{\arraystretch}{1.05}
\centering
{\fontsize{4}{4.6}\selectfont

\begin{subtable}{\textwidth}
\caption{Number of Views}
\centering
\resizebox{\textwidth}{!}{%
\begin{tabular}{lccccc}
\hline
Age Group & Overall & Gambling & Loot Box & Non-gaming & Other Games \\ \hline
Adult     & \makecell{1422.123\\(3441.849)} & \makecell{58.069\\(594.511)}  & \makecell{627.242\\(1700.795)} & \makecell{381.167\\(1584.525)} & \makecell{355.645\\(1318.914)} \\
Minor     & \makecell{552.054\\(5199.240)}  & \makecell{5.842\\(143.062)}   & \makecell{291.774\\(3225.082)} & \makecell{179.906\\(1802.774)} & \makecell{74.533\\(608.182)}   \\
Uncertain & \makecell{2704.268\\(6117.913)} & \makecell{94.653\\(1057.383)} & \makecell{1205.476\\(3163.276)}& \makecell{757.785\\(3191.560)}  & \makecell{646.354\\(1900.339)} \\ \hline
\end{tabular}
}
\end{subtable}

\vspace{0.1cm}

\begin{subtable}{\textwidth}
\caption{Number of Chats per View}
\centering
\resizebox{\textwidth}{!}{%
\begin{tabular}{lccccc}
\hline
Age Group & Overall & Gambling & Loot Box & Non-gaming & Other Games \\ \hline
Adult     & \makecell{14.858\\(24.499)} & \makecell{0.842\\(5.175)}  & \makecell{10.753\\(33.062)} & \makecell{6.822\\(16.147)} & \makecell{7.290\\(12.520)} \\
Minor     & \makecell{7.035\\(9.162)}   & \makecell{0.297\\(2.949)}  & \makecell{4.033\\(7.379)}   & \makecell{3.659\\(8.785)}  & \makecell{3.714\\(7.440)}  \\
Uncertain & \makecell{6.459\\(4.378)}   & \makecell{0.374\\(2.091)}  & \makecell{3.774\\(4.052)}   & \makecell{3.275\\(4.873)}  & \makecell{3.612\\(4.824)}  \\ \hline
\end{tabular}
}
\end{subtable}

\vspace{0.1cm}

\begin{subtable}{\textwidth}
\caption{Average Emotes per Message}
\centering
\resizebox{\textwidth}{!}{%
\begin{tabular}{lccccc}
\hline
Age Group & Overall & Gambling & Loot Box & Non-gaming & Other Games \\ \hline
Adult     & \makecell{0.327\\(0.447)} & \makecell{0.215\\(0.271)} & \makecell{0.310\\(0.460)} & \makecell{0.384\\(0.561)} & \makecell{0.319\\(0.402)} \\
Minor     & \makecell{0.285\\(0.515)} & \makecell{0.209\\(0.402)} & \makecell{0.269\\(0.524)} & \makecell{0.356\\(0.642)} & \makecell{0.279\\(0.561)} \\
Uncertain & \makecell{0.272\\(0.339)} & \makecell{0.194\\(0.279)} & \makecell{0.254\\(0.317)} & \makecell{0.342\\(0.452)} & \makecell{0.260\\(0.299)} \\ \hline
\end{tabular}
}
\end{subtable}

\vspace{0.1cm}

\begin{subtable}{\textwidth}
\caption{Average Character Count per Message}
\centering
\resizebox{\textwidth}{!}{%
\begin{tabular}{lccccc}
\hline
Age Group & Overall & Gambling & Loot Box & Non-gaming & Other Games \\ \hline
Adult     & \makecell{36.153\\(20.593)} & \makecell{31.267\\(14.730)} & \makecell{36.607\\(22.820)} & \makecell{33.841\\(21.429)} & \makecell{35.877\\(20.817)} \\
Minor     & \makecell{23.990\\(14.629)} & \makecell{22.257\\(20.208)} & \makecell{23.160\\(15.012)} & \makecell{23.394\\(15.583)} & \makecell{24.151\\(14.967)} \\
Uncertain & \makecell{26.578\\(10.122)} & \makecell{23.650\\(8.974)}  & \makecell{25.597\\(10.370)} & \makecell{25.810\\(12.098)} & \makecell{26.514\\(10.410)} \\ \hline
\end{tabular}
}
\end{subtable}

\vspace{0.1cm}

\begin{subtable}{\textwidth}
\caption{Average Token Count per Message}
\centering
\resizebox{\textwidth}{!}{%
\begin{tabular}{lccccc}
\hline
Age Group & Overall & Gambling & Loot Box & Non-gaming & Other Games \\ \hline
Adult     & \makecell{6.090\\(3.381)} & \makecell{5.334\\(2.377)} & \makecell{6.149\\(3.727)} & \makecell{5.565\\(3.410)} & \makecell{6.076\\(3.378)} \\
Minor     & \makecell{4.249\\(2.547)} & \makecell{3.893\\(2.988)} & \makecell{4.118\\(2.649)} & \makecell{3.998\\(2.658)} & \makecell{4.309\\(2.673)} \\
Uncertain & \makecell{4.725\\(1.913)} & \makecell{4.263\\(1.660)} & \makecell{4.562\\(2.045)} & \makecell{4.435\\(2.122)} & \makecell{4.759\\(1.973)} \\ \hline
\end{tabular}
}
\end{subtable}}

\caption{\textbf{Viewing and engagement across age groups and content categories for all sampled streamers}. 
Conditional means by age group with standard deviations in parentheses. 
“Number of views” is the average number of presences of viewers from each age group in all livestreaming sessions of a streamer within a day. 
“Number of chats per view” is the average number of chat messages sent within each presence.}
\label{table:summary-age-category-panels-all}
\end{table}




% =========================
% Table 2: Treated Streamers
% =========================

\begin{table}[!htb]
\renewcommand{\arraystretch}{1.05}
\centering
{\fontsize{4}{4.6}\selectfont

\begin{subtable}{\textwidth}
\caption{Number of Views}
\centering
\resizebox{\textwidth}{!}{%
\begin{tabular}{lccccc}
\hline
Age Group & Overall & Gambling & Loot Box & Non-gaming & Other Games \\ \hline
Adult     & \makecell{2399.607\\(6641.852)} & \makecell{593.662\\(1815.598)} & \makecell{775.845\\(2294.036)} & \makecell{776.054\\(2840.213)} & \makecell{254.046\\(1758.446)} \\
Minor     & \makecell{1422.599\\(13263.742)} & \makecell{66.368\\(478.110)} & \makecell{710.160\\(8284.114)} & \makecell{521.630\\(4419.242)} & \makecell{124.441\\(845.469)} \\
Uncertain & \makecell{4166.734\\(11441.332)} & \makecell{966.821\\(3252.719)} & \makecell{1407.409\\(4489.841)} & \makecell{1364.680\\(4929.623)} & \makecell{427.823\\(3127.616)} \\ \hline
\end{tabular}
}
\end{subtable}

\vspace{0.1cm}

\begin{subtable}{\textwidth}
\caption{Number of Chats per View}
\centering
\resizebox{\textwidth}{!}{%
\begin{tabular}{lccccc}
\hline
Age Group & Overall & Gambling & Loot Box & Non-gaming & Other Games \\ \hline
Adult     & \makecell{16.131\\(16.552)} & \makecell{8.605\\(14.388)} & \makecell{11.141\\(23.011)} & \makecell{8.763\\(12.939)} & \makecell{4.813\\(10.718)} \\
Minor     & \makecell{7.579\\(15.250)} & \makecell{3.371\\(9.406)} & \makecell{3.822\\(9.794)} & \makecell{5.034\\(15.266)} & \makecell{2.167\\(6.025)} \\
Uncertain & \makecell{6.390\\(5.387)} & \makecell{3.817\\(5.616)} & \makecell{3.222\\(4.127)} & \makecell{3.665\\(4.914)} & \makecell{1.867\\(3.273)} \\ \hline
\end{tabular}
}
\end{subtable}

\vspace{0.1cm}

\begin{subtable}{\textwidth}
\caption{Average Emotes per Message}
\centering
\resizebox{\textwidth}{!}{%
\begin{tabular}{lccccc}
\hline
Age Group & Overall & Gambling & Loot Box & Non-gaming & Other Games \\ \hline
Adult     & \makecell{0.247\\(0.316)} & \makecell{0.215\\(0.274)} & \makecell{0.246\\(0.344)} & \makecell{0.285\\(0.387)} & \makecell{0.225\\(0.241)} \\
Minor     & \makecell{0.272\\(0.593)} & \makecell{0.205\\(0.399)} & \makecell{0.272\\(0.569)} & \makecell{0.356\\(0.708)} & \makecell{0.294\\(0.518)} \\
Uncertain & \makecell{0.244\\(0.333)} & \makecell{0.193\\(0.279)} & \makecell{0.238\\(0.305)} & \makecell{0.302\\(0.460)} & \makecell{0.234\\(0.303)} \\ \hline
\end{tabular}
}
\end{subtable}

\vspace{0.1cm}

\begin{subtable}{\textwidth}
\caption{Average Character Count per Message}
\centering
\resizebox{\textwidth}{!}{%
\begin{tabular}{lccccc}
\hline
Age Group & Overall & Gambling & Loot Box & Non-gaming & Other Games \\ \hline
Adult     & \makecell{32.709\\(15.461)} & \makecell{31.380\\(14.604)} & \makecell{35.496\\(18.726)} & \makecell{31.121\\(16.550)} & \makecell{33.759\\(16.861)} \\
Minor     & \makecell{23.331\\(17.966)} & \makecell{22.107\\(20.941)} & \makecell{24.132\\(18.315)} & \makecell{23.682\\(15.809)} & \makecell{23.451\\(18.671)} \\
Uncertain & \makecell{25.042\\(9.954)} & \makecell{23.640\\(8.897)} & \makecell{25.860\\(13.123)} & \makecell{25.312\\(13.148)} & \makecell{24.976\\(12.613)} \\ \hline
\end{tabular}
}
\end{subtable}

\vspace{0.1cm}

\begin{subtable}{\textwidth}
\caption{Average Token Count per Message}
\centering
\resizebox{\textwidth}{!}{%
\begin{tabular}{lccccc}
\hline
Age Group & Overall & Gambling & Loot Box & Non-gaming & Other Games \\ \hline
Adult     & \makecell{5.529\\(2.560)} & \makecell{5.339\\(2.334)} & \makecell{5.975\\(3.162)} & \makecell{5.161\\(2.673)} & \makecell{5.725\\(2.910)} \\
Minor     & \makecell{4.019\\(2.953)} & \makecell{3.819\\(2.947)} & \makecell{4.234\\(3.830)} & \makecell{4.006\\(2.809)} & \makecell{4.107\\(3.755)} \\
Uncertain & \makecell{4.479\\(2.238)} & \makecell{4.255\\(1.642)} & \makecell{4.666\\(3.108)} & \makecell{4.413\\(2.691)} & \makecell{4.489\\(2.994)} \\ \hline
\end{tabular}
}
\end{subtable}}

\caption{\textbf{Viewing and engagement across age groups and content categories for treated streamers.} 
Conditional means by age group with standard deviations in parentheses. 
“Number of views” is the average number of presences of viewers from each age group in all livestreaming sessions of a streamer within a day. 
“Number of chats per view” is the average number of chat messages sent within each presence.}
\label{table:summary-age-category-panels-treated}
\end{table}



% ========================================================
% Additional result tables for viewer analysis.

\begin{table}[!htb]
\renewcommand{\arraystretch}{1.2}
\centering
\resizebox{\textwidth}{!}{%
\begin{threeparttable}
\begin{tabular}{lcccccc}
\toprule
 & \multicolumn{3}{c}{log(Number of Chats per Gambling View + 1)} & \multicolumn{3}{c}{log(Number of Chats per View + 1)} \\
\cmidrule(lr){2-4} \cmidrule(lr){5-7}
 & Minor & Adult & Uncertain & Minor & Adult & Uncertain \\
\midrule
Banned $\times$ Post ($\beta_1$) & 
\begin{tabular}[c]{@{}c@{}}-0.792***\\ (0.072)\end{tabular} & 
\begin{tabular}[c]{@{}c@{}}-1.102***\\ (0.093)\end{tabular} & 
\begin{tabular}[c]{@{}c@{}}-0.782***\\ (0.063)\end{tabular} & 
\begin{tabular}[c]{@{}c@{}}-0.065*\\ (0.039)\end{tabular} & 
\begin{tabular}[c]{@{}c@{}}-0.032\\ (0.025)\end{tabular} & 
\begin{tabular}[c]{@{}c@{}}-0.013\\ (0.024)\end{tabular} \\
$\Delta \%$ & -54.7\% & -66.8\% & -54.3\% & -6.3\% & - & - \\
\addlinespace[0.3em]
Unbanned $\times$ Post ($\beta_3$) & 
\begin{tabular}[c]{@{}c@{}}-0.144***\\ (0.041)\end{tabular} & 
\begin{tabular}[c]{@{}c@{}}-0.116**\\ (0.059)\end{tabular} & 
\begin{tabular}[c]{@{}c@{}}-0.076*\\ (0.045)\end{tabular} & 
\begin{tabular}[c]{@{}c@{}}-0.022\\ (0.023)\end{tabular} & 
\begin{tabular}[c]{@{}c@{}}0.004\\ (0.018)\end{tabular} & 
\begin{tabular}[c]{@{}c@{}}0.004\\ (0.014)\end{tabular} \\
$\Delta \%$ & -13.4\% & -11.0\% & -7.3\% & - & - & - \\
\midrule
Streamer FE & $\checkmark$ & $\checkmark$ & $\checkmark$ & $\checkmark$ & $\checkmark$ & $\checkmark$ \\
Week FE & $\checkmark$ & $\checkmark$ & $\checkmark$ & $\checkmark$ & $\checkmark$ & $\checkmark$ \\
Observations & 80,762 & 92,753 & 92,861 & 80,762 & 92,753 & 92,861 \\
Mean Dependent Variable & 0.842 & 1.522 & 1.113 & 1.821 & 2.641 & 1.882 \\
\bottomrule
\end{tabular}
\end{threeparttable}}
\caption{\small \textbf{The effect of Twitch's banning policy on viewers' log number of chats per view by age group.} 1. Number of chats per view is defined as the average number of chat messages sent in each viewer presence from each age group within each streamer-day. The log variables are transformed as $\log(1+y)$. 2. All standard errors are clustered at the streamer level. The mean dependent variable is calculated based on observations of all treated (banned and unbanned) streamers before the policy announcement. FE = fixed effects. Significance levels: *$p<0.1$; **$p<0.05$; ***$p<0.01$.}
\label{ATT_log_chats_ps}
\end{table}


\begin{table}[!htb]
\renewcommand{\arraystretch}{1.2}
\centering
\resizebox{\textwidth}{!}{%
\begin{threeparttable}
\begin{tabular}{lcccccc}
\toprule
 & \multicolumn{3}{c}{Average Emotes per Message (Gambling)} & \multicolumn{3}{c}{Average Emotes per Message (Overall)} \\
\cmidrule(lr){2-4} \cmidrule(lr){5-7}
 & Minor & Adult & Uncertain & Minor & Adult & Uncertain \\
\midrule
Banned $\times$ Post ($\beta_1$) & 
\begin{tabular}[c]{@{}c@{}}-0.107***\\ (0.022)\end{tabular} & 
\begin{tabular}[c]{@{}c@{}}-0.134***\\ (0.023)\end{tabular} & 
\begin{tabular}[c]{@{}c@{}}-0.113***\\ (0.014)\end{tabular} & 
\begin{tabular}[c]{@{}c@{}}-0.025\\ (0.016)\end{tabular} & 
\begin{tabular}[c]{@{}c@{}}-0.023**\\ (0.011)\end{tabular} & 
\begin{tabular}[c]{@{}c@{}}-0.018\\ (0.012)\end{tabular} \\
\addlinespace[0.3em]
Unbanned $\times$ Post ($\beta_3$) & 
\begin{tabular}[c]{@{}c@{}}-0.031***\\ (0.012)\end{tabular} & 
\begin{tabular}[c]{@{}c@{}}-0.017**\\ (0.008)\end{tabular} & 
\begin{tabular}[c]{@{}c@{}}-0.021**\\ (0.008)\end{tabular} & 
\begin{tabular}[c]{@{}c@{}}-0.015\\ (0.023)\end{tabular} & 
\begin{tabular}[c]{@{}c@{}}-0.005\\ (0.008)\end{tabular} & 
\begin{tabular}[c]{@{}c@{}}0.002\\ (0.008)\end{tabular} \\
\midrule
Streamer FE & $\checkmark$ & $\checkmark$ & $\checkmark$ & $\checkmark$ & $\checkmark$ & $\checkmark$ \\
Week FE & $\checkmark$ & $\checkmark$ & $\checkmark$ & $\checkmark$ & $\checkmark$ & $\checkmark$ \\
Observations & 80,762 & 92,753 & 92,861 & 80,762 & 92,753 & 92,861 \\
Mean Dependent Variable & 0.205 & 0.215 & 0.193 & 0.272 & 0.247 & 0.244 \\
\bottomrule
\end{tabular}
\end{threeparttable}}
\caption{\small \textbf{The effect of Twitch's banning policy on viewers' average emotes per chat message by age group.} 1. The dependent variable measures the average number of emotes contained in each chat message per streamer-day and age group. 2. All standard errors are clustered at the streamer level. The mean dependent variable is calculated based on observations of all treated (banned and unbanned) streamers before the policy announcement. FE = fixed effects. Significance levels: *$p<0.1$; **$p<0.05$; ***$p<0.01$.}
\label{ATT_viewer_avg_emotes}
\end{table}

\begin{table}[!htb]
\renewcommand{\arraystretch}{1.2}
\centering
\resizebox{\textwidth}{!}{%
\begin{threeparttable}
\begin{tabular}{lcccccc}
\toprule
 & \multicolumn{3}{c}{Average Character Count per Message (Gambling)} & \multicolumn{3}{c}{Average Character Count per Message (Overall)} \\
\cmidrule(lr){2-4} \cmidrule(lr){5-7}
 & Minor & Adult & Uncertain & Minor & Adult & Uncertain \\
\midrule
Banned $\times$ Post ($\beta_1$) & 
\begin{tabular}[c]{@{}c@{}}-8.256***\\ (1.124)\end{tabular} & 
\begin{tabular}[c]{@{}c@{}}-10.575***\\ (1.551)\end{tabular} & 
\begin{tabular}[c]{@{}c@{}}-10.646***\\ (0.976)\end{tabular} & 
\begin{tabular}[c]{@{}c@{}}0.282\\ (1.082)\end{tabular} & 
\begin{tabular}[c]{@{}c@{}}0.761\\ (0.860)\end{tabular} & 
\begin{tabular}[c]{@{}c@{}}-0.737\\ (0.557)\end{tabular} \\
\addlinespace[0.3em]
Unbanned $\times$ Post ($\beta_3$) & 
\begin{tabular}[c]{@{}c@{}}-0.935\\ (0.739)\end{tabular} & 
\begin{tabular}[c]{@{}c@{}}-0.067\\ (0.893)\end{tabular} & 
\begin{tabular}[c]{@{}c@{}}-1.109*\\ (0.577)\end{tabular} & 
\begin{tabular}[c]{@{}c@{}}0.017\\ (0.641)\end{tabular} & 
\begin{tabular}[c]{@{}c@{}}1.015**\\ (0.492)\end{tabular} & 
\begin{tabular}[c]{@{}c@{}}-0.561**\\ (0.243)\end{tabular} \\
\midrule
Streamer FE & $\checkmark$ & $\checkmark$ & $\checkmark$ & $\checkmark$ & $\checkmark$ & $\checkmark$ \\
Week FE & $\checkmark$ & $\checkmark$ & $\checkmark$ & $\checkmark$ & $\checkmark$ & $\checkmark$ \\
Observations & 80,762 & 92,753 & 92,861 & 80,762 & 92,753 & 92,861 \\
Mean Dependent Variable & 22.107 & 31.380 & 23.640 & 23.331 & 32.709 & 25.042 \\
\bottomrule
\end{tabular}
\end{threeparttable}}
\caption{\small \textbf{The effect of Twitch's banning policy on viewers' average number of characters per chat message by age group.} 1. The dependent variable measures the average number of characters per chat message sent by viewers in each age group on a streamer-day basis. 2. All standard errors are clustered at the streamer level. The mean dependent variable is calculated based on observations of all treated (banned and unbanned) streamers before the policy announcement. FE = fixed effects. Significance levels: *$p<0.1$; **$p<0.05$; ***$p<0.01$.}
\label{ATT_viewer_char_count}
\end{table}

\begin{table}[!htb]
\renewcommand{\arraystretch}{1.2}
\centering
\resizebox{\textwidth}{!}{%
\begin{threeparttable}
\begin{tabular}{lcccccc}
\toprule
 & \multicolumn{3}{c}{Average Token Count per Message (Gambling)} & \multicolumn{3}{c}{Average Token Count per Message (Overall)} \\
\cmidrule(lr){2-4} \cmidrule(lr){5-7}
 & Minor & Adult & Uncertain & Minor & Adult & Uncertain \\
\midrule
Banned $\times$ Post ($\beta_1$) & 
\begin{tabular}[c]{@{}c@{}}-1.502***\\ (0.212)\end{tabular} & 
\begin{tabular}[c]{@{}c@{}}-1.787***\\ (0.261)\end{tabular} & 
\begin{tabular}[c]{@{}c@{}}-1.958***\\ (0.186)\end{tabular} & 
\begin{tabular}[c]{@{}c@{}}0.135\\ (0.215)\end{tabular} & 
\begin{tabular}[c]{@{}c@{}}0.195\\ (0.156)\end{tabular} & 
\begin{tabular}[c]{@{}c@{}}-0.139\\ (0.130)\end{tabular} \\
\addlinespace[0.3em]
Unbanned $\times$ Post ($\beta_3$) & 
\begin{tabular}[c]{@{}c@{}}-0.157\\ (0.127)\end{tabular} & 
\begin{tabular}[c]{@{}c@{}}-0.009\\ (0.152)\end{tabular} & 
\begin{tabular}[c]{@{}c@{}}-0.197*\\ (0.102)\end{tabular} & 
\begin{tabular}[c]{@{}c@{}}0.011\\ (0.110)\end{tabular} & 
\begin{tabular}[c]{@{}c@{}}0.164*\\ (0.089)\end{tabular} & 
\begin{tabular}[c]{@{}c@{}}-0.118**\\ (0.047)\end{tabular} \\
\midrule
Streamer FE & $\checkmark$ & $\checkmark$ & $\checkmark$ & $\checkmark$ & $\checkmark$ & $\checkmark$ \\
Week FE & $\checkmark$ & $\checkmark$ & $\checkmark$ & $\checkmark$ & $\checkmark$ & $\checkmark$ \\
Observations & 80,762 & 92,753 & 92,861 & 80,762 & 92,753 & 92,861 \\
Mean Dependent Variable & 3.819 & 5.339 & 4.255 & 4.019 & 5.529 & 4.479\\
\bottomrule
\end{tabular}
\end{threeparttable}}
\caption{\small \textbf{The effect of Twitch's banning policy on viewers' average number of tokens per chat message by age group.} 1. The dependent variable measures the average number of tokens per chat message sent by viewers in each age group on a streamer-day basis. 2. All standard errors are clustered at the streamer level. The mean dependent variable is calculated based on observations of all treated (banned and unbanned) streamers before the policy announcement. FE = fixed effects. Significance levels: *$p<0.1$; **$p<0.05$; ***$p<0.01$.}
\label{ATT_viewer_token_count}
\end{table}


\clearpage
% Event study plots.

\begin{figure}[h!]
     \centering
     \begin{subfigure}[b]{0.4\textwidth}
         \centering
         \includegraphics[width=\textwidth]{Figures/Minors/EventStudy_log_num_watched_sessions_gambling_banslots.png}
         \caption{Gambling, Banned}
         \label{fig:token1}
     \end{subfigure}
     \begin{subfigure}[b]{0.4\textwidth}
         \centering
         \includegraphics[width=\textwidth]{Figures/Minors/EventStudy_log_num_watched_sessions_gambling_unbanslots.png}
         \caption{Gambling, Unbanned}
         \label{fig:token2}
     \end{subfigure}
     \vspace{-0.3cm}
     \begin{subfigure}[b]{0.4\textwidth}
         \centering
         \includegraphics[width=\textwidth]{Figures/Minors/EventStudy_log_num_watched_sessions_banslots.png}
         \caption{Total Livestreams, Banned}
         \label{fig:token3}
     \end{subfigure}
     \begin{subfigure}[b]{0.4\textwidth}
         \centering
         \includegraphics[width=\textwidth]{Figures/Minors/EventStudy_log_num_watched_sessions_unbanslots.png}
         \caption{Total Livestreams, Unbanned}
         \label{fig:token4}
     \end{subfigure}
     \vspace{0.5cm}
     \caption{\small \textbf{Time-varying treatment effect on log number of views by age group}. 1. Subfigure (a) and (c) shows the point estimates of $\tilde{\beta}_{1t}$, and subfigure (b) and (d) show the point estimates of $(\tilde{\beta}_{1t}, \tilde{\beta}_{2t})$ in specification \ref{equation-es-specification}. 2. $\tilde{\beta}_{1t}$ in one period ahead of the policy announcement is normalized to zero. 3. The 95\% point-wise confidence intervals are illustrated as the inner bars, and the standard errors are clustered at the streamer level.}
     \label{es_log_num_views}
\end{figure}

\begin{figure}[h!]
     \centering
     \begin{subfigure}[b]{0.4\textwidth}
         \centering
         \includegraphics[width=\textwidth]{Figures/Minors/EventStudy_log_num_chats_ps_gambling_banslots.png}
         \caption{Gambling, Banned}
         \label{fig:token1}
     \end{subfigure}
     \begin{subfigure}[b]{0.4\textwidth}
         \centering
         \includegraphics[width=\textwidth]{Figures/Minors/EventStudy_log_num_chats_ps_gambling_unbanslots.png}
         \caption{Gambling, Unbanned}
         \label{fig:token2}
     \end{subfigure}
     \vspace{-0.3cm}
     \begin{subfigure}[b]{0.4\textwidth}
         \centering
         \includegraphics[width=\textwidth]{Figures/Minors/EventStudy_log_num_chats_ps_unbanslots.png}
         \caption{Total Livestreams, Banned}
         \label{fig:token3}
     \end{subfigure}
     \begin{subfigure}[b]{0.4\textwidth}
         \centering
         \includegraphics[width=\textwidth]{Figures/Minors/EventStudy_log_num_chats_ps_unbanslots.png}
         \caption{Total Livestreams, Unbanned}
         \label{fig:token4}
     \end{subfigure}
     \vspace{0.5cm}
     \caption{\small \textbf{Time-varying treatment effect on log number of chats per view by age group}. 1. Subfigure (a) and (c) shows the point estimates of $\tilde{\beta}_{1t}$, and subfigure (b) and (d) show the point estimates of $\tilde{\beta}_{2t}$ in specification \ref{equation-es-specification}. 2. $(\tilde{\beta}_{1t}, \tilde{\beta}_{2t})$ in one period ahead of the policy announcement is normalized to zero. 3. The 95\% point-wise confidence intervals are illustrated as the inner bars, and the standard errors are clustered at the streamer level.}
     \label{es_log_num_chats_ps}
\end{figure}

\begin{figure}[h!]
     \centering
     \begin{subfigure}[b]{0.4\textwidth}
         \centering
         \includegraphics[width=\textwidth]{Figures/Minors/EventStudy_avg_emotes_gambling_banslots.png}
         \caption{Gambling, Banned}
         \label{fig:token1}
     \end{subfigure}
     \begin{subfigure}[b]{0.4\textwidth}
         \centering
         \includegraphics[width=\textwidth]{Figures/Minors/EventStudy_avg_emotes_gambling_unbanslots.png}
         \caption{Gambling, Unbanned}
         \label{fig:token2}
     \end{subfigure}
     \vspace{-0.3cm}
     \begin{subfigure}[b]{0.4\textwidth}
         \centering
         \includegraphics[width=\textwidth]{Figures/Minors/EventStudy_avg_emotes_banslots.png}
         \caption{Total Livestreams, Banned}
         \label{fig:token3}
     \end{subfigure}
     \begin{subfigure}[b]{0.4\textwidth}
         \centering
         \includegraphics[width=\textwidth]{Figures/Minors/EventStudy_avg_emotes_unbanslots.png}
         \caption{Total Livestreams, Unbanned}
         \label{fig:token4}
     \end{subfigure}
     \vspace{0.5cm}
     \caption{\small \textbf{Time-varying treatment effect on average number of emotes per chat message by age group}. 1. Subfigure (a) and (c) shows the point estimates of $\tilde{\beta}_{1t}$, and subfigure (b) and (d) show the point estimates of $\tilde{\beta}_{2t}$ in specification \ref{equation-es-specification}. 2. $(\tilde{\beta}_{1t}, \tilde{\beta}_{2t})$ in one period ahead of the policy announcement is normalized to zero. 3. The 95\% point-wise confidence intervals are illustrated as the inner bars, and the standard errors are clustered at the streamer level.}
     \label{es_avg_emotes}
\end{figure}

\begin{figure}[h!]
     \centering
     \begin{subfigure}[b]{0.4\textwidth}
         \centering
         \includegraphics[width=\textwidth]{Figures/Minors/EventStudy_avg_char_count_gambling_banslots.png}
         \caption{Gambling, Banned}
         \label{fig:token1}
     \end{subfigure}
     \begin{subfigure}[b]{0.4\textwidth}
         \centering
         \includegraphics[width=\textwidth]{Figures/Minors/EventStudy_avg_char_count_gambling_unbanslots.png}
         \caption{Gambling, Unbanned}
         \label{fig:token2}
     \end{subfigure}
     \vspace{-0.3cm}
     \begin{subfigure}[b]{0.4\textwidth}
         \centering
         \includegraphics[width=\textwidth]{Figures/Minors/EventStudy_avg_char_count_banslots.png}
         \caption{Total Livestreams, Banned}
         \label{fig:token3}
     \end{subfigure}
     \begin{subfigure}[b]{0.4\textwidth}
         \centering
         \includegraphics[width=\textwidth]{Figures/Minors/EventStudy_avg_char_count_unbanslots.png}
         \caption{Total Livestreams, Unbanned}
         \label{fig:token4}
     \end{subfigure}
     \vspace{0.5cm}
     \caption{\small \textbf{Time-varying treatment effect on average character count per view by age group}. 1. Subfigure (a) and (c) shows the point estimates of $\tilde{\beta}_{1t}$, and subfigure (b) and (d) show the point estimates of $\tilde{\beta}_{2t}$ in specification \ref{equation-es-specification}. 2. $(\tilde{\beta}_{1t}, \tilde{\beta}_{2t})$ in one period ahead of the policy announcement is normalized to zero. 3. The 95\% point-wise confidence intervals are illustrated as the inner bars, and the standard errors are clustered at the streamer level.}
     \label{es_avg_char_count}
\end{figure}

\begin{figure}[h!]
     \centering
     \begin{subfigure}[b]{0.4\textwidth}
         \centering
         \includegraphics[width=\textwidth]{Figures/Minors/EventStudy_avg_token_count_gambling_banslots.png}
         \caption{Gambling, Banned}
         \label{fig:token1}
     \end{subfigure}
     \begin{subfigure}[b]{0.4\textwidth}
         \centering
         \includegraphics[width=\textwidth]{Figures/Minors/EventStudy_avg_token_count_gambling_unbanslots.png}
         \caption{Gambling, Unbanned}
         \label{fig:token2}
     \end{subfigure}
     \vspace{-0.3cm}
     \begin{subfigure}[b]{0.4\textwidth}
         \centering
         \includegraphics[width=\textwidth]{Figures/Minors/EventStudy_avg_token_count_banslots.png}
         \caption{Total Livestreams, Banned}
         \label{fig:token3}
     \end{subfigure}
     \begin{subfigure}[b]{0.4\textwidth}
         \centering
         \includegraphics[width=\textwidth]{Figures/Minors/EventStudy_avg_token_count_unbanslots.png}
         \caption{Total Livestreams, Unbanned}
         \label{fig:token4}
     \end{subfigure}
     \vspace{0.5cm}
     \caption{\small \textbf{Time-varying treatment effect on average token count per view by age group}. 1. Subfigure (a) and (c) shows the point estimates of $\tilde{\beta}_{1t}$, and subfigure (b) and (d) show the point estimates of $\tilde{\beta}_{2t}$ in specification \ref{equation-es-specification}. 2. $(\tilde{\beta}_{1t}, \tilde{\beta}_{2t})$ in one period ahead of the policy announcement is normalized to zero. 3. The 95\% point-wise confidence intervals are illustrated as the inner bars, and the standard errors are clustered at the streamer level.}
     \label{es_avg_token_count}
\end{figure}

% ===================================================================

\clearpage

\begin{table}[!htb]
\centering
\renewcommand{\arraystretch}{1.2}
\begin{threeparttable}
\begin{tabular}{ll}
\hline 
Type & Example\\
\midrule
Banned websites by the policy & stake.com \\
Unbanned gambling websites detected by OCR & 1bet.com \\
Unbanned gambling websites not detected by OCR & mcluck.com \\
Unbanned sports betting websites & 1xbet.com\\
Gambling websites with different top-level domains & stake.mba\\
\hline 
\end{tabular}
\end{threeparttable}
\caption{\textbf{Types of gambling websites covered by our sample period}.}
\label{table:traffic-websites}
\end{table}

% ===================================================================

\begin{table}[!htb]
\renewcommand{\arraystretch}{1.2}
\centering

\begin{subtable}{\textwidth}
\caption{Organic Social Visits}
\label{subtable:did_organic}
\centering
\resizebox{\textwidth}{!}{%
\begin{tabular}{lcccc}
\hline
& Worldwide All Devices & Worldwide Desktop & US All Devices & US Desktop \\ \hline
Treated $\times$ After & \,0.175 & \,0.151 & \,0.160 & \,0.159 \\
                     & (0.316) & (0.235) & (0.252) & (0.192) \\ 
%Adj. $R^2$           & 0.787   & 0.784    & 0.637   & 0.582   \\
\hline
Website FE & $\checkmark$ & $\checkmark$ & $\checkmark$ & $\checkmark$ \\
Week FE     & $\checkmark$ & $\checkmark$ & $\checkmark$ & $\checkmark$ \\ 
Observations                    & 2288    & 2288     & 2288    & 2288    \\ 
Mean Dependent Variable                 & 1,743.00 & 848.15   & 39.26   & 36.87   \\ 
\hline

\end{tabular}
}
\end{subtable}

\vspace{0.5cm}

\begin{subtable}{\textwidth}
\caption{Paid Social Visits}
\label{subtable:did_paid}
\centering
\resizebox{\textwidth}{!}{%
\begin{tabular}{lcccc}
\hline
& Worldwide All Devices & Worldwide Desktop & US All Devices & US Desktop \\ \hline
Treated $\times$ After & -0.120 & -0.204\textsuperscript{*} & -0.001 & -0.041 \\
                     & (0.156) & (0.113)                   & (0.084) & (0.093) \\ \hline
%Adj. $R^2$           & 0.477   & 0.468                     & 0.718   & 0.645   \\
Website FE & $\checkmark$ & $\checkmark$ & $\checkmark$ & $\checkmark$ \\ 
Week FE     & $\checkmark$ & $\checkmark$ & $\checkmark$ & $\checkmark$ \\ 
Observations                    & 2288    & 2288     & 2288    & 2288    \\ 
Mean Dependent Variable  & 100.73  & 33.95                     & 35.99   & 13.79   \\ 
\hline
\end{tabular}
}
\end{subtable}

\caption{\textbf{The effect of Twitch's banning policy on traffic to banned gambling websites from social media platforms.} 1. Our referral traffic data from \textit{Semrush} classifies Twitch as a social media platform and breaks down the traffic by visits from organic social (Panel a) and paid social (Panel b). Organic social refers to website visits originating from unpaid, non-sponsored posts on social media platforms, whereas paid social refers to visits from social media ads. 2. All standard errors are clustered at the website level. The mean dependent variable is calculated based on observations of website visits before the policy announcement. FE = fixed effects. Significance level: *$p<0.1$; **$p<0.05$; ***$p<0.01$.}
\label{table:did_traffic_visits_breakdown}
\end{table}

% ===============================================================

\end{document}

\subsection{Demand-Side Effect}

In this section, we turn to the impact of the banning policy on viewer demand for online streams. Specifically, we are interested in whether the decrease in streams of gambling content also resulted in a decline in demand for streamers in the treated group. This potential reduction might stem not only from the decrease in streaming length itself, but also from the redistribution of streaming content, making these streams less appealing to viewers whose main interests were in online gambling. This could potentially induced these viewers to switch to other streaming platforms. However, if these viewers perceive video games with loot boxes as a substitute for original online gambling content, then we should not expect a significant reduction in demand.

To estimate the causal effect on demand, we focus on two outcome variables: log of total hours watched by all streamers in each week (``log of hours watched" hereafter) and log of hourly demand, whereas the hourly demand is calculated from dividing total hours watched by all viewers by the total streaming hours in each week. We apply the TWFE, Synthetic DiD and doubly-robust estimators on these outcome variables and report the results in Table \ref{table_demand}.

We consistently observe insignificant estimates for the log of hour watched, suggesting that the implementation of the banning policy did not affect total viewer demand for the treated streamers. Conversely, we observe mixed results in log of hourly demand. The DiD and SynthDiD estimators imply a slight increase of approximately 4.7\% in viewer demand, while the doubly-robust estimator reports an insignificant effect. Although the results are less robust compared to those in the supply-side estimations, none of our estimators show evidence of a reduction in viewer demand. Therefore, the strategic redistribution of treated streamers seems to have successfully prevented them from losing viewers. This could be attributed to either the substitution between streams featuring loot boxes and gambling content, or because livestreams of games with loot boxes attracted more new viewers compared to the loss in old viewers with an interest in gambling.

\begin{table}[h!]
\renewcommand{\arraystretch}{2}
\centering
\begin{tabular}{lccc}
\hline
                     &                                                              & Method                                                    &                                                           \\ \cline{2-4} 
                     & TWFE                                                         & SynthDiD                                                  & DR                                                        \\
Log of Hours Watched & \begin{tabular}[c]{@{}c@{}}0.0058\\ (0.0097)\end{tabular}    & \begin{tabular}[c]{@{}c@{}}0.0067\\ (0.018)\end{tabular}  & \begin{tabular}[c]{@{}c@{}}0.0017\\ (0.0173)\end{tabular} \\
Log of Hourly Demand & \begin{tabular}[c]{@{}c@{}}0.0457***\\ (0.0075)\end{tabular} & \begin{tabular}[c]{@{}c@{}}0.0473*\\ (0.023)\end{tabular} & \begin{tabular}[c]{@{}c@{}}0.038\\ (0.0204)\end{tabular}  \\ \hline
\end{tabular}
\caption{Treatment Effects in Viewer Demand}
\label{table_demand}
\end{table}

\subsection{Mechanism}

While our main results show that the banning policy resulted in a redistribution of gambling livestreams and livestreams featuring games with loot boxes, we are interested in the underlying mechanism. Specifically, we aim to examine whether the substitution from gambling content to games with loot boxes is confounded by the incentive to stream popular video games on the platform. This is crucial as several of the most popular games, such as the Call of Duty franchise, incorporate loot box content. To explore this potential confounding effect, we first use our streamer-level dataset to rank the popularity of streaming content based on the total hours watched by viewers before the policy announcement. We then assign indicators to the $N$ most popular types of content among them, for $N = 50, 100$ or $200$. We run the following game-level DiD regression:

\begin{equation}
    y_{gt} = \alpha_g + \gamma_t + \theta \cdot D_{gt} + \epsilon_{it}
\end{equation}
whereas $D_{gt} \equiv 1\{\text{Popular Game}\}_g \cdot 1\{\text{Post}\}_t$ is the indicator for popular game $g$ in the post-treatment period $t$. We use the log of total streaming hours of game $g$ in week $t$ as the outcome variable, and estimate this model with three different samples: a full sample including all games released before the implementation of the policy, a sub-sample including only games with loot boxes, and a sub-sample excluding all games with loot boxes. For robustness checks, we also construct two sets of popularity indicators: one from $N$ most popular content list without non-gaming content, and one from the list including non-gaming content. The results are reported in "Gaming Top" columns and "Non Gaming Top" columns respectively.

\begin{table}[h!]
\footnotesize
\renewcommand{\arraystretch}{1.5}
\centering
\begin{tabular}{lrrrrrr}
\hline
        & \multicolumn{2}{c}{Full Sample}                                                                                               & \multicolumn{2}{c}{Loot Box Sample}                                                                                                            & \multicolumn{2}{c}{Non-Loot Sample}                                                                                           \\
        & \multicolumn{1}{c}{Gaming Top}                                & \multicolumn{1}{c}{Non Gaming Top}                            & \multicolumn{1}{c}{Gaming Top}                                                 & \multicolumn{1}{c}{Non Gaming Top}                            & \multicolumn{1}{c}{Gaming Top}                                & \multicolumn{1}{c}{Non Gaming Top}                            \\ \cline{2-7} 
Top 50  & \begin{tabular}[c]{@{}r@{}}-0.1533***\\ (0.0128)\end{tabular} & \begin{tabular}[c]{@{}r@{}}-0.1433***\\ (0.0126)\end{tabular} & \begin{tabular}[c]{@{}r@{}}-0.0300$\cdot$\\ (0.0158)\end{tabular} & \begin{tabular}[c]{@{}r@{}}-0.0314*\\ (0.0159)\end{tabular}   & \begin{tabular}[c]{@{}r@{}}-0.2347***\\ (0.0183)\end{tabular} & \begin{tabular}[c]{@{}r@{}}-0.2031***\\ (0.0174)\end{tabular} \\
Top 100 & \begin{tabular}[c]{@{}r@{}}-0.1374***\\ (0.0104)\end{tabular} & \begin{tabular}[c]{@{}r@{}}-0.1244***\\ (0.0103)\end{tabular} & \begin{tabular}[c]{@{}r@{}}-0.0192\\ (0.0155)\end{tabular}                     & \begin{tabular}[c]{@{}r@{}}-0.0161\\ (0.0155)\end{tabular}    & \begin{tabular}[c]{@{}r@{}}-0.1859***\\ (0.0133)\end{tabular} & \begin{tabular}[c]{@{}r@{}}-0.1641***\\ (0.0129)\end{tabular} \\
Top 200 & \begin{tabular}[c]{@{}r@{}}-0.0916***\\ (0.0095)\end{tabular} & \begin{tabular}[c]{@{}r@{}}-0.1859***\\ (0.0133)\end{tabular} & \begin{tabular}[c]{@{}r@{}}-0.0754***\\ (0.0174)\end{tabular}                  & \begin{tabular}[c]{@{}r@{}}-0.0749***\\ (0.0173)\end{tabular} & \begin{tabular}[c]{@{}r@{}}-0.0912***\\ (0.0113)\end{tabular} & \begin{tabular}[c]{@{}r@{}}-0.0909***\\ (0.0112)\end{tabular} \\ \hline
\end{tabular}
\caption{The impact of the banning Policy on Popular Games}
\label{table_mech}
\end{table}

Table \ref{table_mech} summarizes the results from our game-level DiD regressions. The estimates show a negative effect of approximately 5\%-10\% decrease in streaming hours of more popular content after the implementation of the policy. There are two main reasons for the negative estimates: first, popular games are more likely to be released after October to gain benefit from holiday sales battles. Second, streamers may decrease streaming hours in general at the end of the year due to family activities. If the substitution between livestreams of online gambling and games with loot boxes was primarily driven by the popularity of the games, then we would expect very similar results from the sample of games with loot boxes and the sample of games without loot boxes. However, by comparing the average treatment effects, we find that popular games with loot boxes are less impacted by the banning policy. This result provides partial evidence suggesting that the feature of loot boxes compensates for the decrease in supply of these popular content, because these games were perceived as substitutes for the prohibited gambling content after the policy announcement. 


\afterpage{
\begin{landscape}
\begin{table}[]
\renewcommand{\arraystretch}{1.5}
\centering
\begin{tabular}{llccc}
\hline
                                        &                            & Overall Livestreams                                           & Gambling Livestreams                                            & Livestreams of Loot Boxes                                     \\ \cline{3-5} 
\multirow{4}{*}{Log of Streaming Hours} & Treated $\times$ Post         & \begin{tabular}[c]{@{}c@{}}-0.0922***\\ (0.0235)\end{tabular} & \begin{tabular}[c]{@{}c@{}}-0.5711***\\ (0.0529)\end{tabular}   & \begin{tabular}[c]{@{}c@{}}0.2326***\\ (0.0444)\end{tabular}  \\
                                        & Treated $\times$ Announcement & \begin{tabular}[c]{@{}c@{}}0.0231\\ (0.0200)\end{tabular}     & \begin{tabular}[c]{@{}c@{}}0.2110***\\ (0.0463)\end{tabular}    & \begin{tabular}[c]{@{}c@{}}0.0209\\ (0.0352)\end{tabular}     \\
                                        & Log of Streaming Hours     & -                                                             & \begin{tabular}[c]{@{}c@{}}0.0293***\\ (0.0042)\end{tabular}    & \begin{tabular}[c]{@{}c@{}}0.6339***\\ (0.0110)\end{tabular}  \\
                                        & Log of Average Viewership  & \begin{tabular}[c]{@{}c@{}}0.3213***\\ (0.0018)\end{tabular}  & \begin{tabular}[c]{@{}c@{}}-0.0019 cdot\\ (0.0010)\end{tabular} & \begin{tabular}[c]{@{}c@{}}-0.0205***\\ (0.0032)\end{tabular} \\ \cline{2-5} 
\multirow{4}{*}{Share of Content}       & Treated $\times$ Post         & -                                                             & \begin{tabular}[c]{@{}c@{}}-0.1240***\\ (0.0136)\end{tabular}   & \begin{tabular}[c]{@{}c@{}}0.0876***\\ (0.0137)\end{tabular}  \\
                                        & Treated $\times$ Announcement & -                                                             & \begin{tabular}[c]{@{}c@{}}0.04662***\\ (0.0110)\end{tabular}   & \begin{tabular}[c]{@{}c@{}}0.0029\\ (0.0105)\end{tabular}     \\
                                        & Log of Streaming Hours     & -                                                             & \begin{tabular}[c]{@{}c@{}}0.0021 $\cdot$\\ (0.0011)\end{tabular}  & \begin{tabular}[c]{@{}c@{}}-0.0144***\\ (0.0034)\end{tabular} \\
                                        & Log of Average Viewership  & -                                                             & \begin{tabular}[c]{@{}c@{}}0.0016***\\ (0.0004)\end{tabular}    & \begin{tabular}[c]{@{}c@{}}0.0655***\\ (0.0017)\end{tabular}  \\ \hline
\end{tabular}
\caption{Treated Effects in Streaming Hours and Shares of Content}
\label{table_did}
\end{table}
\end{landscape}}



\end{document}